
\providecommand{\dR}{\mathrm{dR}}
\providecommand{\op}{\mathrm{op}}
\providecommand{\Map}{\operatorname{Map}}
\providecommand{\map}{\operatorname{map}}
\providecommand{\Hom}{\operatorname{Hom}}
\providecommand{\Spec}{\operatorname{Spec}}
\providecommand{\Sp}{\mathrm{Sp}}
\providecommand{\Tot}{\operatorname*{Tot}}
\providecommand{\CGrp}{\operatorname{CGrp}}
\providecommand{\CAlg}{\operatorname{CAlg}}
\providecommand{\Fun}{\operatorname{Fun}}
\providecommand{\fib}{\operatorname{fib}}
\providecommand{\hofib}{\operatorname{hofib}}
\providecommand{\cofib}{\operatorname{cofib}}
\providecommand{\colim}{\operatorname*{colim}}
\providecommand{\limop}{\operatorname*{lim}}
\providecommand{\Sym}{\operatorname{Sym}}
\providecommand{\gr}{\operatorname{gr}}
\providecommand{\id}{\operatorname{id}}
\providecommand{\DR}{\operatorname{DR}}
\providecommand{\Sm}{\operatorname{Sm}}
\providecommand{\fp}{\mathrm{fp}}
\providecommand{\Nis}{\mathrm{Nis}}
\providecommand{\cl}{\mathrm{cl}}
\providecommand{\add}{\mathrm{add}}
\providecommand{\Am}{\operatorname{Am}}
\providecommand{\sk}{\operatorname{sk}}
\providecommand{\OmegaInf}{\Omega_{\mathrm{inf}}}

\providecommand{\SH}{\mathrm{SH}}
\providecommand{\Disc}{\operatorname{Disc}}
\providecommand{\RealDiff}{\operatorname{Real}^{0}}
\providecommand{\ObsBNV}{\operatorname{Obs}_{\mathrm{BNV}}}
\providecommand{\Trdiff}{\operatorname{tr}^{\partial}}
\providecommand{\IntBNV}{\int^{\mathrm{BNV}}}
\providecommand{\Transdiff}{\operatorname{Trans}^{\partial}}
\providecommand{\Shv}{\mathrm{Shv}}



\chapter{Brave-New Differential Forms, de Rham Theory, and Differential Umkehr}
\label{chap:brave-new-monadic-de-rham}
\label{ch:brave-new-modal-de-rham}
\label{ch:brave-new-differential-forms}
\label{chap:brave-new-de-rham-umkehr}

\section{Purpose, continuity, and scope}
\label{sec:monadic-de-rham-purpose}

Chapter~\ref{chap:brave-new-infinitesimal-geometry} supplied both pieces of
ambient structure needed here.

On the infinitesimal side it constructed the big affine spectral Nisnevich
\(\infty\)-topos
\[
    \mathcal X_{\Sp},
\]
the intrinsic left-exact de Rham localization \(L_{\dR}\), its relative slice
localizations, the relative de Rham objects
\[
    X_{\dR/S}=S\times_{S_{\dR}}X_{\dR},
\]
their units
\[
    \eta_{X/S}:X\longrightarrow X_{\dR/S},
\]
the arbitrary pullback base-change equivalences, the notions of formal
\(\acute{e}\)taleness and formal smoothness, and the theorem that the relative
de Rham unit of a represented smooth spectral morphism is an effective
epimorphism.  We also retain the finite affine Nisnevich basis, affine density
in the big topos and its slices, and the stable Nisnevich coefficient category
on the smooth site.

On the stable differential side the same chapter constructed the two-stage
coefficient architecture
\[
\boxed{
    \mathcal P_S^{\mathrm{pre}}
    \xrightarrow{\,L_{\mathrm{Th},S}\,}
    \mathcal R_S^0
    \xrightarrow{\,\mathbb H_S\,}
    \SH(S).
}
\]
The first arrow imposes Thom stabilization and no affine-line localization.
The second imposes stable \(\mathbb A^1\)-invariance and sits in the exact BNV
triple
\[
    \mathbb H_S\dashv\Disc_S\dashv\mathbb S_S.
\]
Thus the BNV deformation and cycle functors
\[
    \mathcal A_S,\qquad \mathcal Z_S
\]
measure the affine-line defect inside an already Thom-stable theory.

For every smooth separated finite-presentation morphism
\[
    p:X\longrightarrow Y,
\]
Chapter~\ref{chap:brave-new-infinitesimal-geometry} also constructed the
Thom-stable exceptional lift
\[
\boxed{
    \widetilde p_!
    :=
    p_\sharp^0\Sigma^{-T_p,0}
    \dashv
    \Sigma^{T_p,0}p^{*,0}
    =:
    \widetilde p^{\,!}
}
\]
and proved the motivic comparisons
\[
\boxed{
    \mathbb H_Y\widetilde p_!
    \simeq
    p_!\mathbb H_X,
    \qquad
    \mathbb H_X\widetilde p^{\,!}
    \simeq
    p^!\mathbb H_Y.
}
\]
It further constructed the BNV cycle defect
\[
    \omega_p:
    \mathcal Z_Y\widetilde p_!
    \longrightarrow
    \mathcal Z_Y\widetilde p_!\mathcal Z_X
\]
and its obstruction object
\[
\boxed{
    \ObsBNV(p)
    :=
    \cofib(\omega_p)
    \simeq
    \Sigma\,\mathcal Z_Y\widetilde p_!\mathcal S_X.
}
\]
The obstruction vanishes for split finite \(\acute{e}\)tale morphisms and,
by Nisnevich descent, for the Nisnevich-locally split finite
\(\acute{e}\)tale class proved there.  No general smooth vanishing theorem is
assumed.

The present chapter has two consecutive tasks.

The first is intrinsic de Rham geometry.  We begin with the inherited
relative de Rham unit
\[
    \eta_{X/S}:X\longrightarrow X_{\dR/S}
\]
and construct its effective image
\[
    X\xrightarrow{\epsilon_{X/S}}Q_{X/S}
      \xrightarrow{\jmath_{X/S}}X_{\dR/S}.
\]
The effective epimorphism \(X\to Q_{X/S}\) has a \v{C}ech nerve.  That nerve
is the intrinsic effective relation groupoid determined by the relative de
Rham unit.  For a morphism \(Z\to Y\) over \(S\) we also construct the
intrinsic relative de Rham completion
\[
    \widehat Y^{\,\dR}_{Z/S}
    :=
    Y\times_{Y_{\dR/S}}Z_{\dR/S}.
\]
The infinitesimal relation simplex \(R_{X/S,n}\) is then identified with the
intrinsic de Rham completion of the diagonal
\[
    X\longrightarrow X^{\times_S(n+1)}.
\]

Two scalar-valued de Rham objects are kept distinct.  The full or modal object
is the algebra of additive scalar functions on \(X_{\dR/S}\).  The effective
or \v{C}ech object is the algebra of additive scalar functions on
\(Q_{X/S}\), equivalently the totalization of the relation-nerve cochains.
There is a canonical comparison
\[
    \mathbf{DR}^{\mathrm{full}}_{X/S}
    \longrightarrow
    \mathbf{DR}^{\mathrm{eff}}_{X/S},
\]
and smooth effectivity makes it an equivalence for represented smooth spectral
morphisms.

The scalar object is the already constructed brave-new affine line
\[
    \mathbb A^1_S
    =
    \Spec_S\Sym_{\mathcal O_S}(\mathcal O_S)
\]
when \(S\) is represented.  We do not introduce a second affine line.  The
same free-algebra universal property supplies its additive law and additive
inverse, and stabilizing the additive commutative group gives
\[
    \mathbb O_S^{\add}.
\]

For an arbitrary \(T\in\mathcal X_{\Sp}\) we use the parameterized stable
category
\[
    \mathcal E_T:=\Sp(\mathcal X_{\Sp,/T}).
\]
This category is an internal tool: effective quotients, relative de Rham
objects, and infinitesimal simplices need not be represented or smooth.  It is
not a replacement for \(\mathcal R_S^0\).  At the smooth-site interface the
intrinsic coefficient objects are first placed in
\[
    \mathcal C_S
    :=
    \Shv_{\Nis}(\Sm^{\fp}_{/S};\Sp)
\]
and then Thom-stabilized into \(\mathcal R_S^0\).

There is an unavoidable notational collision which we resolve once and for
all.  The symbol
\[
    Q_{X/S}
\]
continues to mean the \emph{effective infinitesimal quotient}.  The preceding
chapter also used \(Q_S\) for the level-zero Thom-stable realization functor.
In the present chapter we rename that functor
\[
\boxed{
    \RealDiff_S
    :=
    L_{\mathrm{Th},S}F_{0,S}:
    \mathcal C_S\longrightarrow\mathcal R_S^0.
}
\]
Thus \(\RealDiff_S\) is precisely the functor called \(Q_S\) in
Chapter~\ref{chap:brave-new-infinitesimal-geometry}.  It is exact,
colimit-preserving and strong symmetric monoidal, but it is not asserted to
be fully faithful.

The normalized objects extracted from the relation nerve are deliberately
called \emph{normalized infinitesimal de Rham cochains}.  They retain all
finite infinitesimal orders present in the de Rham completion.  They are not
defined to be exterior powers of the cotangent complex and their totalization
filtration is not called the Hodge filtration.  For represented spectral
morphisms we construct a canonical first-order linearization of the degree-one
normalized cochain to the cotangent complex and prove that the normalized
infinitesimal differential linearizes to the universal derivation.  An
explicit smooth affine-line calculation shows that higher infinitesimal terms
survive in the normalized cochain and are killed by this first-order
linearization.  Thus no spectral HKR or exterior-cotangent comparison is
silently assumed.

The second task is integration.  After the intrinsic de Rham package has been
realized by \(\RealDiff_S\), the already constructed exceptional adjunction
\[
    \widetilde p_!\dashv\widetilde p^{\,!}
\]
acts on those coefficients.  Its counit
\[
    \widetilde p_!\widetilde p^{\,!}D
    \longrightarrow
    D
\]
is the differential exceptional trace.  Its motivic shadow is the counit
\[
    p_!p^!\mathbb H_YD
    \longrightarrow
    \mathbb H_YD.
\]
When the constructed BNV obstruction of \(p\) vanishes, the same differential
Umkehr functor preserves the universal BNV interval primitive.  Thus the
motivic exceptional integral and the BNV secondary integral are not
identified as maps of the same type; they are the primary and secondary
layers of one differential transfer.

For a general smooth morphism we do not assert that the intrinsic source de
Rham coefficient is canonically the exceptional pullback of the target
coefficient.  Accordingly, the general construction integrates
\emph{differential dualizing de Rham coefficients}
\[
    \widetilde p^{\,!}D.
\]
For represented spectral \(\acute{e}\)tale change of base, the intrinsic
relative de Rham objects, effective quotients, relation nerves, and scalar-valued
de Rham coefficient systems are invariant under insertion of the
\(\acute{e}\)tale base.  Consequently the smooth-site coefficient exchange is
an equivalence for every represented spectral \(\acute{e}\)tale morphism of
finite presentation.  If such a morphism is separated, then its tangent class
vanishes and the differential exceptional trace becomes a genuine pushforward
of the intrinsic source de Rham coefficients themselves.  In particular this
applies to every finite \(\acute{e}\)tale morphism.  The stronger BNV
compatibility remains controlled by the previously constructed BNV obstruction;
its vanishing is used below only in the verified Nisnevich-locally split finite
\(\acute{e}\)tale class.

The logical progression of the chapter is therefore
\[
\boxed{
\begin{gathered}
X\xrightarrow{\eta_{X/S}}X_{\dR/S}
\longmapsto
\bigl(Q_{X/S},R_{X/S,\bullet}\bigr)
\longmapsto
\underline{\mathbf C}_{\dR}^{\,\bullet}(X/S)
\longleftrightarrow
\mathbb M_{X/S}
\\[4pt]
\longmapsto
\left(
    \mathbf{DR}^{\mathrm{full}}_{X/S}
    \longrightarrow
    \mathbf{DR}^{\mathrm{eff}}_{X/S}
\right)
\longmapsto
\bigl(
    F_{\mathrm{Tot}}^\bullet,
    \underline{\OmegaInf}^{\,\bullet}(X/S),
    d_{\inf}
\bigr)
\\[4pt]
\longmapsto
\mathcal C_S
\xrightarrow{\ \RealDiff_S\ }
\mathcal R_S^0
\xrightarrow{\ \widetilde p_!\dashv\widetilde p^{\,!}\ }
\mathcal R_Y^0
\xrightarrow{\ \mathbb H_Y\ }
\SH(Y).
\end{gathered}
}
\]

The arrow to the descent monad is a presentation theorem: the de Rham
cochains are defined geometrically as additive functions on the infinitesimal
nerve and are then identified with the Amitsur object of the infinitesimal
descent monad.

The complete \v{C}ech-totalization tower supplies the filtration used in the
chapter.  Stable Dold--Kan identifies its successive layers with normalized
infinitesimal cochains.  The complete-filtered/coherent-cochain correspondence
packages the connecting morphisms into the coherent infinitesimal
differential, its specified nullhomotopy of the square, and all higher coherent
relations.  In the homotopy category, ordered Alexander--Whitney normalization
gives the normalized cochains their canonical associative unital product and
the infinitesimal differential satisfies the graded Leibniz identity.  The
total and partially totalized effective de Rham objects retain the canonical
\(E_\infty\)-algebra structures obtained as limits of the cosimplicial
commutative algebra.  No homotopy-coherent \(E_1\)-lift of the un-totalized
Alexander--Whitney product, no multiplicative totalization filtration, no
Hodge identification, and no symmetric-shuffle commutative algebra structure
on the un-totalized normalized cochains is asserted.

\begin{convention}[Universes and sheaves]
\label{conv:monadic-de-rham-universes}
Use the universe convention of
Chapter~\ref{chap:brave-new-infinitesimal-geometry}, and enlarge it once.
Thus, if that convention fixes Grothendieck universes
\[
    \mathcal U\subset\mathcal V\subset\mathcal W,
\]
fix in addition a Grothendieck universe
\[
    \mathcal W\subset\mathcal Z.
\]
Affine spectral tests are \(\mathcal U\)-small, the fixed native site skeleta
are \(\mathcal V\)-small, and the native big topos and its objects are formed
in \(\mathcal W\).  For an arbitrary \(T\in\mathcal X_{\Sp}\), the
affine-point category \(\mathcal C^{\mathrm{aff}}_{/T}\) is understood through
a chosen \(\mathcal W\)-small skeleton, and its sheaf category, the
corresponding parameterized stable category, and the functor categories built
from them are formed in \(\mathcal Z\).  The same enlargement is used when
choosing sets of square-zero generators in arbitrary slices.  Sheaves are not
hypercompleted unless explicitly stated otherwise.
\end{convention}

\begin{convention}[Geometric scope]
\label{conv:monadic-de-rham-scope}
The intrinsic construction is defined for arbitrary morphisms \(X\to S\) in
\(\mathcal X_{\Sp}\).  Representability, smoothness, or the standing motivic
base conditions are imposed only in statements which explicitly use them.  In
the terminology of the preceding chapters, a smooth spectral morphism already
includes local finite presentation; global finite presentation additionally
includes the quasi-compactness and quasi-separatedness conditions of the
preceding chapters.
\end{convention}

\begin{convention}[Stagewise exactness and complete totalization]
\label{conv:monadic-de-rham-stagewise-exactness}
The differential Umkehr constructions below are applied to the already
constructed totalization stages, their finite fibre/cofiber sequences, and the
coherent normalized cochain object extracted from the complete filtration.
Exactness of \(\RealDiff_S\), \(\widetilde p_!\), or
\(\widetilde p^{\,!}\) is used only to preserve those finite stable diagrams.

No commutation of \(\widetilde p_!\) with the sequential inverse limit defining
the complete totalization is inferred merely from exactness.
\end{convention}

\subsection{Chapter-level output}
\label{subsec:monadic-de-rham-output}

Retaining the infinitesimal and Thom-stable differential framework of
Chapter~\ref{chap:brave-new-infinitesimal-geometry}, the present chapter
constructs:
\begin{enumerate}
\item the effective quotient \(Q_{X/S}\), effective infinitesimal relation
groupoid, and intrinsic relative de Rham completions
\[
    R_{X/S,\bullet}
    \simeq
    \check C_\bullet(X\to Q_{X/S}),
    \qquad
    \widehat Y^{\,\dR}_{Z/S}
    =
    Y\times_{Y_{\dR/S}}Z_{\dR/S};
\]
\item parameterized stable coefficients
\[
    \mathcal E_T=\Sp(\mathcal X_{\Sp,/T})
\]
with the stable adjoint triple, Beck--Chevalley, projection, closed base
change, and effective descent;
\item the additive stabilization
\[
    \mathbb O_S^{\add}\in\CAlg(\mathcal E_S)
\]
of the existing brave-new affine line;
\item the geometric cosimplicial de Rham algebra
\[
    \underline{\mathbf C}_{\dR}^{\,\bullet}(X/S)
\]
of additive scalar functions on infinitesimal simplices and its
internal-function presentation;
\item the infinitesimal descent monad
\[
    \mathbb M_{X/S}=\Pi_{\epsilon_{X/S}}\epsilon_{X/S}^*
\]
and the identification of its Amitsur object with the geometric de Rham
cochains;
\item the full and effective de Rham algebras
\[
    \mathbf{DR}^{\mathrm{full}}_{X/S},
    \qquad
    \mathbf{DR}^{\mathrm{eff}}_{X/S},
\]
their canonical comparison, and their equality under smooth effectivity;
\item the complete \v{C}ech-totalization tower, normalized infinitesimal de
Rham cochains, coherent infinitesimal differential, ordered
Alexander--Whitney product in the homotopy category, coherently closed
cochains, exact-cochain classifiers, primitive spaces, and
additive-delooping presentation;
\item the first structured infinitesimal neighbourhood, stabilized cotangent
coefficient, and first-order linearization
\[
    \lambda^1_{X/S}:
    \underline{\OmegaInf}^{\,1}(X/S)
    \longrightarrow
    \Pi_p\mathbb L^{\mathrm{st}}_{X/S},
\]
with
\[
    \lambda^1_{X/S}d_{\inf}^0
    \simeq
    d^{\mathrm{st}}_{X/S},
\]
together with an explicit affine-line class showing that
\(\lambda^1_{X/S}\) need not be an equivalence;
\item contravariant functoriality, arbitrary parameterized base change,
spectral Nisnevich descent, external products, formal-\(\acute{e}\)tale
Cartesianity, formal-\(\acute{e}\)tale invariance under insertion of an
intermediate base, and local linear completion models;
\item smooth-site coefficient objects in \(\mathcal C_S\), followed by their
correct Thom-stable differential realization
\[
\boxed{
    \RealDiff_S
    =
    L_{\mathrm{Th},S}F_{0,S}:
    \mathcal C_S\longrightarrow\mathcal R_S^0,
}
\]
with no full-faithfulness assertion;
\item stabilized evaluation and smooth-direct-image exchange formulas for
\(\RealDiff_S\);
\item differential exceptional traces on the realized de Rham stage system,
whose homotopifications are the motivic exceptional traces;
\item the exact BNV obstruction controlling compatibility of these traces
with universal BNV interval integration;
\item unconditional intrinsic de Rham integration and motivic shadow for every
separated spectral \(\acute{e}\)tale morphism of finite presentation, together
with BNV compatibility for the verified Nisnevich-locally split finite
\(\acute{e}\)tale class;
\item differential transgression for smooth correspondences once an explicit
coefficient trace into the differential dualizing coefficient has been
constructed.
\end{enumerate}

\section{From relative de Rham geometry to the effective infinitesimal nerve}
\label{sec:monadic-de-rham-modality}
\label{sec:modal-infinitesimal-nerve}

Fix \(S\in\mathcal X_{\Sp}\) and a morphism \(X\to S\).  The relative de Rham object
\(X_{\dR/S}\) and the unit \(\eta_{X/S}\) are those of the preceding chapter.

\begin{definition}[Effective infinitesimal quotient]
\label{def:monadic-de-rham-effective-quotient}
Factor the inherited relative de Rham unit in the effective-epimorphism--monomorphism factorization
system of the slice:
\[
\boxed{
    X\xrightarrow{\epsilon_{X/S}}Q_{X/S}
    \xrightarrow{\jmath_{X/S}}X_{\dR/S}.
}
\]
Thus \(\epsilon_{X/S}\) is an effective epimorphism and \(\jmath_{X/S}\) is a monomorphism.  We call
\(Q_{X/S}\) the effective infinitesimal quotient of \(X\) over \(S\).
\end{definition}

\begin{proposition}[Functoriality and base change]
\label{prop:monadic-de-rham-quotient-base-change}
The assignment \(X\mapsto Q_{X/S}\) is functorial in the slice
\(\mathcal X_{\Sp,/S}\).  For every \(c:S'\to S\), with
\(X'=X\times_SS'\), there is a canonical equivalence
\[
\boxed{
    Q_{X'/S'}\xrightarrow{\simeq}Q_{X/S}\times_SS'
}
\]
compatible with identities, composition, and higher pullback-pasting coherences.
\end{proposition}

\begin{proof}
Naturality of the relative de Rham unit gives a functorial arrow
\(X\to X_{\dR/S}\).  Functorial factorization by the
effective-epimorphism--monomorphism factorization system gives the functoriality of its image.
The preceding chapter proves
\[
    X'_{\dR/S'}\simeq X_{\dR/S}\times_SS'.
\]
Effective epimorphisms and monomorphisms in an \(\infty\)-topos are stable under pullback.  Pulling
back the image factorization of \(\eta_{X/S}\) therefore gives the image factorization of
\(\eta_{X'/S'}\), which yields the displayed equivalence.  The coherence is inherited from the
coherent base-change equivalence for relative de Rham objects and from functorial image
factorization.
\end{proof}

\begin{theorem}[Smooth effectivity]
\label{thm:monadic-de-rham-smooth-effectivity}
If \(X\to S\) is represented by a smooth spectral morphism, then
\[
\boxed{
    Q_{X/S}\xrightarrow{\simeq}X_{\dR/S}.
}
\]
\end{theorem}

\begin{proof}
The smooth-effectivity theorem of the preceding chapter says that
\(\eta_{X/S}:X\to X_{\dR/S}\) is an effective epimorphism.  Its effective image is therefore its
whole target.
\end{proof}

\begin{definition}[Infinitesimal relation nerve]
For \(n\geq0\), define
\[
\boxed{
    R_{X/S,n}
    :=
    \underbrace{
        X\times_{X_{\dR/S}}\cdots\times_{X_{\dR/S}}X
    }_{n+1\text{ factors}}.
}
\]
Deletion and repetition of factors give a simplicial object
\[
    R_{X/S,\bullet}:\Delta^{\op}\longrightarrow\mathcal X_{\Sp,/S}.
\]
Write
\[
    a_n:R_{X/S,n}\longrightarrow S
\]
for its structural morphism.
\end{definition}

\begin{convention}[Ordered \v{C}ech faces in degree one]
\label{conv:monadic-de-rham-cech-face-order}
In degree one write
\[
    R_{X/S,1}
    =
    X\times_{X_{\dR/S}}X
\]
with projections
\[
    \operatorname{pr}_1,\operatorname{pr}_2:R_{X/S,1}\rightrightarrows X.
\]
The two geometric face maps of the simplicial \v{C}ech nerve are ordered by
\[
    \partial_0=\operatorname{pr}_2,
    \qquad
    \partial_1=\operatorname{pr}_1;
\]
that is, \(\partial_i\) deletes the \(i\)-th factor when the two factors are
numbered \(0,1\).  After applying scalar-valued functions, the induced
cosimplicial cofaces are therefore
\[
    d^0=\operatorname{pr}_2^*,
    \qquad
    d^1=\operatorname{pr}_1^*,
\]
and the degree-zero alternating coface difference is
\[
\boxed{
    d^0-d^1
    =
    \operatorname{pr}_2^*-\operatorname{pr}_1^*.
}
\]
This convention is retained throughout the chapter.
\end{convention}

\begin{proposition}[Simplex structure]
\label{prop:modal-simplex-structure}
The simplicial object \(R_{X/S,\bullet}\) is a groupoid object.  Its face maps delete a factor and
its degeneracy maps repeat a factor.  The construction is functorial in \(X\to S\).
\end{proposition}

\begin{proof}
It is the \v{C}ech nerve of the inherited map \(X\to X_{\dR/S}\).  The standard \v{C}ech face and
degeneracy maps give the stated formulas and all simplicial identities.  Every \v{C}ech nerve in an
\(\infty\)-topos is a groupoid object.
\end{proof}

\begin{theorem}[Effective infinitesimal presentation]
\label{thm:effective-infinitesimal-presentation}
There are canonical equivalences of simplicial objects
\[
\boxed{
    R_{X/S,\bullet}
    \simeq
    \check C_\bullet(\epsilon_{X/S}:X\to Q_{X/S})
}
\]
and a canonical realization equivalence
\[
\boxed{
    |R_{X/S,\bullet}|\simeq Q_{X/S}.
}
\]
For represented smooth \(X\to S\), this is the effective \v{C}ech presentation of
\(X_{\dR/S}\).
\end{theorem}

\begin{proof}
Because \(\jmath_{X/S}:Q_{X/S}\to X_{\dR/S}\) is a monomorphism,
\[
    Q_{X/S}\times_{X_{\dR/S}}Q_{X/S}\simeq Q_{X/S}.
\]
Using \(X\to Q_{X/S}\to X_{\dR/S}\), we obtain
\[
\begin{aligned}
    X\times_{X_{\dR/S}}X
    &\simeq
    X\times_{Q_{X/S}}
    \bigl(Q_{X/S}\times_{X_{\dR/S}}Q_{X/S}\bigr)
    \times_{Q_{X/S}}X
    \\
    &\simeq
    X\times_{Q_{X/S}}X.
\end{aligned}
\]
The same argument in every iterated degree identifies the two \v{C}ech nerves, compatibly with all
faces and degeneracies.  Since \(\epsilon_{X/S}\) is an effective epimorphism, its \v{C}ech nerve
realizes to \(Q_{X/S}\).  The final assertion follows from
Theorem~\ref{thm:monadic-de-rham-smooth-effectivity}.
\end{proof}

\begin{theorem}[Base change of the infinitesimal nerve]
\label{thm:modal-nerve-base-change}
For \(c:S'\to S\) and \(X'=X\times_SS'\), there are canonical equivalences
\[
\boxed{
    R_{X'/S',n}\simeq R_{X/S,n}\times_SS',
    \qquad
    Q_{X'/S'}\simeq Q_{X/S}\times_SS'
}
\]
compatible with every face, degeneracy, identity, composition, and higher base-change coherence.
\end{theorem}

\begin{proof}
The quotient statement is Proposition~\ref{prop:monadic-de-rham-quotient-base-change}.  The
preceding chapter gives
\(X'_{\dR/S'}\simeq X_{\dR/S}\times_SS'\).  Pullback commutes with the iterated fibre products
defining \(R_{X/S,n}\), which gives the first equivalence and all simplicial coherences.
\end{proof}

\begin{proposition}[Functoriality in the geometric input]
\label{prop:modal-nerve-functoriality}
A morphism \(f:X\to Y\) over \(S\) induces a simplicial morphism
\[
    R_{X/S,\bullet}\longrightarrow R_{Y/S,\bullet}
\]
and a compatible morphism
\[
    Q_{X/S}\longrightarrow Q_{Y/S}.
\]
These assignments preserve identities, composition, and base change.
\end{proposition}

\begin{proof}
Naturality of the relative de Rham unit gives a commutative square from
\(X\to X_{\dR/S}\) to \(Y\to Y_{\dR/S}\).  Applying the induced maps to all factors gives the
simplicial morphism, while functoriality of image factorization gives the quotient morphism.
\end{proof}


\subsection{Intrinsic relative de Rham completions}
\label{subsec:monadic-de-rham-completion}

\begin{definition}[Intrinsic relative de Rham completion]
\label{def:intrinsic-relative-dr-completion}
Let \(Z\to Y\to S\) be morphisms in \(\mathcal X_{\Sp}\).  Define the intrinsic relative de Rham
completion of \(Y\) along \(Z\) over \(S\) by
\[
\boxed{
    \widehat Y^{\,\dR}_{Z/S}
    :=
    Y\times_{Y_{\dR/S}}Z_{\dR/S}.
}
\]
It is an object over \(Y\) equipped with a canonical map from \(Z\).
\end{definition}

\begin{proposition}[Base change of intrinsic completions]
\label{prop:intrinsic-relative-dr-completion-base-change}
For every \(c:S'\to S\), put
\[
    Y':=Y\times_SS',
    \qquad
    Z':=Z\times_SS'.
\]
Then there is a canonical equivalence
\[
\boxed{
    \widehat{Y'}^{\,\dR}_{Z'/S'}
    \xrightarrow{\simeq}
    \widehat Y^{\,\dR}_{Z/S}\times_SS'.
}
\]
These equivalences satisfy identity, composition, and all higher pullback-pasting coherences.
\end{proposition}

\begin{proof}
The relative de Rham base-change theorem of the preceding chapter gives
\[
    Y'_{\dR/S'}\simeq Y_{\dR/S}\times_SS',
    \qquad
    Z'_{\dR/S'}\simeq Z_{\dR/S}\times_SS'.
\]
Substituting these equivalences into the defining pullback of
\(\widehat{Y'}^{\,\dR}_{Z'/S'}\) gives the result.  The coherence is inherited from the coherent
base-change equivalences for relative de Rham objects and from functoriality of finite limits.
\end{proof}


\begin{lemma}[Target locality of relative de Rham equivalences]
\label{lem:relative-dr-equivalence-target-local}
Fix \(S\in\mathcal X_{\Sp}\).  Let
\[
    f:A\longrightarrow B
\]
be a morphism in \(\mathcal X_{\Sp,/S}\), let
\[
    e:B'\longrightarrow B
\]
be an effective epimorphism, and put
\[
    A':=A\times_BB'.
\]
If the pullback
\[
    f':A'\longrightarrow B'
\]
is inverted by \(L_{\dR,/S}\), then \(f\) is inverted by \(L_{\dR,/S}\).
\end{lemma}

\begin{proof}
Let \(B'_\bullet\to B\) be the \v{C}ech nerve of \(e\).  Pullback stability of effective
epimorphisms gives an effective epimorphism
\[
    A'=A\times_BB'\longrightarrow A
\]
whose \v{C}ech nerve is
\[
    A'_\bullet\simeq A\times_BB'_\bullet.
\]
Since \(L_{\dR,/S}\) is left exact, its local-equivalence class is stable under pullback.  Hence for
every \(n\), the map
\[
    A\times_BB'_n\longrightarrow B'_n
\]
is a relative de Rham equivalence.

Let \(F\to S\) be \(L_{\dR,/S}\)-local.  Effective descent in the slice
\(\mathcal X_{\Sp,/S}\) gives
\[
    \Map_S(B,F)
    \simeq
    \Tot_n\Map_S(B'_n,F)
\]
and
\[
    \Map_S(A,F)
    \simeq
    \Tot_n\Map_S(A\times_BB'_n,F).
\]
The degreewise comparison maps are equivalences because the corresponding maps
\(A\times_BB'_n\to B'_n\) are relative de Rham equivalences.  Taking totalizations therefore gives
\[
    \Map_S(B,F)\xrightarrow{\simeq}\Map_S(A,F).
\]
This holds for every relative de Rham-local \(F\), so \(f\) is inverted by \(L_{\dR,/S}\).
\end{proof}

\begin{lemma}[Relative image-locality criterion]
\label{lem:relative-dr-image-locality}
Fix \(S\in\mathcal X_{\Sp}\).  Let
\[
    m:Q\longrightarrow Y
\]
be a monomorphism in \(\mathcal X_{\Sp,/S}\), and suppose that \(Y\) is
\(L_{\dR,/S}\)-local.  Assume that for every affine spectral square-zero immersion
\[
    V\longrightarrow T
\]
over \(S\) and every \(q_0:V\to Q\), the unique extension
\[
    a:T\longrightarrow Y
\]
of \(mq_0\) supplied by locality of \(Y\) factors Nisnevich-locally through \(Q\).  Then \(Q\) is
\(L_{\dR,/S}\)-local.
\end{lemma}

\begin{proof}
Fix an affine spectral square-zero immersion \(V\to T\) over \(S\).  Since \(Y\) is local,
restriction is an equivalence
\[
    \Map_S(T,Y)\xrightarrow{\simeq}\Map_S(V,Y).
\]
For \(a:T\to Y\), let \(a_0\) denote its restriction and define
\[
    \Phi_T(a)
    :=
    \fib_a
    \left(
        \Map_S(T,Q)\longrightarrow\Map_S(T,Y)
    \right),
\]
\[
    \Phi_V(a_0)
    :=
    \fib_{a_0}
    \left(
        \Map_S(V,Q)\longrightarrow\Map_S(V,Y)
    \right).
\]
Because \(m\) is a monomorphism, both fibres are \((-1)\)-truncated and are therefore either empty
or contractible.  Restriction gives
\[
    \Phi_T(a)\longrightarrow\Phi_V(a_0).
\]
If the source is nonempty, then the target is nonempty.  Conversely, suppose
\(\Phi_V(a_0)\) is nonempty and choose the corresponding factorization \(q_0:V\to Q\).  By
hypothesis, \(a\) factors through \(Q\) after a Nisnevich cover of \(T\).  On overlaps, two local
factorizations have the same composite to \(Y\); since \(Q\to Y\) is a monomorphism, the space of
identifications between them is contractible.  The local factorizations therefore descend to a
global factorization \(T\to Q\).  Hence \(\Phi_T(a)\) is nonempty.

Thus
\[
    \Phi_T(a)\xrightarrow{\simeq}\Phi_V(a_0)
\]
for every \(a\).  Together with the equivalence on the \(Y\)-mapping spaces, this gives
\[
    \Map_S(T,Q)\xrightarrow{\simeq}\Map_S(V,Q).
\]
The slice-locality criterion for \(L_{\dR,/S}\) therefore shows that \(Q\) is relative
de Rham-local.
\end{proof}

\begin{theorem}[Relative de Rham preservation of effective epimorphisms]
\label{thm:relative-dr-preserves-effective-epis}
Fix \(S\in\mathcal X_{\Sp}\).  If
\[
    e:A\longrightarrow B
\]
is an effective epimorphism in \(\mathcal X_{\Sp,/S}\), then the induced morphism
\[
\boxed{
    e_{\dR/S}:
    A_{\dR/S}
    \longrightarrow
    B_{\dR/S}
}
\]
is an effective epimorphism in the ambient slice \(\mathcal X_{\Sp,/S}\).  The construction is
compatible with arbitrary base change in \(S\).
\end{theorem}

\begin{proof}
Factor \(e_{\dR/S}\) in the effective-epimorphism--monomorphism factorization system of the
ambient slice:
\[
    A_{\dR/S}
    \longrightarrow
    I
    \xrightarrow{m}
    B_{\dR/S}.
\]
We first prove that \(I\) is \(L_{\dR,/S}\)-local.

Let \(V\to T\) be an affine spectral square-zero immersion over \(S\), and let
\(q_0:V\to I\).  Pulling back the effective epimorphism \(A_{\dR/S}\to I\) gives an effective
epimorphism
\[
    V\times_I A_{\dR/S}\longrightarrow V.
\]
Effective epimorphisms in a slice of an \(\infty\)-topos are detected on the underlying morphisms in
the ambient \(\infty\)-topos.  Thus, after forgetting the structure map to \(S\), the preceding map
is still an effective epimorphism over the affine object \(V\).  The local-splitting theorem for
effective epimorphisms over affines from the preceding chapter therefore supplies a finite affine
spectral Nisnevich basis cover
\[
    \{V_\alpha\longrightarrow V\}_\alpha
\]
and sections
\[
    s_\alpha:V_\alpha\longrightarrow A_{\dR/S}
\]
over \(I\).  One-stage \(\acute{e}\)tale rigidity across the square-zero immersion \(V\to T\)
transports this cover, through a contractible space of compatible choices, to a finite affine
spectral Nisnevich basis cover
\[
    \{T_\alpha\longrightarrow T\}_\alpha
\]
with
\[
    V_\alpha\simeq V\times_TT_\alpha.
\]

The object \(A_{\dR/S}\) is relative de Rham-local, so each \(s_\alpha\) extends through a
contractible space of choices to
\[
    \widetilde s_\alpha:T_\alpha\longrightarrow A_{\dR/S}.
\]
Let
\[
    a:T\longrightarrow B_{\dR/S}
\]
be the unique extension of \(mq_0\).  On \(V_\alpha\), the two maps
\[
    a|_{T_\alpha},
    \qquad
    e_{\dR/S}\widetilde s_\alpha
\]
agree.  Since \(B_{\dR/S}\) is relative de Rham-local, restriction
\[
    \Map_S(T_\alpha,B_{\dR/S})
    \xrightarrow{\simeq}
    \Map_S(V_\alpha,B_{\dR/S})
\]
is an equivalence, so this agreement extends through a contractible space of homotopies over
\(T_\alpha\).  Composing each \(\widetilde s_\alpha\) with
\(A_{\dR/S}\to I\) gives a factorization of \(a|_{T_\alpha}\) through \(I\).
On every iterated overlap, any two such factorizations have the same composite to
\(B_{\dR/S}\); since \(m:I\to B_{\dR/S}\) is a monomorphism, the space of compatible
identifications between them is contractible.  Hence these maps carry the coherent Nisnevich-local
factorization data required by the image-locality criterion.  Thus \(a\) factors
Nisnevich-locally through \(I\).  By
Lemma~\ref{lem:relative-dr-image-locality}, \(I\) is \(L_{\dR,/S}\)-local.

Now form
\[
    J:=B\times_{B_{\dR/S}}I.
\]
The morphism \(J\to B\) is a monomorphism.  Naturality of the relative de Rham unit gives
\[
    e_{\dR/S}\eta_{A/S}\simeq\eta_{B/S}e,
\]
and the left side factors through \(I\).  Hence the effective epimorphism \(e:A\to B\) factors
through the subobject \(J\to B\).  The image of an effective epimorphism is its whole target, so
\(J\to B\) is an equivalence.  Consequently the relative de Rham unit factors as
\[
    B\xrightarrow{j}I\xrightarrow{m}B_{\dR/S}.
\]

Since \(I\) is local, the universal property of the localization unit
\(\eta_{B/S}:B\to B_{\dR/S}\) gives a canonical morphism
\[
    r:B_{\dR/S}\longrightarrow I
\]
with
\[
    r\eta_{B/S}\simeq j.
\]
Therefore
\[
    mr\eta_{B/S}
    \simeq
    mj
    \simeq
    \eta_{B/S}.
\]
Because \(B_{\dR/S}\) is local, precomposition with \(\eta_{B/S}\) is an equivalence on mapping
spaces into \(B_{\dR/S}\), and hence
\[
    mr\simeq\id_{B_{\dR/S}}.
\]
Thus the monomorphism \(m\) is also a split epimorphism, so it is an equivalence.  Therefore the
ambient image of \(e_{\dR/S}\) is all of \(B_{\dR/S}\), proving that \(e_{\dR/S}\) is an
effective epimorphism.

Compatibility with arbitrary base change follows from pullback stability of effective epimorphisms
and the coherent base-change equivalences for relative de Rham objects.
\end{proof}

\begin{corollary}[Relative de Rham \v{C}ech continuity for effective epimorphisms]
\label{cor:relative-dr-cech-effective-epi}
Fix \(S\in\mathcal X_{\Sp}\).  Let \(e:A\to B\) be an effective epimorphism in
\(\mathcal X_{\Sp,/S}\), and let \(A_\bullet\to B\) be its \v{C}ech nerve.  Then
\[
\boxed{
    (A_n)_{\dR/S}
    \simeq
    \underbrace{
        A_{\dR/S}\times_{B_{\dR/S}}\cdots\times_{B_{\dR/S}}A_{\dR/S}
    }_{n+1\text{ factors}}
}
\]
compatibly with all simplicial operators, and
\[
\boxed{
    |(A_\bullet)_{\dR/S}|
    \simeq
    B_{\dR/S}.
}
\]
These equivalences are compatible with arbitrary base change.
\end{corollary}

\begin{proof}
The degreewise equivalence follows from left exactness of \(L_{\dR,/S}\).  By
Theorem~\ref{thm:relative-dr-preserves-effective-epis}, the morphism
\[
    A_{\dR/S}\longrightarrow B_{\dR/S}
\]
is an effective epimorphism in the ambient slice.  The displayed simplicial object is therefore its
\v{C}ech nerve, and its realization is \(B_{\dR/S}\).  Base-change compatibility follows from the
corresponding coherence for relative de Rham objects and pullback stability of effective
epimorphisms.
\end{proof}

\begin{proposition}[Global structured square-zero thickenings are de Rham equivalences]
\label{prop:global-square-zero-dr-equivalence}
Let
\[
    i:Z\longrightarrow T
\]
be a global structured square-zero thickening of spectral schemes of the kind constructed in
Chapter~\ref{chap:cotangent-spectral-deformation}.  Then \(i\) is inverted by the intrinsic de Rham
localization.  More generally, every finite composite of global structured square-zero thickenings is
a de Rham equivalence.  The same assertion holds in every relative slice over \(S\).
\end{proposition}

\begin{proof}
It is enough to prove the assertion in a fixed relative slice over \(S\); the absolute assertion is
the special case of the terminal base.  Choose an affine Zariski cover
\[
    \{T_\alpha\longrightarrow T\}_\alpha.
\]
The global structured square-zero construction of
Chapter~\ref{chap:cotangent-spectral-deformation} identifies each pullback
\[
    Z_\alpha:=Z\times_TT_\alpha\longrightarrow T_\alpha
\]
with an affine structured spectral square-zero immersion.  Hence every
\(Z_\alpha\to T_\alpha\) is inverted by \(L_{\dR,/S}\).

The coproduct
\[
    E:=\coprod_\alpha T_\alpha\longrightarrow T
\]
is an effective epimorphism in the slice, and its pullback along \(Z\to T\) is
\[
    \coprod_\alpha Z_\alpha
    \longrightarrow
    \coprod_\alpha T_\alpha.
\]
This pullback morphism is a relative de Rham equivalence: for every
\(L_{\dR,/S}\)-local object \(F\), mapping out of the two coproducts gives the product of the
equivalences
\[
    \Map_S(T_\alpha,F)
    \xrightarrow{\simeq}
    \Map_S(Z_\alpha,F).
\]
Lemma~\ref{lem:relative-dr-equivalence-target-local} therefore shows that
\[
    Z\longrightarrow T
\]
is a relative de Rham equivalence.  Finite composites remain relative de Rham equivalences by
composition.  This proves both the absolute and relative assertions.
\end{proof}

\begin{proposition}[Factorization of de Rham-equivalent infinitesimal neighbourhoods]
\label{prop:finite-infinitesimal-factorization-completion}
Let
\[
    Z\longrightarrow T\longrightarrow Y\longrightarrow S
\]
be a diagram in which \(Z\to T\) is inverted by the relative de Rham localization over \(S\).
This applies, in particular, when \(Z\to T\) is an affine finite spectral infinitesimal thickening of
the preceding chapter, or when it is a finite composite of global structured square-zero thickenings
as in Proposition~\ref{prop:global-square-zero-dr-equivalence}.  Then the map \(T\to Y\) factors
canonically through the intrinsic completion:
\[
\boxed{
    T\longrightarrow
    \widehat Y^{\,\dR}_{Z/S}
    \longrightarrow
    Y.
}
\]
The factorization is natural in the diagram and compatible with arbitrary base change.
\end{proposition}

\begin{proof}
By hypothesis,
\[
    Z_{\dR/S}\xrightarrow{\simeq}T_{\dR/S}.
\]
Functorial inversion in the maximal \(\infty\)-groupoid of
\(\mathcal X_{\Sp,/S}\) supplies a contractibly canonical inverse
\[
    T_{\dR/S}\xrightarrow{\simeq}Z_{\dR/S}.
\]
Composing it with the relative de Rham unit gives
\[
    T\longrightarrow T_{\dR/S}
    \xrightarrow{\simeq}
    Z_{\dR/S}.
\]
Its composite with \(Z_{\dR/S}\to Y_{\dR/S}\) is the localization of the given map
\(T\to Y\), by naturality of the localization unit.  Thus the pair of maps
\[
    T\to Y,
    \qquad
    T\to Z_{\dR/S}
\]
defines a map to
\[
    Y\times_{Y_{\dR/S}}Z_{\dR/S}
    =
    \widehat Y^{\,\dR}_{Z/S}.
\]
All maps used are canonical transformations of the relative localization, and inversion of
equivalences is functorial on the core.  Thus functoriality and base-change compatibility follow from
the corresponding coherences of the de Rham reflector.
\end{proof}

\begin{theorem}[Infinitesimal simplices as diagonal completions]
\label{thm:relation-simplices-as-completions}
For \(n\geq0\), let
\[
    Y_n:=X^{\times_S(n+1)}
\]
and let
\[
    \Delta_n:X\longrightarrow Y_n
\]
be the diagonal.  Then there is a canonical equivalence over \(Y_n\)
\[
\boxed{
    R_{X/S,n}
    \xrightarrow{\simeq}
    \widehat{Y_n}^{\,\dR}_{X/S}.
}
\]
These equivalences are compatible with the simplicial operators in \(n\), with functoriality in
\(X\to S\), and with arbitrary base change in \(S\).
\end{theorem}

\begin{proof}
Left exactness of the relative de Rham localization gives
\[
    (Y_n)_{\dR/S}
    \simeq
    (X_{\dR/S})^{\times_S(n+1)}.
\]
The relative de Rham image of the diagonal is the diagonal
\[
    X_{\dR/S}
    \longrightarrow
    (X_{\dR/S})^{\times_S(n+1)}.
\]
Therefore
\[
\begin{aligned}
    \widehat{Y_n}^{\,\dR}_{X/S}
    &\simeq
    X^{\times_S(n+1)}
    \times_{(X_{\dR/S})^{\times_S(n+1)}}
    X_{\dR/S}
    \\
    &\simeq
    \underbrace{
        X\times_{X_{\dR/S}}\cdots\times_{X_{\dR/S}}X
    }_{n+1\text{ factors}}
    \\
    &=R_{X/S,n}.
\end{aligned}
\]
Every simplicial operator is obtained by deleting or repeating factors, and the displayed
identification is constructed solely from those product diagrams.  Hence it is compatible with all
simplicial operators.  Functoriality and base change follow from the same finite-limit description
and Proposition~\ref{prop:intrinsic-relative-dr-completion-base-change}.
\end{proof}


\section{Parameterized stable coefficients on the big spectral Nisnevich topos}
\label{sec:monadic-de-rham-parameterized}

The infinitesimal relation objects need not be represented.  We therefore form stable coefficients
inside the slices of the big topos before restricting to any represented or smooth site.

\begin{definition}[Parameterized stable category]
\label{def:monadic-de-rham-parameterized}
For \(T\in\mathcal X_{\Sp}\), define
\[
\boxed{
    \mathcal E_T:=\Sp(\mathcal X_{\Sp,/T}).
}
\]
Write \(\mathbf 1_T\) for its tensor unit.
\end{definition}

\begin{proposition}[Stable presentability and closed fibrewise tensor product]
\label{prop:monadic-de-rham-stable-presentable}
For every \(T\), the category \(\mathcal E_T\) is stable, presentable, and closed symmetric
monoidal.  Its tensor product preserves colimits separately in each variable and admits an internal
Hom
\[
    \underline{\Hom}_{\mathcal E_T}(-,-).
\]
\end{proposition}

\begin{proof}
The slice \(\mathcal X_{\Sp,/T}\) is an \(\infty\)-topos.  Stabilization of its pointed Cartesian
structure is a presentable stable symmetric monoidal category with fibrewise smash product.
Separate colimit preservation of the smash product gives the internal Hom by the adjoint functor
theorem.
\end{proof}

\begin{construction}[Relative suspension spectra]
Let \(a:R\to S\) be a morphism in \(\mathcal X_{\Sp}\).  Write \(R_+\) for the free pointed object
on \(R\) in the slice \(\mathcal X_{\Sp,/S}\), and define
\[
\boxed{
    \Sigma^\infty_{+,S}R
    :=
    \Sigma^\infty_S R_+
    \in\mathcal E_S.
}
\]
For represented \(R\to S\), this is the ordinary fibrewise suspension spectrum of the object
\(R\) of the slice.  The stable dependent-sum theorem below identifies it canonically with
\(\Sigma_a\mathbf 1_R\).
\end{construction}

\begin{theorem}[Affine-point presentation of slices and stable slices]
\label{thm:monadic-de-rham-affine-slice-presentation}
Let \(T\in\mathcal X_{\Sp}\), and let \(\mathcal C^{\mathrm{aff}}_{/T}\) be the
\(\infty\)-category of affine points
\[
    t:h_U\longrightarrow T,
    \qquad U\in\operatorname{AffSp}.
\]
Equip \(\mathcal C^{\mathrm{aff}}_{/T}\) with the topology in which a family over
\((h_U\to T)\) is covering precisely when the underlying family of affine spectral schemes is an
affine spectral Nisnevich cover.  Restriction to affine points induces canonical equivalences
\[
\boxed{
    \mathcal X_{\Sp,/T}
    \xrightarrow{\simeq}
    \operatorname{Shv}_{\Nis}
    \left(
        \mathcal C^{\mathrm{aff}}_{/T};\mathcal S
    \right)
}
\]
and
\[
\boxed{
    \mathcal E_T
    \xrightarrow{\simeq}
    \operatorname{Shv}_{\Nis}
    \left(
        \mathcal C^{\mathrm{aff}}_{/T};\Sp
    \right).
}
\]
For an affine point \(t:h_U\to T\) and \(E\in\mathcal E_T\), stable Yoneda gives
\[
\boxed{
    \map_{\mathcal E_T}
    \left(
        \Sigma^\infty_{+,T}U,E
    \right)
    \simeq
    E(U,t).
}
\]
The objects
\[
    \Sigma^a\Sigma^\infty_{+,T}U,
    \qquad
    (h_U\to T)\in\mathcal C^{\mathrm{aff}}_{/T},\ a\in\mathbb Z,
\]
generate \(\mathcal E_T\) under small colimits.
\end{theorem}

\begin{proof}
For \(Y\to T\) define a presheaf on affine points by
\[
    \Phi_T(Y)(U,t)
    :=
    \Map_T(h_U,Y)
    \simeq
    \fib_t\bigl(Y(U)\to T(U)\bigr).
\]
Both \(Y\) and \(T\) are affine Nisnevich sheaves, and fibres commute with limits.  Hence
\(\Phi_T(Y)\) satisfies the displayed affine-point Nisnevich descent condition.

Conversely, let \(G\) be a sheaf on \(\mathcal C^{\mathrm{aff}}_{/T}\).  Define a presheaf
\(\Psi_T(G)\) on affine spectral schemes by taking, over each affine \(U\), the total space of the
space-valued family
\[
    t\longmapsto G(U,t),
    \qquad t\in T(U).
\]
Thus there is a natural morphism
\[
    \Psi_T(G)(U)\longrightarrow T(U)
\]
whose fibre over \(t\) is \(G(U,t)\).  Let \(U_\bullet\to U\) be an affine spectral Nisnevich
\v{C}ech nerve.  Since \(T\) is a sheaf,
\[
    T(U)\xrightarrow{\simeq}\Tot T(U_\bullet).
\]
Over a point \(t\in T(U)\), the fibre of the descent comparison for \(\Psi_T(G)\) is precisely
\[
    G(U,t)
    \longrightarrow
    \Tot G(U_\bullet,t|_{U_\bullet}),
\]
which is an equivalence because \(G\) is a sheaf.  Therefore \(\Psi_T(G)\to T\) is an object of
\(\mathcal X_{\Sp,/T}\).  Taking fibres and taking total spaces are inverse constructions, so
\(\Phi_T\) and \(\Psi_T\) are inverse equivalences.

For the stable statement, use the reduced-excisive description of spectrum objects.  Spectrum
objects are defined by finite-limit conditions in a pointed functor category, and the affine-point
sheaf condition is also a limit condition computed objectwise.  Stabilizing the preceding
equivalence therefore gives
\[
    \Sp(\mathcal X_{\Sp,/T})
    \simeq
    \operatorname{Shv}_{\Nis}
    \left(
        \mathcal C^{\mathrm{aff}}_{/T};\Sp
    \right).
\]
Under this equivalence the relative suspension spectrum of an affine point is the suspension-spectrum
Yoneda object.  Stable Yoneda gives the mapping-spectrum formula.  Spectrum-valued sheaves are
generated under colimits by the integer suspensions of the sheafified suspension-spectrum
representables, which proves the final assertion.
\end{proof}

\begin{theorem}[Stable dependent sum, pullback, and dependent product]
\label{thm:monadic-de-rham-stable-adjunction}
Every morphism \(f:T\to S\) induces an exact strong symmetric monoidal pullback
\[
    f^*:\mathcal E_S\longrightarrow\mathcal E_T
\]
which preserves all limits and colimits.  It admits exact adjoints
\[
\boxed{
    \Sigma_f\dashv f^*\dashv \Pi_f.
}
\]
The dependent product \(\Pi_f\) is canonically lax symmetric monoidal.  There is a canonical projection
equivalence
\[
\boxed{
    \Sigma_f\bigl(E\otimes f^*F\bigr)
    \xrightarrow{\simeq}
    \Sigma_f E\otimes F
}
\]
for \(E\in\mathcal E_T\) and \(F\in\mathcal E_S\).  Identity and composition carry their canonical
coherent comparison equivalences.
\end{theorem}

\begin{proof}
Pullback between slices of an \(\infty\)-topos is accessible and preserves all limits and colimits.
It is strong Cartesian monoidal.  Passing to spectrum objects therefore gives an accessible exact
strong symmetric monoidal functor
\[
    f^*:\mathcal E_S\longrightarrow\mathcal E_T
\]
which still preserves all limits and colimits.  The presentable adjoint functor theorem supplies
both adjoints
\[
    \Sigma_f\dashv f^*\dashv \Pi_f.
\]
All three functors are exact because the source and target categories are stable, and the right
adjoint of a strong symmetric monoidal functor carries its canonical lax symmetric monoidal
structure.

We first record the effect of \(\Sigma_f\) on relative suspension spectra.  Let \(r:R\to T\) and
regard \(R\) over \(S\) through \(fr\).  For every \(G\in\mathcal E_S\),
\[
\begin{aligned}
    \Map_{\mathcal E_S}
    \left(
        \Sigma_f\Sigma^\infty_{+,T}R,G
    \right)
    &\simeq
    \Map_{\mathcal E_T}
    \left(
        \Sigma^\infty_{+,T}R,f^*G
    \right)
    \\
    &\simeq
    \Map_{\mathcal X_{\Sp,/T}}
    \left(
        R,\Omega^\infty f^*G
    \right)
    \\
    &\simeq
    \Map_{\mathcal X_{\Sp,/S}}
    \left(
        R,\Omega^\infty G
    \right)
    \\
    &\simeq
    \Map_{\mathcal E_S}
    \left(
        \Sigma^\infty_{+,S}R,G
    \right).
\end{aligned}
\]
Hence
\[
    \Sigma_f\Sigma^\infty_{+,T}R
    \simeq
    \Sigma^\infty_{+,S}R.
\]
Likewise pullback gives
\[
    f^*\Sigma^\infty_{+,S}Y
    \simeq
    \Sigma^\infty_{+,T}(Y\times_ST).
\]

There is a canonical projection morphism
\[
    \Sigma_f(E\otimes f^*F)
    \longrightarrow
    \Sigma_f E\otimes F
\]
obtained by adjunction from the counit and the strong monoidality of \(f^*\).  Both sides preserve
small colimits separately in \(E\) and \(F\).  By
Theorem~\ref{thm:monadic-de-rham-affine-slice-presentation}, it is therefore enough to evaluate on
shifted relative suspension spectra.  Let
\[
    E=\Sigma^a\Sigma^\infty_{+,T}R,
    \qquad
    F=\Sigma^b\Sigma^\infty_{+,S}Y.
\]
Then
\[
\begin{aligned}
    \Sigma_f(E\otimes f^*F)
    &\simeq
    \Sigma^{a+b}
    \Sigma_f\Sigma^\infty_{+,T}
    \left(
        R\times_T(Y\times_ST)
    \right)
    \\
    &\simeq
    \Sigma^{a+b}
    \Sigma^\infty_{+,S}(R\times_SY),
\end{aligned}
\]
while
\[
\begin{aligned}
    \Sigma_f E\otimes F
    &\simeq
    \Sigma^{a+b}
    \Sigma^\infty_{+,S}R
    \otimes
    \Sigma^\infty_{+,S}Y
    \\
    &\simeq
    \Sigma^{a+b}
    \Sigma^\infty_{+,S}(R\times_SY).
\end{aligned}
\]
Under these identifications the projection morphism is the identity.  It is therefore an equivalence
for arbitrary \(E\) and \(F\).  The coherent identity and composition structures come from the
parameterized adjunctions and their mate calculus.
\end{proof}

\begin{theorem}[Stable Beck--Chevalley]
\label{thm:monadic-de-rham-beck-chevalley}
For every Cartesian square
\[
\begin{tikzcd}[column sep=large,row sep=large]
    T' \ar[r,"g'"] \ar[d,"f'"'] & T \ar[d,"f"]\\
    S' \ar[r,"g"'] & S
\end{tikzcd}
\]
there are canonical equivalences
\[
\boxed{
    g^*\Sigma_f\xrightarrow{\simeq}\Sigma_{f'} g'^*,
    \qquad
    g^*\Pi_f\xrightarrow{\simeq}\Pi_{f'} g'^*.
}
\]
They satisfy identity, composition, horizontal and vertical pasting, and all higher
Beck--Chevalley coherences.
\end{theorem}

\begin{proof}
The left Beck--Chevalley transformation is the canonical mate
\[
    g^*\Sigma_f\longrightarrow \Sigma_{f'} g'^*.
\]
Both sides preserve small colimits.  It is therefore enough to evaluate on a shifted affine-point
generator \(\Sigma^a\Sigma^\infty_{+,T}R\).  Using the relative suspension-spectrum formulas from
the preceding proof, both sides identify canonically with
\[
    \Sigma^a\Sigma^\infty_{+,S'}(R\times_SS').
\]
Under these identifications the mate is the identity, so
\[
    g^*\Sigma_f\xrightarrow{\simeq}\Sigma_{f'} g'^*.
\]

For the right Beck--Chevalley transformation, let \(E\in\mathcal E_T\) and
\(K\in\mathcal E_{S'}\).  Applying the left Beck--Chevalley equivalence to the same Cartesian
square with the two legs interchanged gives
\[
    f^*\Sigma_g K\simeq \Sigma_{g'} f'^*K.
\]
Hence
\[
\begin{aligned}
    \map_{\mathcal E_{S'}}(K,g^*\Pi_f E)
    &\simeq
    \map_{\mathcal E_S}(\Sigma_g K,\Pi_f E)
    \\
    &\simeq
    \map_{\mathcal E_T}(f^*\Sigma_g K,E)
    \\
    &\simeq
    \map_{\mathcal E_T}(\Sigma_{g'} f'^*K,E)
    \\
    &\simeq
    \map_{\mathcal E_{T'}}(f'^*K,g'^*E)
    \\
    &\simeq
    \map_{\mathcal E_{S'}}(K,\Pi_{f'} g'^*E).
\end{aligned}
\]
Yoneda gives
\[
    g^*\Pi_f\xrightarrow{\simeq}\Pi_{f'} g'^*.
\]
All transformations used are canonical mates of pullback-pasting equivalences.  Their identity,
composition, horizontal and vertical pasting, and higher coherences therefore follow from the
coherent associativity of pullback and the functoriality of taking mates.
\end{proof}

\begin{proposition}[Closed base change]
\label{prop:monadic-de-rham-closed-base-change}
For every \(f:T\to S\) and \(A,B\in\mathcal E_S\), the canonical map
\[
\boxed{
    f^*\underline{\Hom}_{\mathcal E_S}(A,B)
    \xrightarrow{\simeq}
    \underline{\Hom}_{\mathcal E_T}(f^*A,f^*B)
}
\]
is an equivalence, coherently in \(f,A,B\).
\end{proposition}

\begin{proof}
Let \(K\in\mathcal E_T\).  Using the adjunction \(\Sigma_f\dashv f^*\), closedness, and the projection
formula already proved, there are natural equivalences
\[
\begin{aligned}
    \map_{\mathcal E_T}
    \left(
        K,f^*\underline{\Hom}_{\mathcal E_S}(A,B)
    \right)
    &\simeq
    \map_{\mathcal E_S}
    \left(
        \Sigma_f K,\underline{\Hom}_{\mathcal E_S}(A,B)
    \right)
    \\
    &\simeq
    \map_{\mathcal E_S}(\Sigma_f K\otimes A,B)
    \\
    &\simeq
    \map_{\mathcal E_S}
    \left(
        \Sigma_f(K\otimes f^*A),B
    \right)
    \\
    &\simeq
    \map_{\mathcal E_T}(K\otimes f^*A,f^*B)
    \\
    &\simeq
    \map_{\mathcal E_T}
    \left(
        K,\underline{\Hom}_{\mathcal E_T}(f^*A,f^*B)
    \right).
\end{aligned}
\]
Yoneda gives the displayed equivalence.  Naturality and coherence follow from the coherent
projection formula and adjunction data.
\end{proof}

\begin{definition}[Stable sections]
\label{def:monadic-de-rham-stable-sections}
For \(E\in\mathcal E_T\), define
\[
    \Gamma_T^{\mathrm{st}}(E)
    :=
    \map_{\mathcal E_T}(\mathbf 1_T,E).
\]
\end{definition}

\begin{lemma}[Stabilization commutes with left-exact limits]
\label{lem:monadic-de-rham-stabilization-limits}
Let \(I\) be a small \(\infty\)-category and let
\[
    \mathcal Y\simeq\limop_{i\in I}\mathcal Y_i
\]
be a limit of \(\infty\)-categories with finite limits whose transition functors preserve finite
limits.  Then restriction induces a canonical equivalence
\[
\boxed{
    \Sp(\mathcal Y)
    \xrightarrow{\simeq}
    \limop_{i\in I}\Sp(\mathcal Y_i).
}
\]
If the \(\mathcal Y_i\) are \(\infty\)-topoi and the transition functors are inverse-image
functors, the equivalence is compatible with the fibrewise symmetric monoidal structures.
\end{lemma}

\begin{proof}
Use the reduced-excisive model
\[
    \Sp(\mathcal Y)
    \simeq
    \operatorname{Exc}_*
    \left(
        \mathcal S_*^{\mathrm{fin}},\mathcal Y
    \right).
\]
Functor categories commute with limits in the target.  The conditions of carrying the basepoint to
a terminal object and carrying finite homotopy pushout squares of finite pointed spaces to pullback
squares are finite-limit conditions, so they are detected componentwise in
\(\limop_i\mathcal Y_i\).  Hence the reduced-excisive subcategory is the corresponding limit of the
reduced-excisive subcategories.  For inverse-image functors of \(\infty\)-topoi, the Cartesian
products and their stabilized fibrewise smash products are preserved degreewise, which gives the
monoidal compatibility.
\end{proof}

\begin{theorem}[Effective descent for parameterized spectra]
\label{thm:monadic-de-rham-effective-descent}
Let \(e:X\to Y\) be an effective epimorphism in \(\mathcal X_{\Sp}\), and let
\(X_\bullet\to Y\) be its \v{C}ech nerve.  Restriction induces a symmetric monoidal equivalence
\[
\boxed{
    \mathcal E_Y
    \xrightarrow{\simeq}
    \Tot_{[n]\in\Delta}\mathcal E_{X_n}.
}
\]
The equivalence is compatible with arbitrary base change.
\end{theorem}

\begin{proof}
Effective descent in an \(\infty\)-topos gives an equivalence of slices
\[
    \mathcal X_{\Sp,/Y}
    \xrightarrow{\simeq}
    \Tot_{[n]\in\Delta}\mathcal X_{\Sp,/X_n}.
\]
Every transition functor in this limit is pullback and therefore preserves finite limits.  Applying
Lemma~\ref{lem:monadic-de-rham-stabilization-limits} gives
\[
    \Sp(\mathcal X_{\Sp,/Y})
    \xrightarrow{\simeq}
    \Tot_{[n]\in\Delta}
    \Sp(\mathcal X_{\Sp,/X_n}),
\]
which is the displayed equivalence.  It is symmetric monoidal by the last assertion of the lemma.
Effective epimorphisms and their \v{C}ech nerves are stable under arbitrary pullback, and the slice
descent equivalence is compatible with pullback.  Applying the same stabilization comparison gives
the asserted base-change compatibility.
\end{proof}

\begin{proposition}[Parameterized function-spectrum identity]
\label{prop:parameterized-function-spectrum-identity}
For every \(a:R\to S\) and \(E\in\mathcal E_S\), there is a canonical equivalence
\[
\boxed{
    \Pi_a a^*E
    \xrightarrow{\simeq}
    \underline{\Hom}_{\mathcal E_S}
    \left(
        \Sigma^\infty_{+,S}R,E
    \right).
}
\]
It is natural in \(a\) and \(E\) and compatible with arbitrary base change.
\end{proposition}

\begin{proof}
Let \(K\in\mathcal E_S\).  The projection formula gives
\[
    K\otimes\Sigma^\infty_{+,S}R
    =
    K\otimes \Sigma_a\mathbf 1_R
    \simeq
    \Sigma_a a^*K.
\]
Hence
\[
\begin{aligned}
    \map_{\mathcal E_S}
    \left(
        K,
        \underline{\Hom}_{\mathcal E_S}
        (\Sigma^\infty_{+,S}R,E)
    \right)
    &\simeq
    \map_{\mathcal E_S}
    \left(
        K\otimes\Sigma^\infty_{+,S}R,E
    \right)
    \\
    &\simeq
    \map_{\mathcal E_S}
    \left(
        \Sigma_a a^*K,E
    \right)
    \\
    &\simeq
    \map_{\mathcal E_R}
    \left(
        a^*K,a^*E
    \right)
    \\
    &\simeq
    \map_{\mathcal E_S}
    \left(
        K,\Pi_a a^*E
    \right).
\end{aligned}
\]
Yoneda gives the equivalence.  Base-change compatibility follows from
Theorem~\ref{thm:monadic-de-rham-beck-chevalley}, the base-change compatibility of relative
suspension spectra, and Proposition~\ref{prop:monadic-de-rham-closed-base-change}.
\end{proof}


\section{The additive stabilization of the brave-new affine line}
\label{sec:additive-scalar-spectrum}

The brave-new affine line is not redefined here.  For represented spectral \(S\), retain
\[
    \mathbb A^1_S
    =
    \Spec_S\Sym_{\mathcal O_S}(\mathcal O_S)
\]
from the foundational chapters.  For arbitrary \(S\in\mathcal X_{\Sp}\), the same notation denotes
its pullback from the absolute brave-new affine line to the slice over \(S\).

\begin{construction}[Additive ring structure on the existing affine line]
\label{constr:additive-ring-affine-line}
The already constructed zero and unit sections and multiplication are retained.  For represented
\(S\), the relative affine-line universal property gives, naturally in every represented
\(T=\Spec A\to S\),
\[
    \Map_S(T,\mathbb A^1_S)
    \simeq
    \Omega^\infty A.
\]
The right side is canonically a commutative ring space: addition is the infinite-loop addition of the
underlying spectrum of \(A\), multiplication is the multiplication of the \(E_\infty\)-ring \(A\),
and additive inversion is induced by multiplication by \(-1\).  Naturality in \(A\) and affine
density therefore determine, by Yoneda and descent, a canonical commutative ring object structure on
the existing \(\mathbb A^1_S\).

Equivalently, on the universal coordinate the new additive operations are represented
contravariantly by
\[
    t\longmapsto t_1+t_2,
    \qquad
    t\longmapsto -t,
\]
while the retained operations are represented by
\[
    t\longmapsto 0,
    \qquad
    t\longmapsto t_1t_2,
    \qquad
    t\longmapsto1.
\]
More generally, every polynomial expression
\[
    P\in\mathbb Z[t_1,\ldots,t_n]
\]
defines a canonical map
\[
    P:\mathbb A_S^n\longrightarrow\mathbb A_S^1,
\]
and substitution of polynomial expressions agrees with composition.  The additive coherences are inherited from the grouplike \(E_\infty\)-structure on the underlying
spectrum, the multiplicative coherences from the commutative algebra structure, and distributivity
from the coherent separate additivity of multiplication.  Together these structures make the existing
affine line a commutative ring object.  For arbitrary \(S\in\mathcal X_{\Sp}\), pull back the
absolute construction to the slice over \(S\).  All operations commute coherently with arbitrary
base change.
\end{construction}

\begin{lemma}[Multi-additive tensor and additive stabilization]
\label{lem:additive-stabilization-bilinear}
Let \(\mathcal Y\) be an \(\infty\)-topos.  Transport the connective smash product across the
canonical equivalence
\[
    \CGrp(\mathcal Y)
    \simeq
    \Sp_{\geq0}(\mathcal Y).
\]
Write the resulting tensor product on commutative group objects as \(\otimes_{\add}\).  For every
finite tuple of commutative group objects \(G_1,\ldots,G_r,K\), there is a natural equivalence
\[
    \Map_{\CGrp(\mathcal Y)}
    \left(
        G_1\otimes_{\add}\cdots\otimes_{\add}G_r,
        K
    \right)
    \simeq
    \operatorname{MultiAdd}_{\mathcal Y}(G_1,\ldots,G_r;K),
\]
where the right side is the space of maps
\(G_1\times\cdots\times G_r\to K\) which are morphisms of commutative group objects separately in
each variable, together with their coherent multi-additivity data.  Under this tensor product the
recognition equivalence is symmetric monoidal and therefore induces an equivalence
\[
\boxed{
    \CAlg\bigl(\CGrp(\mathcal Y),\otimes_{\add}\bigr)
    \xrightarrow{\simeq}
    \CAlg\bigl(\Sp_{\geq0}(\mathcal Y)\bigr).
}
\]
Consequently every commutative ring object \(A\) of \(\mathcal Y\), whose additive law is already
grouplike, determines canonically a connective commutative algebra object
\[
    \mathbb A^{\add}\in\CAlg(\Sp(\mathcal Y))
\]
whose infinite-loop commutative group object is the additive group underlying \(A\).
\end{lemma}

\begin{proof}
Transport the connective smash product on \(\Sp_{\geq0}(\mathcal Y)\) across the recognition
equivalence.  By the universal property of the smash product of connective spectrum objects, the
transported \(r\)-fold tensor corepresents coherently multi-additive maps of commutative group
objects.  This gives the displayed multi-additive mapping-space formula.  By construction the
recognition equivalence is symmetric monoidal for this transported tensor, and hence it induces the
displayed equivalence on commutative algebra objects.

Let \(A\) be a commutative ring object.  Its multiplication is additive in each variable, and the
iterated multiplications are coherently multi-additive.  The multi-additive universal property
therefore factors every operation of the commutative operad uniquely through the corresponding
\(\otimes_{\add}\)-tensor power of the additive commutative group of \(A\).  Associativity,
commutativity, unitality, distributivity, and all higher operadic coherences are inherited from the
commutative-ring object.  Thus the additive commutative group of \(A\) is canonically an object of
\(\CAlg(\CGrp(\mathcal Y),\otimes_{\add})\).  Transport across the symmetric monoidal recognition
equivalence gives the asserted connective commutative algebra spectrum object.
\end{proof}

\begin{construction}[Additive scalar spectrum]
\label{constr:additive-scalar-spectrum}
The canonical equivalence
\[
    \Sp_{\geq0}(\mathcal X_{\Sp,/S})
    \simeq
    \CGrp(\mathcal X_{\Sp,/S})
\]
identifies connective spectrum objects with commutative group objects.  Define
\[
\boxed{
    \mathbb O_S^{\add}
    \in
    \Sp_{\geq0}(\mathcal X_{\Sp,/S})
}
\]
to be the connective spectrum object corresponding to the additive group underlying
\(\mathbb A^1_S\):
\[
\boxed{
    \Omega^\infty\mathbb O_S^{\add}
    \simeq
    \mathbb A^1_S
}
\]
as commutative group objects.

Multiplication on \(\mathbb A^1_S\) is additive in each variable.  By
Lemma~\ref{lem:additive-stabilization-bilinear}, it induces
\[
    \mathbb O_S^{\add}\otimes\mathbb O_S^{\add}
    \longrightarrow
    \mathbb O_S^{\add}.
\]
The multiplicative identity section is a global point
\[
    1:S\longrightarrow
    \mathbb A^1_S
    \simeq
    \Omega^\infty\mathbb O_S^{\add}
\]
in the slice.  The suspension-spectrum/infinite-loop adjunction gives a canonical equivalence of
spaces
\[
\boxed{
    \Omega^\infty
    \map_{\mathcal E_S}(\mathbf 1_S,\mathbb O_S^{\add})
    \simeq
    \Map_{\mathcal X_{\Sp,/S}}
    \left(S,\Omega^\infty\mathbb O_S^{\add}\right).
}
\]
The section \(1\) therefore corresponds to the unit morphism
\[
    \mathbf 1_S\longrightarrow\mathbb O_S^{\add}.
\]
Together with the multiplicative maps above, the commutative-ring coherence of
\(\mathbb A^1_S\) supplies the commutative-algebra coherence.  Hence
\[
\boxed{
    \mathbb O_S^{\add}\in\CAlg(\mathcal E_S).
}
\]
\end{construction}

\begin{proposition}[Existence and base change]
\label{prop:additive-scalar-spectrum-base-change}
The commutative algebra \(\mathbb O_S^{\add}\) is connective and functorial in \(S\).  For every
\(c:S'\to S\), there is a canonical equivalence
\[
\boxed{
    c^*\mathbb O_S^{\add}
    \xrightarrow{\simeq}
    \mathbb O_{S'}^{\add}
}
\]
in \(\CAlg(\mathcal E_{S'})\), coherently compatible with identities, composition, and pullback
pasting.
\end{proposition}

\begin{proof}
Base change of the brave-new affine line and of its zero, unit, multiplication, and universal
coordinate was established in the preceding spectral geometry.  The additive and inverse maps are
defined by the same relative free-algebra universal property and therefore commute with base change
as well.  The equivalence between commutative group objects and connective spectrum objects is
natural for left-exact functors, and the bilinear tensor comparison of
Lemma~\ref{lem:additive-stabilization-bilinear} is natural for such pullbacks.  Since parameterized
pullback is strong symmetric monoidal, it identifies the entire commutative algebra structures.
\end{proof}

\begin{remark}[Not an unreduced suspension spectrum]
The coefficient \(\mathbb O_S^{\add}\) is the connective spectrum object associated with the already
grouplike additive structure of \(\mathbb A^1_S\).  It is not
\(\Sigma^\infty_{+,S}\mathbb A^1_S\), which is the free stable object on the underlying object of
the slice.
\end{remark}

\begin{proposition}[Scalar sections on represented spectral schemes]
\label{prop:additive-scalar-represents-functions}
Let \(Y\) be represented by a spectral scheme and let \(y:Y\to S\) be its given morphism.  There is a canonical equivalence of \(E_\infty\)-ring spectra
\[
\boxed{
    \map_{\mathcal E_Y}
    \left(
        \mathbf 1_Y,
        y^*\mathbb O_S^{\add}
    \right)
    \simeq
    \Gamma(Y,\mathcal O_Y).
}
\]
It is natural for represented morphisms and is compatible with the canonical exchange morphisms on
derived global sections.  No arbitrary base-change equivalence for global sections is asserted.
\end{proposition}

\begin{proof}
By Proposition~\ref{prop:additive-scalar-spectrum-base-change},
\[
    y^*\mathbb O_S^{\add}\simeq\mathbb O_Y^{\add}.
\]
Let \(r:U=\Spec A\to Y\) be an affine point.  Stable Yoneda from
Theorem~\ref{thm:monadic-de-rham-affine-slice-presentation} gives
\[
    \map_{\mathcal E_Y}
    \left(
        \Sigma^\infty_{+,Y}U,
        \mathbb O_Y^{\add}
    \right)
    \simeq
    \mathbb O_Y^{\add}(U,r).
\]
The right side is the connective spectrum obtained by stabilizing the additive group of sections of
\(\mathbb A^1_U\).  The relative affine-line universal property identifies that additive group with
\(\Omega^\infty A\), and grouplike recognition therefore gives a canonical equivalence
\[
    \mathbb O_Y^{\add}(U,r)\simeq A
\]
of connective spectra.  The polynomial construction of the ring operations identifies this as an
equivalence of \(E_\infty\)-rings.

Yoneda density for the affine-point presentation of the slice gives
\[
    Y
    \simeq
    \colim_{(U=\Spec A\to Y)}U.
\]
Applying the relative suspension-spectrum functor yields
\[
    \mathbf 1_Y
    \simeq
    \colim_{(U=\Spec A\to Y)}
    \Sigma^\infty_{+,Y}U.
\]
Consequently
\[
\begin{aligned}
    \map_{\mathcal E_Y}
    (\mathbf 1_Y,\mathbb O_Y^{\add})
    &\simeq
    \limop_{(U=\Spec A\to Y)\in(\mathcal C^{\mathrm{aff}}_{/Y})^{\op}}
    \map_{\mathcal E_Y}
    \left(
        \Sigma^\infty_{+,Y}U,
        \mathbb O_Y^{\add}
    \right)
    \\
    &\simeq
    \limop_{(U=\Spec A\to Y)\in(\mathcal C^{\mathrm{aff}}_{/Y})^{\op}}A.
\end{aligned}
\]

By the affine-point definition of quasi-coherent modules in Chapter~1,
\[
    \operatorname{QCoh}(Y)
    \simeq
    \limop_{(U=\Spec A\to Y)^{\op}}\operatorname{Mod}_A,
\]
and the structure sheaf corresponds to the compatible family of regular modules \(A\).  Mapping
spectra in a limit of stable \(\infty\)-categories are computed as limits, so
\[
\begin{aligned}
    \Gamma(Y,\mathcal O_Y)
    &=
    \map_{\operatorname{QCoh}(Y)}
    (\mathcal O_Y,\mathcal O_Y)
    \\
    &\simeq
    \limop_{(U=\Spec A\to Y)\in(\mathcal C^{\mathrm{aff}}_{/Y})^{\op}}
    \map_{\operatorname{Mod}_A}(A,A)
    \\
    &\simeq
    \limop_{(U=\Spec A\to Y)\in(\mathcal C^{\mathrm{aff}}_{/Y})^{\op}}A.
\end{aligned}
\]
The two displayed limits use the same affine-point transition maps, so they give the asserted
canonical equivalence.  Since every affine comparison is multiplicative and limits of commutative
algebras are created in spectra, the equivalence is one of \(E_\infty\)-ring spectra.  Naturality
and compatibility with the canonical exchange morphisms follow from functoriality of affine points,
pullback, and the two limit descriptions.
\end{proof}


\section{Infinitesimal de Rham cochains}
\label{sec:cosimplicial-parameterized-de-rham}

Write
\[
    b=b_{X/S}:Q_{X/S}\longrightarrow S
\]
and retain
\[
    a_n:R_{X/S,n}\longrightarrow S.
\]

\begin{definition}[Geometric de Rham cochains]
\label{def:degreewise-modal-de-rham-cochains}
For \(n\geq0\), define
\[
\boxed{
    \underline{\mathbf C}_{\dR}^{\,n}(X/S)
    :=
    \Pi_{a_n}a_n^*\mathbb O_S^{\add}
    \in
    \mathcal E_S.
}
\]
For a simplicial operator
\[
    r_\alpha:R_{X/S,m}\longrightarrow R_{X/S,n},
\]
the adjunction unit
\[
    a_n^*\mathbb O_S^{\add}
    \longrightarrow
    \Pi_{r_\alpha}r_\alpha^*a_n^*\mathbb O_S^{\add}
\]
and the identity \(a_m=a_nr_\alpha\) induce
\[
    \underline{\mathbf C}_{\dR}^{\,n}(X/S)
    \longrightarrow
    \underline{\mathbf C}_{\dR}^{\,m}(X/S).
\]
Thus
\[
\boxed{
    \underline{\mathbf C}_{\dR}^{\,\bullet}(X/S)
    \in
    \Fun(\Delta,\mathcal E_S)
}
\]
is a cosimplicial object.
\end{definition}

\begin{proposition}[Cosimplicial commutative algebra]
\label{prop:monadic-de-rham-degreewise-calg}
The geometric de Rham cochains refine canonically to
\[
\boxed{
    \underline{\mathbf C}_{\dR}^{\,\bullet}(X/S)
    \in
    \Fun\bigl(\Delta,\CAlg(\mathcal E_S)\bigr).
}
\]
\end{proposition}

\begin{proof}
Each \(a_n^*\) is strong symmetric monoidal and each \(\Pi_{a_n}\) is lax symmetric monoidal, so
\(\Pi_{a_n}a_n^*\mathbb O_S^{\add}\) is a commutative algebra.  The adjunction units used to define
the cosimplicial maps are monoidal natural transformations.  Hence every coface and codegeneracy is
a morphism of commutative algebras.
\end{proof}

\begin{proposition}[Mapping-spectrum interpretation]
\label{prop:modal-de-rham-mapping-spectrum}
For every \(n\),
\[
\boxed{
    \Gamma_S^{\mathrm{st}}
    \left(
        \underline{\mathbf C}_{\dR}^{\,n}(X/S)
    \right)
    \simeq
    \map_{\mathcal E_{R_{X/S,n}}}
    \left(
        \mathbf 1_{R_{X/S,n}},
        a_n^*\mathbb O_S^{\add}
    \right).
}
\]
Under these equivalences the cosimplicial maps are pullback of scalar-valued functions along the
simplicial operators of \(R_{X/S,\bullet}\).
\end{proposition}

\begin{proof}
Adjunction and \(a_n^*\mathbf 1_S\simeq\mathbf 1_{R_{X/S,n}}\) give
\[
\begin{aligned}
    \map_{\mathcal E_S}
    \left(
        \mathbf 1_S,\Pi_{a_n}a_n^*\mathbb O_S^{\add}
    \right)
    &\simeq
    \map_{\mathcal E_{R_{X/S,n}}}
    \left(
        a_n^*\mathbf 1_S,a_n^*\mathbb O_S^{\add}
    \right)
    \\
    &\simeq
    \map_{\mathcal E_{R_{X/S,n}}}
    \left(
        \mathbf 1_{R_{X/S,n}},a_n^*\mathbb O_S^{\add}
    \right).
\end{aligned}
\]
The mate description of the adjunction unit identifies the structure maps with pullback.
\end{proof}

\begin{construction}[Suspension spectra of infinitesimal simplices]
\label{constr:suspension-modal-nerve}
Define
\[
\boxed{
    A_n(X/S)
    :=
    \Sigma^\infty_{+,S}R_{X/S,n}
    \in\mathcal E_S.
}
\]
The simplicial structure of \(R_{X/S,\bullet}\) gives
\[
    A_\bullet(X/S)\in\Fun(\Delta^{\op},\mathcal E_S).
\]
Each \(A_n(X/S)\) has its canonical cocommutative coalgebra structure induced by the diagonal of
\(R_{X/S,n}\).
\end{construction}

\begin{theorem}[Internal-function presentation]
\label{thm:modal-cochains-internal-functions}
For every \(n\), there is a canonical equivalence
\[
\boxed{
    \underline{\mathbf C}_{\dR}^{\,n}(X/S)
    \simeq
    \underline{\Hom}_{\mathcal E_S}
    \left(
        A_n(X/S),
        \mathbb O_S^{\add}
    \right).
}
\]
These equivalences assemble cosimplicially.  The degreewise commutative multiplication corresponds
to convolution against the cocommutative diagonal of \(A_n(X/S)\).
\end{theorem}

\begin{proof}
Apply Proposition~\ref{prop:parameterized-function-spectrum-identity} to
\(a_n:R_{X/S,n}\to S\) and \(E=\mathbb O_S^{\add}\):
\[
    \Pi_{a_n}a_n^*\mathbb O_S^{\add}
    \simeq
    \underline{\Hom}_{\mathcal E_S}
    \left(
        \Sigma^\infty_{+,S}R_{X/S,n},
        \mathbb O_S^{\add}
    \right).
\]
Naturality of the parameterized function-spectrum identity gives the cosimplicial compatibility.
Under the adjunction, the lax-monoidal direct-image multiplication is the convolution product
obtained from
\[
    A_n\longrightarrow A_n\otimes A_n
\]
and multiplication on \(\mathbb O_S^{\add}\).
\end{proof}

\begin{proposition}[Represented function-spectrum formula]
\label{prop:represented-modal-cochain-functions}
If \(R_{X/S,n}\) is represented, then
\[
\boxed{
    \Gamma_S^{\mathrm{st}}
    \left(
        \underline{\mathbf C}_{\dR}^{\,n}(X/S)
    \right)
    \simeq
    \Gamma(R_{X/S,n},\mathcal O_{R_{X/S,n}})
}
\]
as \(E_\infty\)-ring spectra, naturally in every represented simplicial operator.
\end{proposition}

\begin{proof}
Combine Proposition~\ref{prop:modal-de-rham-mapping-spectrum} with
Proposition~\ref{prop:additive-scalar-represents-functions}.
\end{proof}


\section{The infinitesimal descent monad and the Amitsur presentation}
\label{sec:monadic-de-rham-descent-monad}

Write
\[
    \epsilon:=\epsilon_{X/S}:X\longrightarrow Q:=Q_{X/S}.
\]

\begin{definition}[Infinitesimal descent monad]
\label{def:monadic-de-rham-descent-monad}
Define
\[
\boxed{
    \mathbb M_{X/S}:=\Pi_\epsilon\epsilon^*:
    \mathcal E_Q\longrightarrow\mathcal E_Q.
}
\]
Its monad unit and multiplication are those of \(\epsilon^*\dashv\Pi_\epsilon\).
\end{definition}

\begin{definition}[Monadic Amitsur object]
\label{def:monadic-de-rham-amitsur}
For \(A\in\mathcal E_Q\), define
\[
    \Am_{\mathbb M}^n(A):=\mathbb M^{n+1}A.
\]
The monad unit gives the cofaces and the monad multiplication gives the codegeneracies.  We regard
\[
    A\longrightarrow\Am_{\mathbb M}^{\bullet}(A)
\]
as the augmented monadic Amitsur object.
\end{definition}

\begin{theorem}[\v{C}ech--monad comparison]
\label{thm:monadic-de-rham-cech-monad}
Let
\[
    e_n:R_{X/S,n}\longrightarrow Q
\]
be the structural map of the \v{C}ech nerve of \(\epsilon\).  For every \(A\in\mathcal E_Q\), there
are canonical equivalences
\[
\boxed{
    \mathbb M_{X/S}^{n+1}(A)
    \xrightarrow{\simeq}
    \Pi_{e_n}e_n^*A
}
\]
which assemble into an equivalence of cosimplicial objects
\[
\boxed{
    \Am_{\mathbb M_{X/S}}^\bullet(A)
    \simeq
    \bigl(\Pi_{e_n}e_n^*A\bigr)_{n\geq0}.
}
\]
\end{theorem}

\begin{proof}
The case \(n=0\) is \(\Pi_\epsilon\epsilon^*A\).  For the first iteration, form the Cartesian
self-pullback
\[
\begin{tikzcd}
    X\times_QX \ar[r,"\operatorname{pr}_2"] \ar[d,"\operatorname{pr}_1"'] & X \ar[d,"\epsilon"]\\
    X \ar[r,"\epsilon"'] & Q.
\end{tikzcd}
\]
Right Beck--Chevalley identifies
\[
    \epsilon^*\Pi_\epsilon
    \simeq
    \Pi_{\operatorname{pr}_1}(\operatorname{pr}_2)^*
\]
after the evident source and target identifications.  Substituting this into
\(\Pi_\epsilon\epsilon^*\Pi_\epsilon\epsilon^*A\) yields
\[
    \mathbb M^2A
    \simeq
    \Pi_{e_1}e_1^*A.
\]
Repeating the same base-change calculation for the successive self-pullbacks adds one \v{C}ech
factor at each stage and gives
\[
    \mathbb M^{n+1}A\simeq\Pi_{e_n}e_n^*A.
\]
Under mate calculus, insertion of a monad unit is deletion of one factor on the geometric side and
therefore a coface after applying functions; multiplication of two adjacent monad factors is induced
by a diagonal and therefore a codegeneracy.  Hence the degreewise equivalences are cosimplicial.
\end{proof}

\begin{corollary}[Monadic descent]
\label{cor:monadic-de-rham-descent}
For every \(A\in\mathcal E_Q\),
\[
\boxed{
    A\xrightarrow{\simeq}
    \Tot\Am_{\mathbb M_{X/S}}^\bullet(A).
}
\]
\end{corollary}

\begin{proof}
The map \(\epsilon:X\to Q\) is an effective epimorphism.  By
Theorem~\ref{thm:monadic-de-rham-cech-monad}, the Amitsur object is its stable \v{C}ech descent
object.  Apply Theorem~\ref{thm:monadic-de-rham-effective-descent}.
\end{proof}

\begin{theorem}[Monadic presentation of geometric de Rham cochains]
\label{thm:cosimplicial-modal-de-rham-object}
Put
\[
    A_{X/S}:=b^*\mathbb O_S^{\add}\in\CAlg(\mathcal E_Q).
\]
There is a canonical equivalence of cosimplicial commutative algebras
\[
\boxed{
    \underline{\mathbf C}_{\dR}^{\,\bullet}(X/S)
    \simeq
    \Pi_b\Am_{\mathbb M_{X/S}}^\bullet(A_{X/S}).
}
\]
Thus the monadic Amitsur object is a presentation of the geometrically defined de Rham cochains.
\end{theorem}

\begin{proof}
Theorem~\ref{thm:monadic-de-rham-cech-monad} gives
\[
    \mathbb M^{n+1}(A_{X/S})
    \simeq
    \Pi_{e_n}e_n^*b^*\mathbb O_S^{\add}.
\]
Applying \(\Pi_b\) and using \(a_n=be_n\) gives
\[
    \Pi_b\mathbb M^{n+1}(A_{X/S})
    \simeq
    \Pi_{a_n}a_n^*\mathbb O_S^{\add}
    =
    \underline{\mathbf C}_{\dR}^{\,n}(X/S).
\]
The \v{C}ech--monad theorem identifies all cosimplicial structure maps, and all functors involved
carry their canonical lax or strong monoidal structures, so the equivalence is multiplicative.
\end{proof}


\section{Full and effective de Rham algebras}
\label{sec:total-modal-de-rham}

Write
\[
    q_{\dR}=q_{X/S}:X_{\dR/S}\longrightarrow S,
    \qquad
    b=b_{X/S}:Q_{X/S}\longrightarrow S.
\]

\begin{definition}[Full or modal de Rham algebra]
\label{def:full-modal-de-rham-object}
Define
\[
\boxed{
    \mathbf{DR}^{\mathrm{full}}_{X/S}
    :=
    \Pi_{q_{\dR}}q_{\dR}^*\mathbb O_S^{\add}
    \in
    \CAlg(\mathcal E_S).
}
\]
Its global mapping spectrum is
\[
\boxed{
    \mathbf R\Gamma_{\dR}^{\mathrm{full}}(X/S)
    :=
    \Gamma_S^{\mathrm{st}}
    \left(\mathbf{DR}^{\mathrm{full}}_{X/S}\right).
}
\]
\end{definition}

\begin{definition}[Effective or \v{C}ech de Rham algebra]
\label{def:total-modal-de-rham-object}
Define
\[
\boxed{
    \mathbf{DR}^{\mathrm{eff}}_{X/S}
    :=
    \Tot_{[n]\in\Delta}
    \underline{\mathbf C}_{\dR}^{\,n}(X/S)
    \in
    \mathcal E_S.
}
\]
Since the cosimplicial object takes values in commutative algebras, the same limit defines
canonically
\[
    \mathbf{DR}^{\mathrm{eff}}_{X/S}
    \in
    \CAlg(\mathcal E_S).
\]
Define
\[
\boxed{
    \mathbf R\Gamma_{\dR}^{\mathrm{eff}}(X/S)
    :=
    \Gamma_S^{\mathrm{st}}
    \left(\mathbf{DR}^{\mathrm{eff}}_{X/S}\right).
}
\]
For notational continuity in the sections devoted to the \v{C}ech-totalization tower, we write
\[
    \mathbf{DR}_{X/S}:=\mathbf{DR}^{\mathrm{eff}}_{X/S},
    \qquad
    \mathbf R\Gamma_{\dR}(X/S)
    :=
    \mathbf R\Gamma_{\dR}^{\mathrm{eff}}(X/S)
\]
whenever no superscript is displayed.
\end{definition}

\begin{theorem}[Effective-quotient de Rham theorem]
\label{thm:effective-quotient-de-rham}
There are canonical multiplicative equivalences
\[
\boxed{
    \mathbf{DR}^{\mathrm{eff}}_{X/S}
    \simeq
    \Pi_b b^*\mathbb O_S^{\add}
}
\]
and
\[
\boxed{
    \mathbf{DR}^{\mathrm{eff}}_{X/S}
    \simeq
    \underline{\Hom}_{\mathcal E_S}
    \left(
        \Sigma^\infty_{+,S}Q_{X/S},
        \mathbb O_S^{\add}
    \right).
}
\]
The second object carries convolution multiplication against the diagonal of \(Q_{X/S}\).
Consequently,
\[
\boxed{
    \mathbf R\Gamma_{\dR}^{\mathrm{eff}}(X/S)
    \simeq
    \map_{\mathcal E_{Q_{X/S}}}
    \left(
        \mathbf 1_{Q_{X/S}},
        b^*\mathbb O_S^{\add}
    \right).
}
\]
\end{theorem}

\begin{proof}
Let
\[
    e_n:R_{X/S,n}\longrightarrow Q_{X/S}
\]
be the \(n\)-th structural map of the \v{C}ech nerve of
\(\epsilon:X\to Q_{X/S}\).  Effective descent in \(\mathcal E_{Q_{X/S}}\) gives
\[
    b^*\mathbb O_S^{\add}
    \xrightarrow{\simeq}
    \Tot_n
    \Pi_{e_n}e_n^*b^*\mathbb O_S^{\add}.
\]
Applying \(\Pi_b\), which preserves limits, and using \(a_n=be_n\), we obtain
\[
    \Pi_b b^*\mathbb O_S^{\add}
    \simeq
    \Tot_n
    \Pi_{a_n}a_n^*\mathbb O_S^{\add}
    =
    \mathbf{DR}^{\mathrm{eff}}_{X/S}.
\]
The descent equivalence is an equivalence of commutative algebras because pullback is strong
symmetric monoidal, dependent product is lax symmetric monoidal, and limits of commutative algebra
objects are created in the underlying category.

For the internal-function formula, Theorem~\ref{thm:modal-cochains-internal-functions} gives
\[
\begin{aligned}
    \mathbf{DR}^{\mathrm{eff}}_{X/S}
    &\simeq
    \Tot_n
    \underline{\Hom}_{\mathcal E_S}
    \left(
        A_n(X/S),\mathbb O_S^{\add}
    \right)
    \\
    &\simeq
    \underline{\Hom}_{\mathcal E_S}
    \left(
        |A_\bullet(X/S)|,\mathbb O_S^{\add}
    \right).
\end{aligned}
\]
Internal Hom converts colimits in its first variable to limits.  Since fibrewise suspension spectra
preserve geometric realizations and
\[
    |R_{X/S,\bullet}|\simeq Q_{X/S},
\]
we have
\[
    |A_\bullet(X/S)|
    \simeq
    \Sigma^\infty_{+,S}Q_{X/S}.
\]
The cocommutative diagonal on the suspension spectrum of \(Q_{X/S}\) induces the convolution
commutative algebra structure, and the preceding equivalence identifies it with pointwise
multiplication on the totalized cochains.  Applying stable sections and the adjunction
\(b^*\dashv\Pi_b\) gives the global formula.
\end{proof}

\begin{proposition}[Internal-function presentation of the full de Rham algebra]
\label{prop:full-dr-internal-functions}
There is a canonical multiplicative equivalence
\[
\boxed{
    \mathbf{DR}^{\mathrm{full}}_{X/S}
    \simeq
    \underline{\Hom}_{\mathcal E_S}
    \left(
        \Sigma^\infty_{+,S}X_{\dR/S},
        \mathbb O_S^{\add}
    \right).
}
\]
Consequently,
\[
\boxed{
    \mathbf R\Gamma_{\dR}^{\mathrm{full}}(X/S)
    \simeq
    \map_{\mathcal E_{X_{\dR/S}}}
    \left(
        \mathbf1_{X_{\dR/S}},
        q_{\dR}^*\mathbb O_S^{\add}
    \right).
}
\]
\end{proposition}

\begin{proof}
Apply Proposition~\ref{prop:parameterized-function-spectrum-identity} to
\(q_{\dR}:X_{\dR/S}\to S\) and \(E=\mathbb O_S^{\add}\).  Multiplicativity is the convolution
multiplication induced by the diagonal of \(X_{\dR/S}\).  The global formula is adjunction.
\end{proof}

\begin{theorem}[Full-to-effective comparison]
\label{thm:full-effective-dr-comparison}
The monomorphism
\[
    \jmath_{X/S}:Q_{X/S}\longrightarrow X_{\dR/S}
\]
induces a canonical morphism of commutative algebras
\[
\boxed{
    \rho_{X/S}:
    \mathbf{DR}^{\mathrm{full}}_{X/S}
    \longrightarrow
    \mathbf{DR}^{\mathrm{eff}}_{X/S}.
}
\]
It is functorial in \(X\to S\), compatible with arbitrary base change, and is an equivalence whenever
\(\eta_{X/S}:X\to X_{\dR/S}\) is an effective epimorphism.  In particular, if \(X\to S\) is
represented by a smooth spectral morphism, then
\[
\boxed{
    \mathbf{DR}^{\mathrm{full}}_{X/S}
    \xrightarrow{\simeq}
    \mathbf{DR}^{\mathrm{eff}}_{X/S}.
}
\]
\end{theorem}

\begin{proof}
Under Propositions~\ref{prop:full-dr-internal-functions} and
Theorem~\ref{thm:effective-quotient-de-rham}, precomposition with
\[
    \Sigma^\infty_{+,S}\jmath_{X/S}:
    \Sigma^\infty_{+,S}Q_{X/S}
    \longrightarrow
    \Sigma^\infty_{+,S}X_{\dR/S}
\]
gives
\[
    \underline{\Hom}
    \left(
        \Sigma^\infty_{+,S}X_{\dR/S},\mathbb O_S^{\add}
    \right)
    \longrightarrow
    \underline{\Hom}
    \left(
        \Sigma^\infty_{+,S}Q_{X/S},\mathbb O_S^{\add}
    \right).
\]
The diagonal identity
\[
    \Delta_{X_{\dR/S}}\jmath_{X/S}
    =
    (\jmath_{X/S}\times\jmath_{X/S})\Delta_{Q_{X/S}}
\]
shows that precomposition preserves convolution multiplication, so the map is one of commutative
algebras.

Equivalently, it is obtained by applying \(\Pi_{q_{\dR}}\) to the unit
\[
    q_{\dR}^*\mathbb O_S^{\add}
    \longrightarrow
    \Pi_{\jmath_{X/S}}\jmath_{X/S}^*q_{\dR}^*\mathbb O_S^{\add}
\]
and using \(q_{\dR}\jmath_{X/S}=b\).  This description proves functoriality and, together with
Beck--Chevalley and base change of the image factorization, arbitrary base-change compatibility.

If \(\eta_{X/S}\) is an effective epimorphism, its image is all of \(X_{\dR/S}\), so
\(\jmath_{X/S}\) is an equivalence.  Smooth effectivity from the preceding chapter gives the final
assertion.
\end{proof}

\begin{proposition}[Augmentation to scalar sections on \(X\)]
\label{prop:modal-de-rham-augmentation}
Let \(p:X\to S\) be the structural morphism.  Evaluation of the effective totalization cone in
degree zero defines a canonical morphism of commutative algebras
\[
\boxed{
    \varepsilon_{\dR}:
    \mathbf{DR}^{\mathrm{eff}}_{X/S}
    \longrightarrow
    \Pi_p p^*\mathbb O_S^{\add}.
}
\]
It is compatible with arbitrary parameterized base change.
\end{proposition}

\begin{proof}
The degree-zero infinitesimal simplex is \(R_{X/S,0}=X\), so
\[
    \underline{\mathbf C}_{\dR}^{\,0}(X/S)
    =
    \Pi_p p^*\mathbb O_S^{\add}.
\]
The totalization cone gives the displayed evaluation map.  Equivalently, under
Theorem~\ref{thm:effective-quotient-de-rham}, it is \(\Pi_b\) applied to the adjunction unit
\[
    b^*\mathbb O_S^{\add}
    \longrightarrow
    \Pi_\epsilon\epsilon^*b^*\mathbb O_S^{\add}.
\]
Both descriptions are multiplicative and commute with base change by Beck--Chevalley.
\end{proof}


\section{The \v{C}ech-totalization tower and coherent normalization}
\label{sec:stable-normalization-brave-new-forms}
\label{sec:modal-hodge-filtration}

The coherent normalized-cochain structure is extracted from the totalization tower of the actual
cosimplicial infinitesimal-function object.  This avoids introducing a second independent category
of ``complexes'' as foundational data.

\begin{convention}[Coherent cochains and Dold--Kan signs]
\label{conv:monadic-de-rham-coherent-cochains}
Let \(\Xi\) be the pointed category with objects \(\mathbb Z\cup\{*\}\), zero morphisms factoring
through \(*\), and generating nonzero morphisms
\[
    \partial_n:n\longrightarrow n-1
\]
subject to \(\partial_{n-1}\partial_n=0\).  Let \(\Xi_{\geq0}\) be the full pointed subcategory on
\(*\) and the nonnegative integers.  For a stable \(\infty\)-category \(\mathcal D\), write
\[
    \operatorname{CCh}_{\geq0}(\mathcal D)
    :=
    \Fun_*\bigl(\Xi_{\geq0}^{\op},\mathcal D\bigr)
\]
for the \(\infty\)-category of coherent nonnegative cochain objects.  Thus an object consists of
terms \(C^q\), maps \(C^q\to C^{q+1}\), specified nullhomotopies of consecutive composites, and
the higher coherence data forced by the pointed source category.  Here \(\Xi_{\geq0}\) is regarded
through its nerve.  In particular, the relation that two consecutive generators compose to the zero
morphism supplies the corresponding specified nullhomotopy of the consecutive cochain composite, and
the higher simplices of the nerve supply the compatibility data among these nullhomotopies.

We normalize stable Dold--Kan so that, for a cosimplicial object, the image of
this differential in the homotopy category is the alternating coface sum
\[
    \sum_i(-1)^id^i.
\]
\end{convention}

\begin{lemma}[Completion and associated-graded detection]
\label{lem:complete-filtration-reflection}
Let \(\mathcal D\) be a stable \(\infty\)-category admitting sequential limits, and let
\[
    \operatorname{Fil}(\mathcal D)
    :=
    \Fun(\mathbb Z^{\op},\mathcal D)
\]
be the category of decreasing \(\mathbb Z\)-indexed filtrations.  For
\(F\in\operatorname{Fil}(\mathcal D)\), put
\[
    F^{+\infty}
    :=
    \limop_{q\to+\infty}F^q
\]
and define
\[
\boxed{
    \widehat F^{\,q}
    :=
    \cofib
    \left(
        F^{+\infty}
        \longrightarrow
        F^q
    \right).
}
\]
Then:
\begin{enumerate}
\item \(\widehat F\) is complete:
\[
    \limop_q\widehat F^{\,q}\simeq0;
\]
\item the canonical morphism \(F\to\widehat F\) exhibits completion as an exact
reflective localization of \(\operatorname{Fil}(\mathcal D)\) onto the full
subcategory of complete filtrations;
\item completion does not change associated graded:
\[
\boxed{
    \gr^q\widehat F\xrightarrow{\simeq}\gr^qF;
}
\]
\item a morphism between complete filtrations is an equivalence if and only if
it is an equivalence on every associated-graded term.
\end{enumerate}
\end{lemma}

\begin{proof}
There is a canonical cofiber sequence of filtrations
\[
    \operatorname{const}(F^{+\infty})
    \longrightarrow
    F
    \longrightarrow
    \widehat F.
\]
Sequential limits are exact in a stable \(\infty\)-category, so
\[
\begin{aligned}
    \limop_q\widehat F^{\,q}
    &\simeq
    \cofib
    \left(
        F^{+\infty}
        \longrightarrow
        \limop_qF^q
    \right)
    \\
    &\simeq
    \cofib(\id_{F^{+\infty}})
    \simeq0.
\end{aligned}
\]
Thus \(\widehat F\) is complete.

Let \(G\) be complete.  In the stable functor category,
\[
\begin{aligned}
    \map_{\operatorname{Fil}(\mathcal D)}
    \left(
        \operatorname{const}(F^{+\infty}),G
    \right)
    &\simeq
    \map_{\mathcal D}
    \left(
        F^{+\infty},
        \limop_qG^q
    \right)
    \\
    &\simeq0.
\end{aligned}
\]
Applying \(\map(-,G)\) to the displayed cofiber sequence therefore gives
\[
    \map(\widehat F,G)
    \xrightarrow{\simeq}
    \map(F,G).
\]
This is the universal property of the reflector onto complete filtrations.
Exactness follows from the defining cofiber formula and exactness of sequential
limits.

Compare the two consecutive cofiber squares
\[
\begin{tikzcd}
    F^{+\infty} \ar[r] \ar[d,equal]
        & F^{q+1} \ar[d]\\
    F^{+\infty} \ar[r]
        & F^q.
\end{tikzcd}
\]
The stable \(3\times3\) lemma identifies the cofiber of
\(\widehat F^{\,q+1}\to\widehat F^{\,q}\) with the cofiber of
\(F^{q+1}\to F^q\), proving
\[
    \gr^q\widehat F\simeq\gr^qF.
\]

Finally, let \(u:F\to G\) be a morphism of complete filtrations whose associated
graded is an equivalence in every degree, and let \(H:=\fib(u)\).  Then \(H\)
is complete and \(\gr^qH\simeq0\) for every \(q\).  Hence every transition
\(H^{q+1}\to H^q\) is an equivalence.  The filtration \(H\) is therefore
constant, while completeness gives
\[
    H^q\simeq\limop_rH^r\simeq0.
\]
Thus \(H\simeq0\) and \(u\) is an equivalence.  The converse is immediate.
\end{proof}

\begin{theorem}[Complete filtrations and coherent cochains]
\label{thm:complete-filtered-coherent-cochains}
Let \(\mathcal D\) be a stable \(\infty\)-category admitting sequential limits.  Let
\(\operatorname{Fil}^{\mathrm{comp}}_{\geq0}(\mathcal D)\) be the full subcategory of decreasing
\(\mathbb Z\)-indexed filtrations
\[
    \cdots\to F^{q+1}\to F^q\to\cdots
\]
for which \(F^q\to F^0\) is an equivalence for every \(q\leq0\), and which are complete in the
sense that
\[
    \limop_qF^q\simeq0.
\]
There is a canonical equivalence
\[
\boxed{
    \Phi_{\mathcal D}:
    \operatorname{Fil}^{\mathrm{comp}}_{\geq0}(\mathcal D)
    \xrightarrow{\simeq}
    \operatorname{CCh}_{\geq0}(\mathcal D).
}
\]
For a complete filtration \(F^\bullet\), the degree-\(q\) term of
\(\Phi_{\mathcal D}(F^\bullet)\) is
\[
    \gr^qF[q],
\]
and its first differential is induced by the connecting morphism of the consecutive cofiber
sequences
\[
    F^{q+1}\longrightarrow F^q\longrightarrow\gr^qF.
\]
The functor is natural in exact functors preserving the sequential limits involved.
\end{theorem}

\begin{proof}
Let
\[
    \operatorname{CCh}(\mathcal D)
    :=
    \Fun_*(\Xi^{\op},\mathcal D)
\]
denote the \(\mathbb Z\)-indexed coherent cochain category, and let
\(\operatorname{Fil}^{\mathrm{comp}}(\mathcal D)\) denote the category of
complete decreasing \(\mathbb Z\)-indexed filtrations.  The pointed indexing
category used in Ariotta's construction is exactly the category \(\Xi\) fixed in
Convention~\ref{conv:monadic-de-rham-coherent-cochains}.  The general
complete-filtered/coherent-cochain equivalence of Ariotta\footnote{Stefano
Ariotta, \emph{Coherent cochain complexes and Beilinson
\(t\)-structures}, arXiv:2109.01017, Definition~2.1, Theorem~4.7, and
Remark~4.9.} gives a canonical exact equivalence
\[
\boxed{
    \operatorname{Fil}^{\mathrm{comp}}(\mathcal D)
    \xrightarrow{\simeq}
    \operatorname{CCh}(\mathcal D)
}
\]
whose value on a complete filtration \(F\) has degree-\(q\) term
\[
\boxed{
    \gr^qF[q].
}
\]
The first cochain structure map is the shifted connecting morphism obtained
from the two consecutive cofiber sequences
\[
    F^{q+1}\longrightarrow F^q\longrightarrow\gr^qF,
    \qquad
    F^{q+2}\longrightarrow F^{q+1}\longrightarrow\gr^{q+1}F.
\]
The pointed indexing category \(\Xi\) retains the specified nullhomotopies of
consecutive composites and all higher coherences.

It remains only to identify the nonnegative restrictions used in the statement.
For a complete decreasing filtration \(F\), the condition
\[
    F^q\xrightarrow{\simeq}F^0
    \qquad(q\leq0)
\]
is equivalent to
\[
    \gr^qF\simeq0
    \qquad(q<0).
\]
Indeed, the displayed equivalences force the negative associated-graded terms
to vanish, while conversely vanishing of all negative associated-graded terms
makes every transition
\[
    F^{q+1}\longrightarrow F^q
    \qquad(q<0)
\]
an equivalence, and finite composition identifies each \(F^q\), \(q\leq0\),
with \(F^0\).

Under the general equivalence, the condition
\(\gr^qF\simeq0\) for \(q<0\) is exactly the condition that the associated
coherent cochain object vanish in negative degrees.  Restriction along
\[
    \Xi_{\geq0}^{\op}\hookrightarrow\Xi^{\op}
\]
identifies the full subcategory of coherent cochain objects vanishing in
negative degrees with
\[
    \Fun_*(\Xi_{\geq0}^{\op},\mathcal D)
    =
    \operatorname{CCh}_{\geq0}(\mathcal D).
\]
Its inverse is extension by the zero object in every negative degree.  All
new structure maps involving a negative term are then the unique zero maps,
and the pointed relations and higher simplices of \(\Xi\) supply the unique
compatible coherent zero data.  Restricting the general equivalence therefore
gives
\[
\boxed{
    \Phi_{\mathcal D}:
    \operatorname{Fil}^{\mathrm{comp}}_{\geq0}(\mathcal D)
    \xrightarrow{\simeq}
    \operatorname{CCh}_{\geq0}(\mathcal D).
}
\]
The degree-\(q\) term and first differential are the restrictions of the
pointwise description above, hence are exactly
\[
    \gr^qF[q]
\]
and the connecting morphism asserted in the theorem.

The construction of the complete-filtered/coherent-cochain equivalence is
functorial for exact functors preserving the sequential limits used to define
completeness.  Restriction to the nonnegative full subcategories is likewise
functorial, which proves the final assertion.
\end{proof}

\begin{definition}[Partial totalizations and \v{C}ech-totalization filtration]
\label{def:modal-hodge-filtration}
For \(q\geq0\), let \(\Delta_{\leq q-1}\subseteq\Delta\) be the full subcategory on
\([0],\ldots,[q-1]\), with \(\Delta_{\leq-1}=\varnothing\).  Define
\[
\boxed{
    \mathbf{DR}_{X/S}^{<q}
    :=
    \Tot_{\Delta_{\leq q-1}}
    \underline{\mathbf C}_{\dR}^{\,\bullet}(X/S),
}
\]
where \(\mathbf{DR}_{X/S}^{<0}\simeq0\).  Define the \(q\)-th totalization stage by
\[
\boxed{
    F_{\mathrm{Tot}}^q\mathbf{DR}_{X/S}
    :=
    \fib
    \left(
        \mathbf{DR}_{X/S}
        \longrightarrow
        \mathbf{DR}_{X/S}^{<q}
    \right).
}
\]
Thus \(F_{\mathrm{Tot}}^0\mathbf{DR}_{X/S}\simeq\mathbf{DR}_{X/S}\), and the restriction maps between partial
totalizations induce a decreasing tower
\[
    \cdots
    \longrightarrow
    F_{\mathrm{Tot}}^{q+1}\mathbf{DR}_{X/S}
    \longrightarrow
    F_{\mathrm{Tot}}^q\mathbf{DR}_{X/S}
    \longrightarrow
    \cdots
    \longrightarrow
    \mathbf{DR}_{X/S}.
\]
\end{definition}

\begin{definition}[Normalized infinitesimal de Rham cochains]
\label{def:normalized-modal-de-rham-object}
For \(q\geq0\), let
\[
    M_\sigma^q
    \left(
        \underline{\mathbf C}_{\dR}^{\,\bullet}(X/S)
    \right)
    :=
    \limop_{\substack{\alpha:[q]\twoheadrightarrow[k]\\ k<q}}
    \underline{\mathbf C}_{\dR}^{\,k}(X/S)
\]
be the codegeneracy matching object.  Define the parameterized normalized infinitesimal \(q\)-cochain spectrum by
\[
\boxed{
    \underline{\OmegaInf}^{\,q}(X/S)
    :=
    \fib
    \left(
        \underline{\mathbf C}_{\dR}^{\,q}(X/S)
        \longrightarrow
        M_\sigma^q
        \left(
            \underline{\mathbf C}_{\dR}^{\,\bullet}(X/S)
        \right)
    \right).
}
\]
For \(n\in\mathbb Z\), define the global spectrum of shifted normalized infinitesimal \(q\)-cochains by
\[
\boxed{
    \OmegaInf^q(X/S;n)
    :=
    \map_{\mathcal E_S}
    \left(
        \mathbf 1_S,
        \underline{\OmegaInf}^{\,q}(X/S)[n]
    \right).
}
\]
Write \(\OmegaInf^q(X/S)\) for \(n=0\).
\end{definition}

\begin{proposition}[Explicit normalization]
\label{prop:explicit-modal-normalized-terms}
The object \(\underline{\OmegaInf}^{\,q}(X/S)\) is equivalently the total fibre of the finite cube
generated by the elementary codegeneracies
\[
    s^i:
    \underline{\mathbf C}_{\dR}^{\,q}(X/S)
    \longrightarrow
    \underline{\mathbf C}_{\dR}^{\,q-1}(X/S).
\]
In particular,
\[
    \underline{\OmegaInf}^{\,0}(X/S)
    \simeq
    \underline{\mathbf C}_{\dR}^{\,0}(X/S),
\]
and there is a fibre sequence
\[
\boxed{
    \underline{\OmegaInf}^{\,1}(X/S)
    \longrightarrow
    \underline{\mathbf C}_{\dR}^{\,1}(X/S)
    \xrightarrow{s^0}
    \underline{\mathbf C}_{\dR}^{\,0}(X/S).
}
\]
\end{proposition}

\begin{proof}
Every nonidentity surjection in \(\Delta\) is a composite of the elementary surjections, and their
relations are precisely the codegeneracy identities.  In a stable category the intersection of the
corresponding kernels is the homotopy fibre of the map to the matching limit, equivalently the total
fibre of the elementary codegeneracy cube.
\end{proof}

\begin{proposition}[Stable Dold--Kan and the complete totalization filtration]
\label{prop:monadic-de-rham-stable-dold-kan-filtration}
Let \(\mathcal D\) be a stable \(\infty\)-category with the sequential limits needed below, and let
\[
    C^\bullet\in\Fun(\Delta,\mathcal D).
\]
For \(q\geq0\), put
\[
    N^qC
    :=
    \fib
    \left(
        C^q\longrightarrow
        \limop_{\substack{\alpha:[q]\twoheadrightarrow[k]\\k<q}}C^k
    \right),
\]
let
\[
    D_q(C):=\Tot_{\Delta_{\leq q}}C^\bullet,
    \qquad
    D_{-1}(C):=0,
\]
and write
\[
    D_\infty(C):=\Tot_\Delta C^\bullet.
\]
Then:
\begin{enumerate}
\item there are canonical fibre equivalences
\[
\boxed{
    \fib\bigl(D_q(C)\to D_{q-1}(C)\bigr)
    \simeq
    N^qC[-q];
}
\]
\item the decreasing filtration
\[
    F^qC
    :=
    \fib\bigl(D_\infty(C)\to D_{q-1}(C)\bigr)
\]
is complete and has associated graded
\[
\boxed{
    \gr^qF^\bullet C\simeq N^qC[-q];
}
\]
\item after extending the filtration constantly to nonpositive degrees, the equivalence of Theorem~\ref{thm:complete-filtered-coherent-cochains} sends
\(F^\bullet C\) to a coherent nonnegative cochain object whose degree-\(q\) term is \(N^qC\).
Its first structure maps are the connecting morphisms
\[
    N^qC\longrightarrow N^{q+1}C,
\]
their composites carry specified nullhomotopies, and the coherent cochain object retains all higher
compatibility data.  In the homotopy category its differential is represented by
\[
    \sum_{i=0}^{q+1}(-1)^id^i.
\]
\end{enumerate}
The construction is natural in \(C^\bullet\) and in exact functors which preserve the limits used in
the construction.
\end{proposition}

\begin{proof}
For the layer calculation use the standard stable Dold--Kan equivalence, independently of
Theorem~\ref{thm:complete-filtered-coherent-cochains}.  Apply it in the opposite stable
\(\infty\)-category \(\mathcal D^{\op}\) to the simplicial object corresponding to
\(C^\bullet\).  The skeletal filtration of that simplicial object has successive cofibers given by
its normalized terms shifted by the simplicial degree.  Under passage back to \(\mathcal D\), the
normalized cofiber in \(\mathcal D^{\op}\) is the codegeneracy matching fibre \(N^qC\), the
skeletal colimit tower becomes the partial-totalization limit tower, and cofibers become fibres.
Thus
\[
    \fib\bigl(D_q(C)\to D_{q-1}(C)\bigr)
    \simeq
    N^qC[-q].
\]

The simplex category is the filtered union of the full subcategories on
\([0],\ldots,[q]\), hence
\[
    D_\infty(C)
    \simeq
    \limop_qD_q(C).
\]
Therefore
\[
    \limop_qF^qC
    \simeq
    \fib
    \left(
        D_\infty(C)\longrightarrow\limop_qD_{q-1}(C)
    \right)
    \simeq0,
\]
so the filtration is complete.  Comparing the fibre sequences defining \(F^{q+1}C\) and
\(F^qC\), the stable \(3\times3\) lemma gives
\[
    \cofib(F^{q+1}C\to F^qC)
    \simeq
    \fib(D_q(C)\to D_{q-1}(C))
    \simeq
    N^qC[-q].
\]

Extend \(F^\bullet C\) to a decreasing \(\mathbb Z\)-filtration by setting
\(F^qC=F^0C\) for \(q\leq0\).  Its negative associated-graded terms vanish.  Under Theorem~\ref{thm:complete-filtered-coherent-cochains}, the term in degree
\(q\geq0\) is
\[
    \gr^qF[q]
    \simeq
    N^qC[-q][q]
    \simeq
    N^qC.
\]
The first structure map is obtained by composing the connecting map
\[
    \gr^qF\longrightarrow F^{q+1}[1]
\]
with the shifted quotient to \(\gr^{q+1}F[1]\), and then undoing the grading shifts.  Because this
construction is performed inside the coherent-cochain/complete-filtration equivalence, it includes
the specified nullhomotopies of two consecutive differentials and all higher compatibility data.
With the sign convention fixed in Convention~\ref{conv:monadic-de-rham-coherent-cochains}, stable
Dold--Kan identifies its image in the homotopy category with the ordinary normalized cochain
differential
\[
    \sum_{i=0}^{q+1}(-1)^id^i.
\]
Naturality follows from functoriality of partial totalization, normalization, exact triangles, and
Theorem~\ref{thm:complete-filtered-coherent-cochains}.
\end{proof}

\begin{proposition}[Canonical degree-zero normalized differential]
\label{prop:degree-zero-normalized-differential}
Let \(\mathcal D\) be a stable \(\infty\)-category and let
\[
    C^\bullet\in\Fun(\Delta,\mathcal D).
\]
Under the canonical identifications
\[
    N^0C\simeq C^0,
    \qquad
    N^1C\simeq\fib\left(s^0:C^1\longrightarrow C^0\right),
\]
the cosimplicial identities
\[
    s^0d^0\simeq\id_{C^0},
    \qquad
    s^0d^1\simeq\id_{C^0}
\]
determine a canonical nullhomotopy
\[
    s^0(d^0-d^1)\simeq0
\]
and hence a canonical factorization
\[
\boxed{
    \delta_C:
    C^0
    \longrightarrow
    N^1C
    \longrightarrow
    C^1
}
\]
whose composite with \(N^1C\to C^1\) is \(d^0-d^1\).  The morphism
\[
    \delta_C:N^0C\longrightarrow N^1C
\]
is the degree-zero structure map of the coherent normalized cochain object associated with the
complete totalization filtration in
Proposition~\ref{prop:monadic-de-rham-stable-dold-kan-filtration}.

In particular, for
\[
    C^\bullet
    =
    \underline{\mathbf C}_{\dR}^{\,\bullet}(X/S),
\]
the coherent infinitesimal differential
\[
    d_{\inf}^0:
    \underline{\OmegaInf}^{\,0}(X/S)
    \longrightarrow
    \underline{\OmegaInf}^{\,1}(X/S)
\]
is canonically this factorization of \(d^0-d^1\).
\end{proposition}

\begin{proof}
The two cosimplicial identities displayed above identify the composites of \(d^0\) and \(d^1\)
with \(s^0\) through the identity of \(C^0\).  Since mapping spaces in a stable
\(\infty\)-category are grouplike, their difference therefore carries the canonical nullhomotopy
\[
    s^0(d^0-d^1)\simeq0.
\]
The universal property of the fibre of \(s^0\) gives the canonical factorization
\(\delta_C:C^0\to N^1C\).

Let
\[
    E:=\fib(\delta_C).
\]
For every \(K\in\mathcal D\), a map \(K\to E\) is a map \(x:K\to C^0\) together with a
nullhomotopy of \(\delta_Cx\).  After the inclusion \(N^1C\to C^1\), this is equivalently a
specified homotopy
\[
    d^0x\simeq d^1x
\]
whose restriction along \(s^0\) is the canonical identity homotopy.  Together with the
cosimplicial identities, this is exactly the data of a cone from \(K\) to the restriction of
\(C^\bullet\) to \(\Delta_{\leq1}\).  Hence
\[
    E\simeq D_1(C)=\Tot_{\Delta_{\leq1}}C^\bullet
\]
compatibly with the projection to
\[
    D_0(C)=C^0.
\]
Thus the fibre sequence
\[
    N^1C[-1]\longrightarrow D_1(C)\longrightarrow C^0
\]
has connecting morphism \(\delta_C:C^0\to N^1C\).  By the construction of the coherent cochain
object from the complete totalization filtration, this connecting morphism is its degree-zero
structure map.  The final assertion follows by specializing to the geometric de Rham cochains.
\end{proof}

\begin{theorem}[Stable Dold--Kan layer formula and coherent infinitesimal differential]
\label{thm:monadic-de-rham-dold-kan}
For every \(q\geq0\), the partial totalization tower has canonical fibre equivalences
\[
\boxed{
    \fib
    \left(
        \mathbf{DR}_{X/S}^{<q+1}
        \longrightarrow
        \mathbf{DR}_{X/S}^{<q}
    \right)
    \simeq
    \underline{\OmegaInf}^{\,q}(X/S)[-q].
}
\]
Consequently there are canonical cofiber sequences
\[
\boxed{
    F_{\mathrm{Tot}}^{q+1}\mathbf{DR}_{X/S}
    \longrightarrow
    F_{\mathrm{Tot}}^q\mathbf{DR}_{X/S}
    \longrightarrow
    \underline{\OmegaInf}^{\,q}(X/S)[-q].
}
\]
The complete filtered object \(F_{\mathrm{Tot}}^\bullet\mathbf{DR}_{X/S}\) canonically determines boundary maps
\[
\boxed{
    d_{\inf}^q:
    \underline{\OmegaInf}^{\,q}(X/S)
    \longrightarrow
    \underline{\OmegaInf}^{\,q+1}(X/S).
}
\]
Their images in the homotopy category are represented by the alternating coface formula
\[
\boxed{
    d_{\inf}^q
    =
    \sum_{i=0}^{q+1}(-1)^id^i.
}
\]
There is a specified coherent nullhomotopy
\[
    d_{\inf}^{q+1}d_{\inf}^{q}\simeq0,
\]
and the complete totalization filtration retains the full hierarchy of higher compatibility data
supplied by the coherent normalized cochain object associated with the original cosimplicial object.
\end{theorem}

\begin{proof}
Apply Proposition~\ref{prop:monadic-de-rham-stable-dold-kan-filtration} to
\[
    C^\bullet
    =
    \underline{\mathbf C}_{\dR}^{\,\bullet}(X/S).
\]
Its normalized term \(N^qC\) is exactly \(\underline{\OmegaInf}^{\,q}(X/S)\) by
Definition~\ref{def:normalized-modal-de-rham-object}, and
\[
    D_{q-1}(C)=\mathbf{DR}_{X/S}^{<q},
    \qquad
    D_\infty(C)=\mathbf{DR}_{X/S}.
\]
The first two displayed assertions therefore follow directly from the layer and associated-graded
statements of that proposition.

For clarity, write
\[
    G_{\mathrm{Tot}}^q
    :=
    \gr_{\mathrm{Tot}}^q\mathbf{DR}_{X/S}
    \simeq
    \underline{\OmegaInf}^{\,q}(X/S)[-q].
\]
The cofiber sequence
\[
    F_{\mathrm{Tot}}^{q+1}\mathbf{DR}
    \longrightarrow
    F_{\mathrm{Tot}}^q\mathbf{DR}
    \longrightarrow
    G_{\mathrm{Tot}}^q
\]
has its canonical connecting morphism
\[
    \partial_q:G_{\mathrm{Tot}}^q\longrightarrow F_{\mathrm{Tot}}^{q+1}\mathbf{DR}[1].
\]
Composing with the shifted quotient map gives
\[
    G_{\mathrm{Tot}}^q
    \xrightarrow{\partial_q}
    F_{\mathrm{Tot}}^{q+1}\mathbf{DR}[1]
    \longrightarrow
    G_{\mathrm{Tot}}^{q+1}[1].
\]
Since
\[
    G_{\mathrm{Tot}}^q\simeq\underline{\OmegaInf}^{\,q}(X/S)[-q],
    \qquad
    G_{\mathrm{Tot}}^{q+1}[1]\simeq\underline{\OmegaInf}^{\,q+1}(X/S)[-q],
\]
shifting by \(q\) gives \(d_{\inf}^q\).  Proposition~\ref{prop:monadic-de-rham-stable-dold-kan-filtration}
identifies these connecting morphisms with the first maps of the coherent normalized cochain object;
it therefore supplies the specified nullhomotopy of \(d_{\inf}^{q+1}d_{\inf}^q\), all higher coherent
relations, and the alternating-coface formula in the homotopy category.
\end{proof}

\begin{proposition}[The coherent infinitesimal boundary]
\label{prop:coherent-modal-de-rham-boundary}
The maps \(d_{\inf}^q\) of Theorem~\ref{thm:monadic-de-rham-dold-kan} are natural in \(X\to S\).
They are not additional structure: they are the connecting morphisms of the \v{C}ech-totalization
tower, assembled into a coherent cochain object by the complete-filtration correspondence of
Proposition~\ref{prop:monadic-de-rham-stable-dold-kan-filtration}.
\end{proposition}

\begin{proof}
Every map of cosimplicial objects induces a map of partial totalization towers and therefore a map of
the associated complete filtered objects.  Connecting morphisms in stable categories are natural.
\end{proof}

\begin{theorem}[Associated graded and totalization completeness]
\label{thm:modal-hodge-completeness}
For every \(q\geq0\),
\[
\boxed{
    \gr_{\mathrm{Tot}}^q\mathbf{DR}_{X/S}
    :=
    \cofib
    \left(
        F_{\mathrm{Tot}}^{q+1}\mathbf{DR}_{X/S}
        \longrightarrow
        F_{\mathrm{Tot}}^q\mathbf{DR}_{X/S}
    \right)
    \simeq
    \underline{\OmegaInf}^{\,q}(X/S)[-q].
}
\]
The \v{C}ech-totalization filtration is complete and separated:
\[
\boxed{
    \limop_qF_{\mathrm{Tot}}^q\mathbf{DR}_{X/S}\simeq0
}
\]
and
\[
\boxed{
    \mathbf{DR}_{X/S}
    \xrightarrow{\simeq}
    \limop_q\mathbf{DR}_{X/S}^{<q}.
}
\]
Each \(\mathbf{DR}_{X/S}^{<q}\) is canonically a commutative algebra, the transition maps are maps
of commutative algebras, and the completion equivalence is an equivalence in
\(\CAlg(\mathcal E_S)\).
\end{theorem}

\begin{proof}
The associated-graded formula is the second cofiber sequence of
Theorem~\ref{thm:monadic-de-rham-dold-kan}.  By definition,
\[
    \mathbf{DR}_{X/S}
    =
    \limop_{\Delta}
    \underline{\mathbf C}_{\dR}^{\,\bullet}(X/S).
\]
The simplex category is the union of its finite full subcategories, so Fubini for limits gives
\[
    \mathbf{DR}_{X/S}
    \simeq
    \limop_q
    \mathbf{DR}_{X/S}^{<q}.
\]
Therefore
\[
\begin{aligned}
    \limop_qF_{\mathrm{Tot}}^q\mathbf{DR}_{X/S}
    &\simeq
    \fib
    \left(
        \mathbf{DR}_{X/S}
        \longrightarrow
        \limop_q\mathbf{DR}_{X/S}^{<q}
    \right)
    \\
    &\simeq0.
\end{aligned}
\]
Since \(\underline{\mathbf C}_{\dR}^{\,\bullet}(X/S)\) is a cosimplicial commutative algebra, every
partial totalization is a limit in \(\CAlg(\mathcal E_S)\).  Such limits are created in the
underlying stable category, so the same Fubini argument proves the completion formula in
commutative algebras.
\end{proof}

\begin{proposition}[Parameterized base change of normalized cochains and the totalization tower]
\label{prop:normalized-forms-base-change}
For \(c:S'\to S\) and \(X'=X\times_SS'\), pullback induces canonical equivalences
\[
\boxed{
    c^*\underline{\OmegaInf}^{\,q}(X/S)
    \simeq
    \underline{\OmegaInf}^{\,q}(X'/S'),
}
\]
\[
\boxed{
    c^*\mathbf{DR}_{X/S}^{<q}
    \simeq
    \mathbf{DR}_{X'/S'}^{<q},
    \qquad
    c^*F_{\mathrm{Tot}}^q\mathbf{DR}_{X/S}
    \simeq
    F_{\mathrm{Tot}}^q\mathbf{DR}_{X'/S'}.
}
\]
They intertwine the complete coherent infinitesimal differential and satisfy all pullback-pasting
coherences.  Stable sections give restriction morphisms on global normalized-cochain spectra; these restriction
maps are not asserted to be equivalences for arbitrary \(c\).
\end{proposition}

\begin{proof}
Theorem~\ref{thm:modal-nerve-base-change},
Proposition~\ref{prop:additive-scalar-spectrum-base-change}, and right Beck--Chevalley give
\[
    c^*\underline{\mathbf C}_{\dR}^{\,\bullet}(X/S)
    \simeq
    \underline{\mathbf C}_{\dR}^{\,\bullet}(X'/S').
\]
The pullback \(c^*\) on parameterized spectra preserves all limits and colimits.  It therefore
commutes with the finite matching limits defining \(\underline{\OmegaInf}^{\,q}\), all partial and complete
totalizations, and the fibres defining \(F_{\mathrm{Tot}}^q\).  Naturality of the Dold--Kan filtration gives the
boundary compatibility.
\end{proof}


\section{First-order linearization and the cotangent complex}
\label{sec:first-order-linearization-cotangent}

Assume throughout this section that \(p:X\to S\) is represented by a morphism of spectral schemes.
The constructions above retain all finite infinitesimal orders.  We now extract their canonical
first-order quotient and compare it with the cotangent complex constructed in
Chapter~\ref{chap:cotangent-spectral-deformation}.

\begin{construction}[First structured infinitesimal neighbourhood of the diagonal]
\label{constr:first-diagonal-thickening}
Put
\[
    Y:=X\times_SX
\]
and let
\[
    \Delta:X\longrightarrow Y
\]
be the diagonal.  Regard \(Y\) over \(X\) through the first projection.  The composite
\[
    X\xrightarrow{\Delta}Y\xrightarrow{\operatorname{pr}_1}X
\]
is the identity.  Cotangent base change gives
\[
    L_{Y/X}
    \simeq
    \operatorname{pr}_2^*L_{X/S},
\]
and hence
\[
    \Delta^*L_{Y/X}\simeq L_{X/S}.
\]
The transitivity triangle for the displayed composite is
\[
    \Delta^*L_{Y/X}
    \longrightarrow
    L_{X/X}\simeq0
    \longrightarrow
    L_{X/Y}.
\]
We fix the resulting rotation by defining
\[
\boxed{
    \partial_\Delta:
    L_{X/Y}
    \xrightarrow{\simeq}
    \Delta^*L_{Y/X}[1]
    \xrightarrow{\simeq}
    L_{X/S}[1]
}
\]
to be the connecting equivalence of this transitivity triangle followed by
cotangent base change.  This fixes the sign of the first-order diagonal class.

The global structured square-zero construction of
Chapter~\ref{chap:cotangent-spectral-deformation} applied to
\(\partial_\Delta\) produces a canonical structured square-zero thickening
\[
\boxed{
    i_1:X\longrightarrow D^{(1)}_{X/S}
    \xrightarrow{\nu_1}
    X\times_SX
}
\]
with coefficient \(L_{X/S}\), classified specifically by
\(\partial_\Delta\).  Write
\[
    h_1:D^{(1)}_{X/S}\longrightarrow X
\]
for the composite with the first projection.  Then \(h_1i_1=\id_X\).
\end{construction}

\begin{proposition}[First-order neighbourhood maps to the intrinsic relation]
\label{prop:first-order-neighbourhood-to-relation}
There is a canonical morphism over \(X\times_SX\)
\[
\boxed{
    \ell_1:
    D^{(1)}_{X/S}
    \longrightarrow
    R_{X/S,1}.
}
\]
It carries the section \(i_1:X\to D^{(1)}_{X/S}\) to the diagonal degeneracy
\(X\to R_{X/S,1}\), and is compatible with arbitrary base change of represented morphisms.
\end{proposition}

\begin{proof}
The morphism \(i_1:X\to D^{(1)}_{X/S}\) is a one-stage structured spectral square-zero thickening.
Apply Proposition~\ref{prop:finite-infinitesimal-factorization-completion} to the diagram
\[
    X\xrightarrow{i_1}D^{(1)}_{X/S}
    \xrightarrow{\nu_1}X\times_SX.
\]
By Theorem~\ref{thm:relation-simplices-as-completions}, the target intrinsic completion is
\[
    \widehat{X\times_SX}^{\,\dR}_{X/S}
    \simeq
    R_{X/S,1}.
\]
The construction of the factorization preserves the given map from \(X\), hence identifies the
section with the diagonal degeneracy.  Base-change compatibility follows from base change of the
cotangent complex, structured square-zero extensions, and intrinsic completions.
\end{proof}

\begin{definition}[Stabilized cotangent coefficient]
\label{def:stabilized-cotangent-coefficient}
Define
\[
\boxed{
    \mathbb L^{\mathrm{st}}_{X/S}
    :=
    \fib
    \left(
        \Pi_{h_1}h_1^*\mathbb O_X^{\add}
        \longrightarrow
        \mathbb O_X^{\add}
    \right)
    \in
    \mathcal E_X,
}
\]
where the displayed map is restriction along the section \(i_1\).
\end{definition}

\begin{proposition}[The stabilized cotangent coefficient realizes the cotangent sheaf]
\label{prop:stabilized-cotangent-is-cotangent}
Under the affine-point presentation
\[
    \mathcal E_X
    \simeq
    \operatorname{Shv}_{\Nis}
    \left(\mathcal C^{\mathrm{aff}}_{/X};\Sp\right),
\]
the object \(\mathbb L^{\mathrm{st}}_{X/S}\) is the spectrum-valued sheaf
\[
    (U=\Spec A\to X)
    \longmapsto
    \Gamma(U,L_{X/S}|_U).
\]
In particular, for every affine point \(U\to X\),
\[
\boxed{
    \mathbb L^{\mathrm{st}}_{X/S}(U)
    \simeq
    \Gamma(U,L_{X/S}).
}
\]
The equivalences are natural under affine restriction and arbitrary base change.
\end{proposition}

\begin{proof}
Let \(U=\Spec A\to X\) be an affine point and base change the structured thickening
\(i_1:X\to D^{(1)}_{X/S}\) along \(U\to X\).  The global square-zero base-change theorem of
Chapter~\ref{chap:cotangent-spectral-deformation} identifies the result with the structured
square-zero extension
\[
    A'\longrightarrow A
\]
whose coefficient is the connective \(A\)-module
\[
    M:=\Gamma(U,L_{X/S}|_U).
\]
The defining fibre sequence of a structured square-zero extension is
\[
    M\longrightarrow A'\longrightarrow A.
\]
By Proposition~\ref{prop:additive-scalar-represents-functions}, evaluation of
\(\Pi_{h_1}h_1^*\mathbb O_X^{\add}\) on \(U\) is \(A'\), while evaluation of
\(\mathbb O_X^{\add}\) is \(A\).  Evaluating the defining fibre of
\(\mathbb L^{\mathrm{st}}_{X/S}\) therefore gives
\[
    \mathbb L^{\mathrm{st}}_{X/S}(U)
    \simeq
    \fib(A'\to A)
    \simeq
    M.
\]
Structured square-zero base change supplies the compatibility with further affine restriction and
with arbitrary base change of \(X\to S\).  Affine density then identifies the entire
spectrum-valued sheaf as stated.
\end{proof}

\begin{corollary}[Stable sections of the stabilized cotangent coefficient]
\label{cor:stabilized-cotangent-global-sections}
There is a canonical equivalence
\[
\boxed{
    \Gamma_X^{\mathrm{st}}
    \left(
        \mathbb L^{\mathrm{st}}_{X/S}
    \right)
    \simeq
    \Gamma(X,L_{X/S}).
}
\]
More generally, for every represented base change \(c:T\to S\), with
\(X_T:=X\times_ST\) and structural morphism \(p_T:X_T\to T\), Beck--Chevalley and cotangent
base change identify
\[
    c^*\Pi_p\mathbb L^{\mathrm{st}}_{X/S}
    \simeq
    \Pi_{p_T}\mathbb L^{\mathrm{st}}_{X_T/T},
\]
and hence
\[
\boxed{
    \Gamma_T^{\mathrm{st}}
    \left(
        c^*\Pi_p\mathbb L^{\mathrm{st}}_{X/S}
    \right)
    \simeq
    \Gamma(X_T,L_{X_T/T}).
}
\]
\end{corollary}

\begin{proof}
Under the affine-point presentation of \(\mathcal E_X\), Proposition~\ref{prop:stabilized-cotangent-is-cotangent}
identifies the value of \(\mathbb L^{\mathrm{st}}_{X/S}\) at an affine point
\(U=\Spec A\to X\) with
\[
    \Gamma(U,L_{X/S}|_U).
\]
Affine density gives
\[
    \mathbf1_X
    \simeq
    \colim_{(U\to X)}
    \Sigma^\infty_{+,X}U.
\]
Therefore
\[
\begin{aligned}
    \Gamma_X^{\mathrm{st}}
    \left(
        \mathbb L^{\mathrm{st}}_{X/S}
    \right)
    &\simeq
    \limop_{(U\to X)^{\op}}
    \Gamma(U,L_{X/S}|_U)
    \\
    &\simeq
    \Gamma(X,L_{X/S}),
\end{aligned}
\]
where the second equivalence is affine descent for the quasi-coherent cotangent complex.

For the base-change statement, right Beck--Chevalley gives
\[
    c^*\Pi_p\mathbb L^{\mathrm{st}}_{X/S}
    \simeq
    \Pi_{p_T}g^*\mathbb L^{\mathrm{st}}_{X/S},
\]
where \(g:X_T\to X\).  Proposition~\ref{prop:stabilized-cotangent-is-cotangent} and global cotangent
base change identify
\[
    g^*\mathbb L^{\mathrm{st}}_{X/S}
    \simeq
    \mathbb L^{\mathrm{st}}_{X_T/T}.
\]
Applying stable sections and the first assertion over \(T\) gives the final equivalence.
\end{proof}

\begin{construction}[First-order linearization of normalized one-cochains]
\label{constr:first-order-linearization}
Restriction of scalar functions along
\[
    \ell_1:D^{(1)}_{X/S}\longrightarrow R_{X/S,1}
\]
gives
\[
    \underline{\mathbf C}_{\dR}^{\,1}(X/S)
    \longrightarrow
    \Pi_p\Pi_{h_1}h_1^*\mathbb O_X^{\add}.
\]
The two maps to degree zero are both restriction along the diagonal section.  Taking the fibres of
those restriction maps therefore defines the canonical first-order linearization
\[
\boxed{
    \lambda^1_{X/S}:
    \underline{\OmegaInf}^{\,1}(X/S)
    \longrightarrow
    \Pi_p\mathbb L^{\mathrm{st}}_{X/S}.
}
\]
\end{construction}

\begin{construction}[Stabilized universal derivation]
\label{constr:stabilized-universal-derivation}
The stabilized universal derivation is naturally parameterized over the base \(S\).  Let
\[
    t:T=\Spec R\longrightarrow S
\]
be an affine point and put
\[
    X_T:=X\times_ST.
\]
Right Beck--Chevalley, Proposition~\ref{prop:additive-scalar-represents-functions}, and
Corollary~\ref{cor:stabilized-cotangent-global-sections} give canonical equivalences
\[
\boxed{
    \Gamma_T^{\mathrm{st}}
    \left(
        t^*\Pi_p\mathbb O_X^{\add}
    \right)
    \simeq
    \Gamma(X_T,\mathcal O_{X_T})
}
\]
and
\[
\boxed{
    \Gamma_T^{\mathrm{st}}
    \left(
        t^*\Pi_p\mathbb L^{\mathrm{st}}_{X/S}
    \right)
    \simeq
    \Gamma(X_T,L_{X_T/T}).
}
\]
The global cotangent theory of Chapter~\ref{chap:cotangent-spectral-deformation} supplies the
universal relative derivation of structure sheaves on \(X_T\).  Forgetting module-linearity and
taking global sections gives a canonical map of spectra
\[
    \Gamma(X_T,\mathcal O_{X_T})
    \longrightarrow
    \Gamma(X_T,L_{X_T/T}).
\]
These maps are natural in affine points \(T\to S\): for a morphism \(T'\to T\) over \(S\),
global cotangent base change identifies
\[
    L_{X_{T'}/T'}
    \simeq
    L_{X_T/T}|_{X_{T'}},
\]
and the universal derivation commutes with this pullback.  By the affine-point presentation of
\(\mathcal E_S\), the resulting natural transformation defines a canonical morphism
\[
\boxed{
    d^{\mathrm{st}}_{X/S}:
    \Pi_p\mathbb O_X^{\add}
    \longrightarrow
    \Pi_p\mathbb L^{\mathrm{st}}_{X/S}
}
\]
in \(\mathcal E_S\).  It is compatible with arbitrary represented base change.
\end{construction}

\begin{theorem}[The infinitesimal differential linearizes to the universal derivation]
\label{thm:infinitesimal-differential-linearizes}
The composite
\[
    \Pi_p\mathbb O_X^{\add}
    =
    \underline{\OmegaInf}^{\,0}(X/S)
    \xrightarrow{d_{\inf}^0}
    \underline{\OmegaInf}^{\,1}(X/S)
    \xrightarrow{\lambda^1_{X/S}}
    \Pi_p\mathbb L^{\mathrm{st}}_{X/S}
\]
is canonically equivalent to the parameterized stabilized universal derivation:
\[
\boxed{
    \lambda^1_{X/S}d_{\inf}^0
    \simeq
    d^{\mathrm{st}}_{X/S}.
}
\]
For every affine base test \(T=\Spec R\to S\), after base change to \(T\) and restriction to an
affine open \(U=\Spec A\subseteq X_T\), the induced map is the universal derivation
\[
\boxed{
    A\longrightarrow L_{A/R}.
}
\]
\end{theorem}

\begin{proof}
By Proposition~\ref{prop:degree-zero-normalized-differential}, the morphism
\(d_{\inf}^0\) is canonically the factorization through
\[
    \underline{\OmegaInf}^{\,1}(X/S)
    =
    \fib(s^0)
\]
of the actual coface difference
\[
    d^0-d^1:
    \underline{\mathbf C}_{\dR}^{\,0}(X/S)
    \longrightarrow
    \underline{\mathbf C}_{\dR}^{\,1}(X/S).
\]
By Convention~\ref{conv:monadic-de-rham-cech-face-order}, this is pullback
along the second projection minus pullback along the first projection.

We first perform the affine calculation.  Let
\[
    S=\Spec R,
    \qquad
    X=\Spec A,
    \qquad
    M:=L_{A/R},
\]
and put
\[
    B:=A\otimes_RA.
\]
Write
\[
    j_0,j_1:A\rightrightarrows B,
    \qquad
    j_0(a)=a\otimes1,
    \qquad
    j_1(a)=1\otimes a,
\]
and let
\[
    \mu:B\longrightarrow A
\]
be multiplication.  Regard \(B\) as an \(A\)-algebra through \(j_0\).  Cotangent
base change gives
\[
    L_{B/A}
    \simeq
    B\otimes_{A,j_1}M.
\]
Consequently
\[
\boxed{
    P:=A\otimes_BL_{B/A}
    \xrightarrow{\simeq}
    M.
}
\]

The composite
\[
    A\xrightarrow{j_0}B\xrightarrow{\mu}A
\]
is the identity.  Its transitivity triangle is therefore
\[
    P
    \longrightarrow
    L_{A/A}\simeq0
    \longrightarrow
    L_{A/B}.
\]
Let
\[
\boxed{
    \partial:
    L_{A/B}
    \xrightarrow{\simeq}
    P[1]
    \xrightarrow{\simeq}
    M[1]
}
\]
be the connecting equivalence, with the same rotation convention as
\(\partial_\Delta\) in
Construction~\ref{constr:first-diagonal-thickening}.  This is exactly the
affine classifying morphism obtained by restricting \(\partial_\Delta\).

Let
\[
    q:A'\longrightarrow A
\]
be the structured square-zero \(B\)-algebra extension with coefficient \(M\)
classified by \(\partial\).  Write
\[
    \beta:B\longrightarrow A'
\]
for its \(B\)-algebra structure map.  The two tensor-factor maps induce
sections
\[
    \alpha_i:=\beta j_i:
    A\longrightarrow A',
    \qquad i=0,1,
\]
of \(q\).

Restrict the structured extension along \(j_0:A\to B\).  Its relative
classifying cotangent complex is \(L_{A/A}\simeq0\).  The structured
square-zero classification therefore identifies the resulting extension,
through a contractible space of compatible choices, with the split
\(A\)-algebra square-zero extension
\[
    A'\simeq A\oplus M
\]
in such a way that
\[
    \alpha_0(a)=(a,0).
\]
There is consequently a unique \(R\)-derivation
\[
    \delta:A\longrightarrow M
\]
for which
\[
    \alpha_1(a)=(a,\delta(a)).
\]

Under the same splitting, the map \(\beta:B\to A\oplus M\) has the form
\[
    \beta(b)=(\mu(b),D(b)),
\]
where
\[
    D:B\longrightarrow M
\]
is a derivation relative to the first tensor factor \(j_0\); equivalently,
\(D j_0=0\).  Let
\[
    \overline D:
    P=A\otimes_BL_{B/A}
    \longrightarrow
    M
\]
be the \(A\)-linear map classifying this relative derivation.

Apply \(\map_{\operatorname{Mod}_A}(-,M[1])\) to the transitivity triangle
\[
    P\longrightarrow0\longrightarrow L_{A/B}.
\]
The connecting equivalence gives a canonical equivalence of mapping spaces
\[
\boxed{
    \Map_{\operatorname{Mod}_A}(P,M)
    \xrightarrow{\simeq}
    \Map_{\operatorname{Mod}_A}(L_{A/B},M[1]).
}
\]
Under structured square-zero classification, the extension \(A'\to A\)
corresponds on the right to the classifying map \(\partial\).  By the
definition of \(\partial\), its inverse image on the left is the identity of
\(P\), after the fixed identification \(P\simeq M\).  The relative derivation
carried by the split \(B\)-algebra structure is therefore classified by
\[
\boxed{
    \overline D=\id_M.
}
\]

The derivation \(\delta\) is the restriction of \(D\) along the second tensor
factor:
\[
    \delta=Dj_1.
\]
Under the cotangent base-change identification
\[
    P=A\otimes_BL_{B/A}\simeq L_{A/R}=M,
\]
restriction along \(j_1\) carries the classifying map \(\overline D\) to the
\(A\)-linear morphism
\[
    L_{A/R}\longrightarrow M
\]
classifying \(\delta\).  Since \(\overline D=\id_M\), this classifying
morphism is the identity of \(L_{A/R}\).  By the universal property of the
cotangent complex,
\[
\boxed{
    \delta=d_{A/R}:A\longrightarrow L_{A/R}.
}
\]
Equivalently,
\[
\boxed{
    \alpha_1-\alpha_0=d_{A/R}
}
\]
after identifying the fibre of \(A'\to A\) with \(M\).

By the fixed \v{C}ech face convention,
\[
    d^0-d^1
    =
    j_1-j_0
\]
on affine scalar functions.  Restriction along
\[
    \ell_1:D^{(1)}_{X/S}\longrightarrow R_{X/S,1}
\]
takes its normalized factorization to
\(\alpha_1-\alpha_0\).  Hence, in the affine case,
\[
\boxed{
    \lambda^1_{X/S}d_{\inf}^0
    =
    d_{A/R}.
}
\]
The sign is therefore fixed by the ordered \v{C}ech faces and by the actual
connecting morphism in the transitivity triangle; no independent sign choice
remains.

For the general represented morphism \(X\to S\), evaluate the two morphisms in
\(\mathcal E_S\) at an arbitrary affine point \(T=\Spec R\to S\).
Parameterized base change reduces the comparison to the map
\[
    \Gamma(X_T,\mathcal O_{X_T})
    \longrightarrow
    \Gamma(X_T,L_{X_T/T}).
\]
The diagonal square-zero construction, the map \(\ell_1\), the first-order
linearization, and the cotangent complex commute with restriction to affine
opens because they are formed from inverse image, finite limits, cotangent
base change, and the global structured square-zero construction.  Hence, on
every affine open \(U=\Spec A\subseteq X_T\), the affine calculation above
identifies the restriction of \(\lambda^1d_{\inf}^0\) with
\[
    d_{A/R}:A\longrightarrow L_{A/R}.
\]
The global cotangent complex is computed by affine descent, so these compatible
affine comparisons identify the resulting global map with the universal
relative derivation
\[
    \Gamma(X_T,\mathcal O_{X_T})
    \longrightarrow
    \Gamma(X_T,L_{X_T/T}).
\]
By Construction~\ref{constr:stabilized-universal-derivation}, this is the value
of \(d^{\mathrm{st}}_{X/S}\) at the affine point \(T\to S\).  Affine-point
density in \(\mathcal E_S\) therefore gives
\[
    \lambda^1_{X/S}d_{\inf}^0
    \simeq
    d^{\mathrm{st}}_{X/S}.
\]
\end{proof}

\begin{proposition}[Higher infinitesimal terms survive normalization]
\label{prop:normalized-one-cochains-strictly-nonlinear}
Let
\[
    S=\Spec H\mathbb Q,
    \qquad
    X=\Spec H(\mathbb Q[t])
\]
be the discrete spectral enhancement of the ordinary affine line over \(\mathbb Q\).  Then the
first-order linearization
\[
    \pi_0\Gamma_S^{\mathrm{st}}
    \left(\underline{\OmegaInf}^{\,1}(X/S)\right)
    \longrightarrow
    \pi_0\Gamma(X,L_{X/S})
\]
has nonzero kernel.  In particular, the canonical first-order linearization
\[
\boxed{
    \lambda^1_{X/S}:
    \underline{\OmegaInf}^{\,1}(X/S)
    \longrightarrow
    \Pi_p\mathbb L^{\mathrm{st}}_{X/S}
}
\]
is not an equivalence.
\end{proposition}

\begin{proof}
Write
\[
    Y=X\times_SX
    =
    \Spec H(\mathbb Q[t_0,t_1])
\]
and put \(u=t_1-t_0\).  Consider
\[
    T_2
    :=
    \Spec H\left(\mathbb Q[t,u]/(u^3)\right)
\]
with the map to \(Y\) defined by
\[
    t_0\longmapsto t,
    \qquad
    t_1\longmapsto t+u.
\]
The closed immersion \(X\to T_2\) factors as
\[
    X
    \longrightarrow
    \Spec H\left(\mathbb Q[t,u]/(u^2)\right)
    \longrightarrow
    T_2.
\]
Both stages have square-zero kernel: respectively \((u)/(u^2)\) and
\((u^2)/(u^3)\).  By the discrete square-zero recognition theorem of the
preceding infinitesimal geometry, they are structured spectral square-zero
thickenings.  Hence \(X\to T_2\) is a finite spectral infinitesimal thickening,
and Proposition~\ref{prop:finite-infinitesimal-factorization-completion} gives
a factorization
\[
    T_2\longrightarrow
    \widehat Y^{\,\dR}_{X/S}
    \simeq
    R_{X/S,1}.
\]

The polynomial
\[
    (t_1-t_0)^2
\]
defines a represented scalar morphism
\[
    Y=X\times_SX
    \longrightarrow
    \mathbb A^1_S
    \simeq
    \Omega^\infty\mathbb O_S^{\add}.
\]
Composing with the canonical morphism
\[
    R_{X/S,1}\longrightarrow Y
\]
therefore gives a section of \(a_1^*\mathbb O_S^{\add}\) over the generally
nonrepresented object \(R_{X/S,1}\).  Equivalently, by
Proposition~\ref{prop:modal-de-rham-mapping-spectrum}, it determines a point of
\[
    \Omega^\infty
    \Gamma_S^{\mathrm{st}}
    \left(
        \underline{\mathbf C}_{\dR}^{\,1}(X/S)
    \right).
\]
Denote its class in \(\pi_0\) by \(u^2\).  Its restriction along the diagonal
degeneracy is zero.  The defining fibre sequence
\[
    \underline{\OmegaInf}^{\,1}(X/S)
    \longrightarrow
    \underline{\mathbf C}_{\dR}^{\,1}(X/S)
    \xrightarrow{s^0}
    \underline{\mathbf C}_{\dR}^{\,0}(X/S)
\]
therefore supplies a lift
\[
    \omega_2
    \in
    \pi_0\Gamma_S^{\mathrm{st}}
    \left(\underline{\OmegaInf}^{\,1}(X/S)\right)
\]
whose image in degree-one cochains is the class \(u^2\).  Its further
restriction to \(T_2\) is
\[
    u^2\in\mathbb Q[t,u]/(u^3),
\]
which is nonzero.  Hence the degree-one cochain class is nonzero, and therefore
\[
\boxed{
    \omega_2\neq0.
}
\]
No uniqueness of the lift is needed.

It remains to prove that every such lift is killed by first-order
linearization.  Put
\[
    A=\mathbb Q[t],
    \qquad
    M=L_{A/\mathbb Q}\simeq A\,dt,
\]
and let
\[
    q:A'\longrightarrow A
\]
be the affine structured square-zero extension representing the first
structured diagonal neighbourhood.  Its defining fibre sequence of spectra is
\[
    M\longrightarrow A'\xrightarrow{q}A.
\]
Because \(A\) is discrete,
\[
    \pi_1(A)=0.
\]
The long exact homotopy sequence therefore gives an injection
\[
\boxed{
    \pi_0M\hookrightarrow\pi_0A'.
}
\]

On \(D^{(1)}_{X/S}\), the function \(u=t_1-t_0\) lies in the square-zero
coefficient ideal: its restriction along the section \(i_1\) is zero, and the
multiplication on the coefficient ideal is null.  Consequently the restriction
of \(u^2\) to \(A'\) is zero in \(\pi_0A'\).

By construction of \(\lambda^1_{X/S}\), the image
\[
    \lambda^1_{X/S}(\omega_2)
    \in
    \pi_0M
    \simeq
    \pi_0\Gamma(X,L_{X/S})
\]
maps under the displayed injection \(\pi_0M\hookrightarrow\pi_0A'\) to the
restriction of the degree-one cochain \(u^2\) to the first structured diagonal
neighbourhood.  That restriction is zero, so injectivity forces
\[
\boxed{
    \lambda^1_{X/S}(\omega_2)=0.
}
\]
Thus the first-order linearization has nonzero kernel on \(\pi_0\) and cannot
be an equivalence.
\end{proof}

\begin{remark}[Normalized infinitesimal cochains are not cotangent forms]
\label{rem:normalized-cochains-not-cotangent-forms}
The preceding proposition shows that the distinction is mathematical, not terminological.  The
normalized relation cochains retain quadratic and higher finite-infinitesimal information, while the
cotangent complex is their first-order linearization.  Accordingly, this chapter does not identify
\(\underline{\OmegaInf}^{\,q}(X/S)\) with \(\wedge^qL_{X/S}\), and the
\v{C}ech-totalization filtration is not identified with the Hodge filtration.  Any such comparison
requires additional linearization or HKR-type input beyond the definitions constructed here.
\end{remark}


\section{Multiplicative structures}
\label{sec:multiplicative-modal-de-rham}

The cosimplicial de Rham object is already a cosimplicial \(E_\infty\)-algebra.  Two multiplicative
statements extracted from it must be distinguished.  Partial and complete totalizations are limits of
commutative algebras and therefore carry canonical \(E_\infty\)-algebra structures.  Normalized
infinitesimal cochains admit the ordered Alexander--Whitney maps; in this chapter we use their associative, unital,
and Leibniz properties after passage to the homotopy category.  We do not assert that these ordered
maps lift to an \(E_1\)-algebra object in coherent cochains, and consequently we do not assert a
multiplicative \v{C}ech-totalization filtration.

\begin{proposition}[Monoidal complete-filtration correspondence]
\label{prop:monoidal-complete-filtration-correspondence}
Let \(\mathcal D\) be a stable presentably symmetric monoidal \(\infty\)-category whose tensor
product preserves colimits separately.  Complete decreasing filtrations whose associated graded is
concentrated in nonnegative degrees carry the completed Day-convolution tensor induced by addition of
filtration degrees.  Under Theorem~\ref{thm:complete-filtered-coherent-cochains}, this tensor
transports to a symmetric monoidal structure on coherent nonnegative cochains.  More precisely,
\[
\boxed{
    F\widehat\otimes G
    :=
    \widehat{F\otimes_{\mathrm{Day}}G},
}
\]
and
\[
\boxed{
    \gr^n(F\widehat\otimes G)
    \simeq
    \bigoplus_{a+b=n}
    \gr^aF\otimes\gr^bG.
}
\]
If
\[
    C=\Phi_{\mathcal D}(F),
    \qquad
    D=\Phi_{\mathcal D}(G),
\]
then the transported tensor satisfies
\[
\boxed{
    (C\otimes D)^n
    \simeq
    \bigoplus_{a+b=n}C^a\otimes D^b.
}
\]
Thus any filtered \(E_1\)-algebra structure, when separately supplied, induces degree-additive
multiplication on associated graded objects and a coherent derivation structure on its connecting
differential.
\end{proposition}

\begin{proof}
Let
\[
    \operatorname{Seq}(\mathcal D)
    :=
    \Fun(\mathbb Z^{\op},\mathcal D)
\]
with Day convolution for addition on \(\mathbb Z\).  Since the tensor product in
\(\mathcal D\) preserves colimits separately, Day convolution exists and
preserves colimits separately.

Let
\[
    L_{\mathrm{comp}}:
    \operatorname{Seq}(\mathcal D)
    \longrightarrow
    \operatorname{Fil}^{\mathrm{comp}}(\mathcal D)
\]
denote the completion reflector of
Lemma~\ref{lem:complete-filtration-reflection}.  A morphism \(u:F\to G\)
becomes an equivalence after completion if and only if it is an equivalence on
every associated-graded term.  Indeed, completion preserves associated
graded, so one implication is immediate.  Conversely, if \(u\) is a graded
equivalence, then \(L_{\mathrm{comp}}u\) is a graded equivalence between
complete filtrations and is therefore an equivalence by the graded-detection
statement of Lemma~\ref{lem:complete-filtration-reflection}.

We next compute associated graded for Day convolution.  If
\[
    H:=F\otimes_{\mathrm{Day}}G,
\]
the Day-colimit formula for the ordered monoid \(\mathbb Z^{\op}\) gives
\[
\boxed{
    H^n
    \simeq
    \colim_{a+b\geq n}
    F^a\otimes G^b.
}
\]
The transition
\[
    H^{n+1}\longrightarrow H^n
\]
is induced by the inclusion of the subdiagram \(a+b\geq n+1\) into the
subdiagram \(a+b\geq n\).  Its cofiber is therefore obtained from the
boundary \(a+b=n\) by quotienting the two incoming directions from higher
filtration degree.  For a fixed pair \(a+b=n\), the corresponding latching
map is
\[
\begin{aligned}
    &
    \left(F^{a+1}\otimes G^b\right)
    \coprod_{F^{a+1}\otimes G^{b+1}}
    \left(F^a\otimes G^{b+1}\right)
    \\
    &\hspace{4cm}\longrightarrow
    F^a\otimes G^b.
\end{aligned}
\]
Because tensor product is exact separately, the cofiber of this map is
canonically
\[
    \gr^aF\otimes\gr^bG.
\]
After the entire higher region \(a+b\geq n+1\) has been collapsed, distinct
boundary pairs \(a+b=n\) have no remaining transition morphisms between them.
Hence the boundary cofibers add independently, giving
\[
\boxed{
    \gr^n(F\otimes_{\mathrm{Day}}G)
    \simeq
    \bigoplus_{a+b=n}
    \gr^aF\otimes\gr^bG.
}
\]

It follows immediately that graded equivalences form a tensor ideal for Day
convolution.  The symmetric-monoidal localization theorem therefore equips
the complete filtered category with the completed Day-convolution product
\[
    F\widehat\otimes G
    :=
    L_{\mathrm{comp}}
    \left(
        F\otimes_{\mathrm{Day}}G
    \right)
    =
    \widehat{F\otimes_{\mathrm{Day}}G}.
\]
Since completion preserves associated graded,
\[
    \gr^n(F\widehat\otimes G)
    \simeq
    \bigoplus_{a+b=n}
    \gr^aF\otimes\gr^bG.
\]

If the associated gradeds of \(F\) and \(G\) are concentrated in
nonnegative degrees, then for \(n<0\) every summand on the right vanishes,
because at least one of \(a,b\) is negative.  Thus
\(\operatorname{Fil}^{\mathrm{comp}}_{\geq0}(\mathcal D)\) is closed under
\(\widehat\otimes\), and the completed Day-convolution unit belongs to this
subcategory.

Transport the completed tensor product across
Theorem~\ref{thm:complete-filtered-coherent-cochains}.  If
\[
    C=\Phi_{\mathcal D}(F),
    \qquad
    D=\Phi_{\mathcal D}(G),
\]
then
\[
\begin{aligned}
    \Phi_{\mathcal D}(F\widehat\otimes G)^n
    &\simeq
    \gr^n(F\widehat\otimes G)[n]
    \\
    &\simeq
    \bigoplus_{a+b=n}
    \left(
        \gr^aF\otimes\gr^bG
    \right)[a+b]
    \\
    &\simeq
    \bigoplus_{a+b=n}
    \gr^aF[a]\otimes\gr^bG[b]
    \\
    &\simeq
    \bigoplus_{a+b=n}
    C^a\otimes D^b.
\end{aligned}
\]
This is the asserted degree-additive tensor product on coherent cochains.

Finally, suppose \(F\) is an \(E_1\)-algebra for the completed tensor.  Its
multiplication induces
\[
    \gr^pF\otimes\gr^qF
    \longrightarrow
    \gr^{p+q}F.
\]
Under the coherent-cochain correspondence this multiplication is a morphism
of coherent cochain objects.  Compatibility with the first structure maps,
together with the suspension symmetry in the identification
\(\Phi_{\mathcal D}(F)^p=\gr^pF[p]\), gives the standard tensor-cochain sign
\[
    d(x\otimes y)
    =
    dx\otimes y
    +
    (-1)^p x\otimes dy.
\]
Postcomposing with multiplication gives, in the homotopy category,
\[
    d(xy)
    =
    d(x)y+(-1)^p x\,d(y)
\]
for \(x\) of coherent degree \(p\).  No multiplicative structure on the
specific \v{C}ech-totalization filtration is produced unless one is separately
supplied.
\end{proof}

\begin{lemma}[Consecutive-block matching-cube factorization]
\label{lem:aw-consecutive-block-matching-cube}
Let \(\mathcal D\) be a stable symmetric monoidal \(\infty\)-category whose
tensor product is exact separately in each variable.  Let
\[
    C^\bullet,D^\bullet,E^\bullet\in\Fun(\Delta,\mathcal D)
\]
and let
\[
    \mu^\bullet:
    C^\bullet\otimes D^\bullet
    \longrightarrow
    E^\bullet
\]
be a natural pairing.  For every \(p,q\geq0\), with \(N=p+q\), the
consecutive-block map
\[
    C^p\otimes D^q
    \longrightarrow
    E^N
\]
canonically induces an actual morphism in \(\mathcal D\)
\[
\boxed{
    N^pC\otimes N^qD
    \longrightarrow
    N^NE.
}
\]
The factorization is natural in the pairing and retains the complete coherent
zero-cone data over the target codegeneracy matching diagram.
\end{lemma}

\begin{proof}
For \(r\geq0\), put
\[
    I_r:=\{0,\ldots,r-1\}.
\]
For every subset \(A\subseteq I_r\), let
\[
    \sigma_A:[r]\twoheadrightarrow[r-|A|]
\]
be the quotient order map for the equivalence relation generated by
\[
    i\sim i+1
    \qquad(i\in A).
\]
Equivalently,
\[
    \sigma_A(i)=\sigma_A(i+1)
    \quad\Longleftrightarrow\quad
    i\in A.
\]
Every order-preserving surjection out of \([r]\) is uniquely of this form.
If \(A\subseteq A'\), there is a unique order-preserving surjection
\[
    \rho_{A,A'}:
    [r-|A|]\twoheadrightarrow[r-|A'|]
\]
such that
\[
\boxed{
    \sigma_{A'}
    =
    \rho_{A,A'}\sigma_A.
}
\]

For a cosimplicial object \(C^\bullet\), define its elementary codegeneracy
cube
\[
    K_r(C):\mathcal P(I_r)\longrightarrow\mathcal D
\]
by
\[
    K_r(C)(A):=C^{r-|A|}
\]
and by the maps \(C(\rho_{A,A'})\) for inclusions \(A\subseteq A'\).
The punctured cube
\[
    \mathcal P(I_r)\setminus\{\varnothing\}
\]
is canonically the category of nonidentity order-preserving surjections out of
\([r]\).  Hence its limit is the codegeneracy matching object and
\[
\boxed{
    \operatorname{tfib}K_r(C)
    \simeq
    N^rC.
}
\]

Now fix \(p,q\geq0\), put \(N=p+q\), and identify
\[
    I_N
    =
    I_p\amalg(p+I_q).
\]
Thus
\[
    \mathcal P(I_p)\times\mathcal P(I_q)
    \xrightarrow{\simeq}
    \mathcal P(I_N),
    \qquad
    (A,B)\longmapsto A\cup(p+B).
\]
For \(A\subseteq I_p\) and \(B\subseteq I_q\), put
\[
    r_A:=p-|A|,
    \qquad
    s_B:=q-|B|,
    \qquad
    n_{A,B}:=r_A+s_B.
\]
Let
\[
    \iota_{1,A,B}:[r_A]\hookrightarrow[n_{A,B}],
    \qquad
    \iota_{2,A,B}:[s_B]\hookrightarrow[n_{A,B}]
\]
be the order-preserving injections onto the consecutive vertex blocks,
explicitly
\[
    \iota_{1,A,B}(i)=i,
    \qquad
    \iota_{2,A,B}(j)=r_A+j.
\]
At the vertex \((A,B)\), define
\[
\begin{aligned}
    m_{A,B}:\,
    C^{r_A}\otimes D^{s_B}
    &\xrightarrow{
        C(\iota_{1,A,B})
        \otimes
        D(\iota_{2,A,B})
    }
    C^{n_{A,B}}\otimes D^{n_{A,B}}
    \\
    &\xrightarrow{\mu^{n_{A,B}}}
    E^{n_{A,B}}.
\end{aligned}
\]

We verify that the maps \(m_{A,B}\) form a morphism of \(N\)-cubes.  Let
\[
    A\subseteq A',
    \qquad
    B\subseteq B'.
\]
Write
\[
    \rho_A:=\rho_{A,A'},
    \qquad
    \rho_B:=\rho_{B,B'},
\]
and let
\[
    \rho_{A,B}^{A',B'}:
    [n_{A,B}]
    \twoheadrightarrow
    [n_{A',B'}]
\]
be the quotient map characterized by
\[
    \sigma_{A'\cup(p+B')}
    =
    \rho_{A,B}^{A',B'}
    \sigma_{A\cup(p+B)}.
\]
The restrictions of this quotient to the two consecutive blocks satisfy the
ordinal identities
\[
\boxed{
    \rho_{A,B}^{A',B'}\iota_{1,A,B}
    =
    \iota_{1,A',B'}\rho_A,
}
\]
and
\[
\boxed{
    \rho_{A,B}^{A',B'}\iota_{2,A,B}
    =
    \iota_{2,A',B'}\rho_B.
}
\]
Indeed, both sides of the first identity are the quotient of the first
consecutive vertex block by the relations indexed by \(A'\), and both sides
of the second identity are the quotient of the second consecutive vertex
block by the relations indexed by \(B'\).

Naturality of \(\mu^\bullet\) now gives
\[
\begin{aligned}
    E(\rho_{A,B}^{A',B'})m_{A,B}
    &=
    \mu^{n_{A',B'}}
    \bigl(
        C(\rho_{A,B}^{A',B'})
        \otimes
        D(\rho_{A,B}^{A',B'})
    \bigr)
    \\
    &\qquad\circ
    \bigl(
        C(\iota_{1,A,B})
        \otimes
        D(\iota_{2,A,B})
    \bigr)
    \\
    &=
    m_{A',B'}
    \bigl(
        C(\rho_A)\otimes D(\rho_B)
    \bigr).
\end{aligned}
\]
Hence the vertex maps assemble to a canonical morphism of \(N\)-cubes
\[
\boxed{
    K_p(C)\boxtimes K_q(D)
    \longrightarrow
    K_N(E).
}
\]

Total fibre is computed by iterating ordinary fibres in the cube coordinates.
Because tensor product is exact separately in each variable, iterated fibre
commutes with tensor product.  Fubini for the two sets of cube coordinates
therefore gives
\[
\begin{aligned}
    \operatorname{tfib}
    \bigl(
        K_p(C)\boxtimes K_q(D)
    \bigr)
    &\simeq
    \operatorname{tfib}K_p(C)
    \otimes
    \operatorname{tfib}K_q(D)
    \\
    &\simeq
    N^pC\otimes N^qD.
\end{aligned}
\]
The target total fibre is
\[
    \operatorname{tfib}K_N(E)
    \simeq
    N^NE.
\]
Taking total fibres of the cube morphism therefore gives the required
canonical map
\[
    N^pC\otimes N^qD
    \longrightarrow
    N^NE.
\]
Because the construction is made at the level of the complete matching cubes,
the induced map includes the compatible zero-cone data for all codegeneracies
and all their higher overlap coherences.  Naturality follows from functoriality
of the quotient maps, the pairing, and total fibres.
\end{proof}

\begin{proposition}[Ordered Alexander--Whitney normalization]
\label{prop:ordered-monoidal-stable-dold-kan}
Let \(\mathcal D\) be a stable presentably symmetric monoidal \(\infty\)-category whose tensor
product preserves colimits separately, and let
\[
    C^\bullet\in\Fun(\Delta,\CAlg(\mathcal D)).
\]
For \(p,q\geq0\), there is a natural map
\[
\boxed{
    \smile_{\mathrm{AW}}:
    N^pC\otimes N^qC
    \longrightarrow
    N^{p+q}C.
}
\]
After passage to the homotopy category these maps are associative and unital.  With the Dold--Kan
sign convention of Convention~\ref{conv:monadic-de-rham-coherent-cochains}, the normalized
differential satisfies
\[
\boxed{
    d(x\smile_{\mathrm{AW}}y)
    =
    dx\smile_{\mathrm{AW}}y
    +
    (-1)^p x\smile_{\mathrm{AW}}dy
}
\]
for \(x\) of normalized degree \(p\).  In the homotopy category the binary product is represented by
\[
\boxed{
    x\smile_{\mathrm{AW}}y
    =
    \left(
        d^{p+q}\cdots d^{p+1}x
    \right)
    \cdot
    \left(
        (d^0)^py
    \right).
}
\]
The construction is natural in \(C^\bullet\) and in exact strong symmetric monoidal functors.
\end{proposition}

\begin{proof}
Apply Lemma~\ref{lem:aw-consecutive-block-matching-cube} with
\[
    C^\bullet=D^\bullet=E^\bullet
\]
and with \(\mu^\bullet\) the multiplication of the cosimplicial commutative
algebra.  This gives the actual morphism
\[
    N^pC\otimes N^qC
    \longrightarrow
    N^{p+q}C
\]
in \(\mathcal D\), including the coherent factorization through the full
codegeneracy matching fibre.

Put \(N=p+q\), and let
\[
    f_{p,q}:[p]\hookrightarrow[N],
    \qquad
    f_{p,q}(i)=i,
\]
and
\[
    b_{p,q}:[q]\hookrightarrow[N],
    \qquad
    b_{p,q}(j)=p+j.
\]
These are the two consecutive-block injections.  On the cosimplicial object,
\[
    C(f_{p,q})
    =
    d^{p+q}\cdots d^{p+1},
    \qquad
    C(b_{p,q})
    =
    (d^0)^p.
\]
Hence, after passage to the homotopy category, the composite of the normalized
product with \(N^{p+q}C\to C^{p+q}\) is represented by
\[
    x\smile_{\mathrm{AW}}y
    =
    \left(
        d^{p+q}\cdots d^{p+1}x
    \right)
    \cdot
    \left(
        (d^0)^py
    \right).
\]

For three normalized inputs of degrees \(p,q,r\), the two parenthesizations,
after composition with
\[
    N^{p+q+r}C\longrightarrow C^{p+q+r},
\]
are both induced by multiplication after the same three consecutive vertex
blocks
\[
    [0,p],
    \qquad
    [p,p+q],
    \qquad
    [p+q,p+q+r].
\]
The cosimplicial identities and associativity of multiplication identify
these two composites.  Stable Dold--Kan supplies the normalized--degenerate
splitting in the additive homotopy category, so
\[
    N^{m}C\longrightarrow C^{m}
\]
is split monic there for every \(m\).  Hence the two parenthesizations already
agree in \(N^{p+q+r}C\).  The same argument with the constant degree-zero unit
proves unitality.

It remains to prove the Leibniz identity.  For \(x\in N^pC\) and
\(y\in N^qC\), write
\[
    F_{p,q}(x,y)
    :=
    \mu^{p+q}
    \left(
        C(f_{p,q})x,
        C(b_{p,q})y
    \right)
    \in C^{p+q}.
\]
Let
\[
    \delta^k:[p+q]\hookrightarrow[p+q+1]
\]
denote the \(k\)-th coface injection.  For \(0\leq i\leq p\), the ordinal
identities
\[
    \delta^if_{p,q}
    =
    f_{p+1,q}\delta^i,
    \qquad
    \delta^ib_{p,q}
    =
    b_{p+1,q}
\]
give, by naturality of multiplication,
\[
\boxed{
    d^iF_{p,q}(x,y)
    =
    F_{p+1,q}(d^ix,y)
    \qquad(0\leq i\leq p).
}
\]
For \(1\leq j\leq q+1\), the ordinal identities
\[
    \delta^{p+j}f_{p,q}
    =
    f_{p,q+1},
    \qquad
    \delta^{p+j}b_{p,q}
    =
    b_{p,q+1}\delta^j
\]
give
\[
\boxed{
    d^{p+j}F_{p,q}(x,y)
    =
    F_{p,q+1}(x,d^jy)
    \qquad(1\leq j\leq q+1).
}
\]
At the unique break between the two consecutive blocks one has
\[
    f_{p+1,q}\delta^{p+1}
    =
    f_{p,q+1},
    \qquad
    b_{p+1,q}
    =
    b_{p,q+1}\delta^0,
\]
and therefore
\[
\boxed{
    F_{p+1,q}(d^{p+1}x,y)
    =
    F_{p,q+1}(x,d^0y).
}
\]

Now expand the alternating-coface differential:
\[
\begin{aligned}
    dF_{p,q}(x,y)
    &=
    \sum_{i=0}^{p}
    (-1)^i
    F_{p+1,q}(d^ix,y)
    \\
    &\qquad+
    \sum_{j=1}^{q+1}
    (-1)^{p+j}
    F_{p,q+1}(x,d^jy).
\end{aligned}
\]
On the other hand,
\[
\begin{aligned}
    F_{p+1,q}(dx,y)
    &+
    (-1)^pF_{p,q+1}(x,dy)
    \\
    &=
    \sum_{i=0}^{p+1}
    (-1)^i
    F_{p+1,q}(d^ix,y)
    \\
    &\qquad+
    \sum_{j=0}^{q+1}
    (-1)^{p+j}
    F_{p,q+1}(x,d^jy).
\end{aligned}
\]
The only terms in the second display not occurring in the first are
\[
    (-1)^{p+1}
    F_{p+1,q}(d^{p+1}x,y)
\]
and
\[
    (-1)^p
    F_{p,q+1}(x,d^0y),
\]
and they cancel by the break identity.  Hence
\[
    dF_{p,q}(x,y)
    =
    F_{p+1,q}(dx,y)
    +
    (-1)^pF_{p,q+1}(x,dy).
\]
Both sides factor through the normalized term \(N^{p+q+1}C\), by
Lemma~\ref{lem:aw-consecutive-block-matching-cube} and by the normalized
cochain differential.  Since
\[
    N^{p+q+1}C\longrightarrow C^{p+q+1}
\]
is split monic in the homotopy category, the equality after inclusion is
exactly
\[
    d(x\smile_{\mathrm{AW}}y)
    =
    dx\smile_{\mathrm{AW}}y
    +
    (-1)^p x\smile_{\mathrm{AW}}dy
\]
in the homotopy category.

Naturality in \(C^\bullet\) and in exact strong symmetric monoidal functors
follows from the naturality of the matching-cube construction, the
consecutive-block maps, stable Dold--Kan normalization, and multiplication.
\end{proof}

\begin{proposition}[Paired ordered Alexander--Whitney normalization]
\label{prop:paired-ordered-monoidal-stable-dold-kan}
Let \(\mathcal D\) be a stable presentably symmetric monoidal \(\infty\)-category whose tensor
product preserves colimits separately.  Let
\[
    C^\bullet,D^\bullet,E^\bullet\in\Fun(\Delta,\mathcal D)
\]
and let
\[
    \mu^\bullet:
    C^\bullet\otimes D^\bullet
    \longrightarrow
    E^\bullet
\]
be a natural pairing of cosimplicial objects.  For every \(p,q\geq0\), there is a canonical natural
map
\[
\boxed{
    \boxtimes_{\mathrm{AW}}:
    N^pC\otimes N^qD
    \longrightarrow
    N^{p+q}E.
}
\]
In the homotopy category it is represented by
\[
\boxed{
    x\boxtimes_{\mathrm{AW}}y
    =
    \mu^{p+q}
    \left(
        d^{p+q}\cdots d^{p+1}x,
        (d^0)^py
    \right).
}
\]
With the Dold--Kan sign convention of
Convention~\ref{conv:monadic-de-rham-coherent-cochains}, the normalized differentials satisfy
\[
\boxed{
    d_E(x\boxtimes_{\mathrm{AW}}y)
    =
    d_Cx\boxtimes_{\mathrm{AW}}y
    +
    (-1)^p
    x\boxtimes_{\mathrm{AW}}d_Dy
}
\]
in the homotopy category for \(x\) of normalized degree \(p\).  The construction is natural in the
pairing and in exact strong symmetric monoidal functors.
\end{proposition}

\begin{proof}
Apply Lemma~\ref{lem:aw-consecutive-block-matching-cube} to the given pairing
\(\mu^\bullet\).  This gives the canonical actual morphism
\[
    N^pC\otimes N^qD
    \longrightarrow
    N^{p+q}E.
\]

Put \(N=p+q\), and let
\[
    f_{p,q}:[p]\hookrightarrow[N],
    \qquad
    f_{p,q}(i)=i,
\]
and
\[
    b_{p,q}:[q]\hookrightarrow[N],
    \qquad
    b_{p,q}(j)=p+j.
\]
At the initial vertex of the matching-cube construction the map is
\[
    C^p\otimes D^q
    \xrightarrow{
        C(f_{p,q})\otimes D(b_{p,q})
    }
    C^N\otimes D^N
    \xrightarrow{\mu^N}
    E^N.
\]
Since
\[
    C(f_{p,q})
    =
    d^{p+q}\cdots d^{p+1},
    \qquad
    D(b_{p,q})
    =
    (d^0)^p,
\]
the normalized map is represented in the homotopy category by
\[
    x\boxtimes_{\mathrm{AW}}y
    =
    \mu^{p+q}
    \left(
        d^{p+q}\cdots d^{p+1}x,
        (d^0)^py
    \right).
\]

For the differential identity define
\[
    F_{p,q}(x,y)
    :=
    \mu^{p+q}
    \left(
        C(f_{p,q})x,
        D(b_{p,q})y
    \right)
    \in E^{p+q}.
\]
The ordinal calculations used in the preceding proposition depend only on
the two consecutive-block injections and therefore apply to the present
pairing without change.  Explicitly, for \(0\leq i\leq p\),
\[
    \delta^if_{p,q}
    =
    f_{p+1,q}\delta^i,
    \qquad
    \delta^ib_{p,q}
    =
    b_{p+1,q},
\]
while for \(1\leq j\leq q+1\),
\[
    \delta^{p+j}f_{p,q}
    =
    f_{p,q+1},
    \qquad
    \delta^{p+j}b_{p,q}
    =
    b_{p,q+1}\delta^j.
\]
At the break,
\[
    f_{p+1,q}\delta^{p+1}
    =
    f_{p,q+1},
    \qquad
    b_{p+1,q}
    =
    b_{p,q+1}\delta^0.
\]
Naturality of \(\mu^\bullet\) therefore gives
\[
    d^iF_{p,q}(x,y)
    =
    F_{p+1,q}(d_C^ix,y)
    \qquad(0\leq i\leq p),
\]
\[
    d^{p+j}F_{p,q}(x,y)
    =
    F_{p,q+1}(x,d_D^jy)
    \qquad(1\leq j\leq q+1),
\]
and the break identity
\[
    F_{p+1,q}(d_C^{p+1}x,y)
    =
    F_{p,q+1}(x,d_D^0y).
\]

Expanding the alternating coface sums yields
\[
\begin{aligned}
    d_EF_{p,q}(x,y)
    &=
    \sum_{i=0}^{p}
    (-1)^i
    F_{p+1,q}(d_C^ix,y)
    \\
    &\qquad+
    \sum_{j=1}^{q+1}
    (-1)^{p+j}
    F_{p,q+1}(x,d_D^jy),
\end{aligned}
\]
whereas
\[
\begin{aligned}
    F_{p+1,q}(d_Cx,y)
    &+
    (-1)^pF_{p,q+1}(x,d_Dy)
    \\
    &=
    \sum_{i=0}^{p+1}
    (-1)^i
    F_{p+1,q}(d_C^ix,y)
    \\
    &\qquad+
    \sum_{j=0}^{q+1}
    (-1)^{p+j}
    F_{p,q+1}(x,d_D^jy).
\end{aligned}
\]
The only two additional terms in the second display are
\[
    (-1)^{p+1}
    F_{p+1,q}(d_C^{p+1}x,y)
\]
and
\[
    (-1)^p
    F_{p,q+1}(x,d_D^0y),
\]
and these cancel by the break identity.  Thus the two expressions agree in
\(E^{p+q+1}\).

Both expressions factor through \(N^{p+q+1}E\).  Stable Dold--Kan identifies
the normalized term as a direct summand of the corresponding cosimplicial
degree in the additive homotopy category, so
\[
    N^{p+q+1}E\longrightarrow E^{p+q+1}
\]
is split monic there.  The equality after inclusion is therefore precisely
\[
    d_E(x\boxtimes_{\mathrm{AW}}y)
    =
    d_Cx\boxtimes_{\mathrm{AW}}y
    +
    (-1)^p
    x\boxtimes_{\mathrm{AW}}d_Dy
\]
in the homotopy category.

Naturality follows from functoriality of the matching-cube construction, the
pairing, the consecutive-block maps, and stable Dold--Kan normalization.
\end{proof}

\begin{theorem}[Alexander--Whitney product on normalized infinitesimal cochains]
\label{thm:alexander-whitney-normalization}
For \(p,q\geq0\), there are canonical maps
\[
\boxed{
    \smile_{\mathrm{AW}}:
    \underline{\OmegaInf}^{\,p}(X/S)
    \otimes
    \underline{\OmegaInf}^{\,q}(X/S)
    \longrightarrow
    \underline{\OmegaInf}^{\,p+q}(X/S).
}
\]
Their images in the homotopy category are associative and unital and are represented by
\[
\boxed{
    x\smile_{\mathrm{AW}}y
    =
    \left(
        d^{p+q}\cdots d^{p+1}x
    \right)
    \cdot
    \left(
        (d^0)^py
    \right).
}
\]
The coherent infinitesimal differential has the homotopy-category Leibniz identity
\[
\boxed{
    d_{\inf}(x\smile_{\mathrm{AW}}y)
    =
    d_{\inf}x\smile_{\mathrm{AW}}y
    +
    (-1)^p
    x\smile_{\mathrm{AW}}d_{\inf}y.
}
\]
No \(E_1\)-algebra structure on the coherent normalized infinitesimal cochain object, no multiplicative totalization
filtration, and no symmetric-shuffle commutative algebra structure on un-totalized normalized infinitesimal cochains
is asserted.
\end{theorem}

\begin{proof}
Apply Proposition~\ref{prop:ordered-monoidal-stable-dold-kan} to
\[
    C^\bullet
    =
    \underline{\mathbf C}_{\dR}^{\,\bullet}(X/S).
\]
Its normalized terms are \(\underline{\OmegaInf}^{\,q}(X/S)\), and the differential supplied by the
complete \v{C}ech-totalization filtration is the coherent infinitesimal differential of
Theorem~\ref{thm:monadic-de-rham-dold-kan}.  The formulas and the homotopy-category algebra and
Leibniz identities therefore specialize exactly as stated.
\end{proof}

\begin{construction}[Normalized cup product]
\label{constr:normalized-modal-cup-product}
The maps of Theorem~\ref{thm:alexander-whitney-normalization} are called the normalized
Alexander--Whitney cup product.  Whenever multiplicative statements about normalized infinitesimal cochains are made
below, they refer to these natural maps and, for associativity, unitality, and Leibniz identities, to
their images in the homotopy category.
\end{construction}

\begin{theorem}[Commutative algebra structures on total de Rham objects]
\label{thm:multiplicative-modal-de-rham}
Every \(\mathbf{DR}_{X/S}^{<q}\) and \(\mathbf{DR}_{X/S}\) carries its canonical commutative algebra
structure as a partial or complete limit of
\(\underline{\mathbf C}_{\dR}^{\,\bullet}(X/S)\).  The transition maps form a tower in
\(\CAlg(\mathcal E_S)\), and
\[
\boxed{
    \mathbf{DR}_{X/S}
    \xrightarrow{\simeq}
    \limop_q\mathbf{DR}_{X/S}^{<q}
}
\]
is an equivalence in \(\CAlg(\mathcal E_S)\).  The equivalences of
Theorem~\ref{thm:effective-quotient-de-rham} are equivalences of commutative algebras.
\end{theorem}

\begin{proof}
Every partial or complete totalization is a limit of commutative algebra objects, and limits in
\(\CAlg(\mathcal E_S)\) are created in \(\mathcal E_S\).  The completion equivalence is therefore
the commutative-algebra refinement of Theorem~\ref{thm:modal-hodge-completeness}.  The effective-
descent and internal-function equivalences of Theorem~\ref{thm:effective-quotient-de-rham} are
constructed from strong symmetric monoidal pullback, lax symmetric monoidal dependent product,
limits of commutative algebras, and convolution against the cocommutative diagonal, so they are
equivalences in \(\CAlg(\mathcal E_S)\).
\end{proof}

\begin{proposition}[Totalization-stage quotient]
\label{prop:modal-hodge-multiplicativity}
For every \(q>0\), the defining totalization fibre sequence gives
\[
\boxed{
    \cofib
    \left(
        F_{\mathrm{Tot}}^q\mathbf{DR}_{X/S}
        \longrightarrow
        \mathbf{DR}_{X/S}
    \right)
    \simeq
    \mathbf{DR}_{X/S}^{<q}.
}
\]
The right side carries its canonical \(E_\infty\)-algebra structure as a partial totalization.  No
ideal structure on \(F_{\mathrm{Tot}}^q\mathbf{DR}_{X/S}\) and no cofiber construction in
\(\CAlg(\mathcal E_S)\) is asserted.
\end{proposition}

\begin{proof}
This is the defining fibre sequence of \(F_{\mathrm{Tot}}^q\) in the stable category, in which fibres and
cofibers agree up to the canonical stable identification.  The commutative algebra structure on the
partial totalization is Theorem~\ref{thm:multiplicative-modal-de-rham}.
\end{proof}

\begin{proposition}[Partial totalization and the canonical limit product]
\label{prop:lax-monoidal-coherent-totalization}
The tower
\[
    \cdots\longrightarrow
    \mathbf{DR}_{X/S}^{<q+1}
    \longrightarrow
    \mathbf{DR}_{X/S}^{<q}
    \longrightarrow\cdots
\]
is a tower in \(\CAlg(\mathcal E_S)\), and
\[
    \mathbf{DR}_{X/S}
    \simeq
    \limop_q\mathbf{DR}_{X/S}^{<q}
\]
there.
\end{proposition}

\begin{proof}
This is Theorem~\ref{thm:multiplicative-modal-de-rham}.
\end{proof}


\section{The skeletal presentation of the \v{C}ech-totalization filtration}
\label{sec:skeletal-hodge-filtration}

\begin{construction}[Skeletal quotients of the infinitesimal nerve]
\label{constr:modal-skeletal-quotients}
Let
\[
    A_\bullet=A_\bullet(X/S)
    =
    \Sigma^\infty_{+,S}R_{X/S,\bullet}.
\]
For \(q\geq0\), define
\[
\boxed{
    A_{\geq q}(X/S)
    :=
    \cofib
    \left(
        |\sk_{q-1}A_\bullet|
        \longrightarrow
        |A_\bullet|
    \right),
}
\]
where \(\sk_{-1}A_\bullet=0\).

In pointed objects of \(\mathcal X_{\Sp,/S}\), set
\[
    |\sk_{-1}R_{X/S,\bullet}|_+:=*,
\]
and define
\[
\boxed{
    \mathfrak I^{\langle q\rangle}_{X/S}
    :=
    \cofib_{\mathcal X_{\Sp,/S,*}}
    \left(
        |\sk_{q-1}R_{X/S,\bullet}|_+
        \longrightarrow
        Q_{X/S,+}
    \right).
}
\]
Then
\[
\boxed{
    A_{\geq q}(X/S)
    \simeq
    \Sigma^\infty_S
    \mathfrak I^{\langle q\rangle}_{X/S}.
}
\]
\end{construction}

\begin{theorem}[Skeletal totalization presentation]
\label{thm:skeletal-hodge-presentation}
For every \(q\geq0\),
\[
\boxed{
    F_{\mathrm{Tot}}^q\mathbf{DR}_{X/S}
    \simeq
    \underline{\Hom}_{\mathcal E_S}
    \left(
        A_{\geq q}(X/S),
        \mathbb O_S^{\add}
    \right).
}
\]
Equivalently,
\[
\boxed{
    \mathbf{DR}_{X/S}^{<q}
    \simeq
    \underline{\Hom}_{\mathcal E_S}
    \left(
        |\sk_{q-1}A_\bullet(X/S)|,
        \mathbb O_S^{\add}
    \right).
}
\]
The equivalences are natural in \(X\), compatible with arbitrary parameterized base change, and
identify the totalization transition \(F_{\mathrm{Tot}}^{q+1}\to F_{\mathrm{Tot}}^q\) with the map induced contravariantly by the
canonical morphism \(A_{\geq q}\to A_{\geq q+1}\).
\end{theorem}

\begin{proof}
By Theorem~\ref{thm:modal-cochains-internal-functions},
\[
    \underline{\mathbf C}_{\dR}^{\,\bullet}(X/S)
    \simeq
    \underline{\Hom}_{\mathcal E_S}
    \left(
        A_\bullet(X/S),\mathbb O_S^{\add}
    \right).
\]
Realization of the \((q-1)\)-skeleton is the colimit of the simplicial diagram restricted to
degrees at most \(q-1\).  Since internal Hom converts colimits in its first variable to limits,
\[
\begin{aligned}
    \mathbf{DR}_{X/S}^{<q}
    &=
    \Tot_{\Delta_{\leq q-1}}
    \underline{\mathbf C}_{\dR}^{\,\bullet}(X/S)
    \\
    &\simeq
    \underline{\Hom}_{\mathcal E_S}
    \left(
        |\sk_{q-1}A_\bullet(X/S)|,
        \mathbb O_S^{\add}
    \right).
\end{aligned}
\]
Likewise,
\[
    \mathbf{DR}_{X/S}
    \simeq
    \underline{\Hom}_{\mathcal E_S}
    \left(
        |A_\bullet(X/S)|,
        \mathbb O_S^{\add}
    \right).
\]
Applying internal Hom to the cofiber sequence
\[
    |\sk_{q-1}A_\bullet|
    \longrightarrow
    |A_\bullet|
    \longrightarrow
    A_{\geq q}
\]
gives a fibre sequence
\[
    \underline{\Hom}(A_{\geq q},\mathbb O_S^{\add})
    \longrightarrow
    \mathbf{DR}_{X/S}
    \longrightarrow
    \mathbf{DR}_{X/S}^{<q}.
\]
The fibre is \(F_{\mathrm{Tot}}^q\mathbf{DR}_{X/S}\) by definition.  Naturality and base change follow from
functoriality of skeleta, realization, cofibers, and closed base change.
\end{proof}

\begin{remark}[Intrinsic totalization associated graded]
The intrinsic totalization associated graded is
\[
    \gr_{\mathrm{Tot}}^q\mathbf{DR}_{X/S}
    \simeq
    \underline{\OmegaInf}^{\,q}(X/S)[-q].
\]
By Proposition~\ref{prop:normalized-one-cochains-strictly-nonlinear}, even the degree-one
normalized infinitesimal cochain can contain nonzero higher infinitesimal classes killed by its
canonical cotangent linearization.  Thus the totalization associated graded is not identified with a
Hodge associated graded, and no exterior-cotangent identification is asserted.  Any comparison with
exterior powers of the cotangent complex requires additional linearization or HKR-type input.
\end{remark}


\section{Coherently closed normalized cochains, exact cochains, and primitives}
\label{sec:closed-forms-deloopings}

\begin{definition}[Coherently closed normalized infinitesimal cochains]
\label{def:coherently-closed-form-spectrum}
For \(q\geq0\) and \(n\in\mathbb Z\), define the parameterized spectrum of coherently closed shifted normalized infinitesimal
\(q\)-cochains by
\[
\boxed{
    \underline{\OmegaInf}^{\,q,\cl}(X/S;n)
    :=
    F_{\mathrm{Tot}}^q\mathbf{DR}_{X/S}[n+q].
}
\]
Define its global section spectrum by
\[
\boxed{
    \OmegaInf^{q,\cl}(X/S;n)
    :=
    \map_{\mathcal E_S}
    \left(
        \mathbf 1_S,
        \underline{\OmegaInf}^{\,q,\cl}(X/S;n)
    \right).
}
\]
Write
\[
    \underline{\OmegaInf}^{\,q,\cl}(X/S)
    :=
    \underline{\OmegaInf}^{\,q,\cl}(X/S;0),
    \qquad
    \OmegaInf^{q,\cl}(X/S)
    :=
    \OmegaInf^{q,\cl}(X/S;0).
\]
The associated-graded map
\[
    F_{\mathrm{Tot}}^q\mathbf{DR}_{X/S}
    \longrightarrow
    \gr_{\mathrm{Tot}}^q\mathbf{DR}_{X/S}
    \simeq
    \underline{\OmegaInf}^{\,q}(X/S)[-q]
\]
induces the underlying-cochain map
\[
\boxed{
    u_q:
    \underline{\OmegaInf}^{\,q,\cl}(X/S;n)
    \longrightarrow
    \underline{\OmegaInf}^{\,q}(X/S)[n].
}
\]
\end{definition}

\begin{proposition}[Coherent closedness is retained structure]
\label{prop:closedness-coherent-structure}
A point of
\[
    \Omega^\infty\OmegaInf^{q,\cl}(X/S;n)
\]
consists of a shifted normalized infinitesimal \(q\)-cochain together with a specified nullhomotopy of its infinitesimal boundary and
the full hierarchy of higher compatibility data determined by the complete \v{C}ech-totalization filtration.  For a
chosen underlying cochain \(\omega\), the homotopy fibre of the global underlying-cochain map over
\(\omega\) is the space of coherent closedness structures on \(\omega\).
\end{proposition}

\begin{proof}
By Theorem~\ref{thm:monadic-de-rham-dold-kan}, the exact tower
\(F_{\mathrm{Tot}}^\bullet\mathbf{DR}_{X/S}\) has successive shifted cofibers
\(\underline{\OmegaInf}^{\,q}(X/S)\), and its successive connecting morphisms are the coherent infinitesimal
boundaries.  A point of the shifted stage \(F_{\mathrm{Tot}}^q[q]\) therefore determines a degree-\(q\) normalized infinitesimal cochain
together with compatible lifts through all subsequent stages of the tower.  The first compatibility
is a nullhomotopy of the unary boundary, the next is compatibility with the specified nullhomotopy
of \(d_{\inf}^2\), and the remaining lifts are precisely the higher compatibility data retained by
the complete tower.
\end{proof}

\begin{theorem}[Closed normalized cochains as shifted totalization stages]
\label{thm:closed-forms-hodge-tail}
There are canonical equivalences
\[
\boxed{
    \underline{\OmegaInf}^{\,q,\cl}(X/S;n)
    \simeq
    F_{\mathrm{Tot}}^q\mathbf{DR}_{X/S}[n+q]
}
\]
and
\[
\boxed{
    \OmegaInf^{q,\cl}(X/S;n)
    \simeq
    \Gamma_S^{\mathrm{st}}
    \left(
        F_{\mathrm{Tot}}^q\mathbf{DR}_{X/S}[n+q]
    \right).
}
\]
They commute with arbitrary parameterized base change.
\end{theorem}

\begin{proof}
The first equivalence is the definition.  The second is stable sections.  Base-change compatibility
is Proposition~\ref{prop:normalized-forms-base-change}.
\end{proof}

\begin{construction}[Exact-cochain classifier]
\label{constr:exact-form-classifier}
For \(q\geq1\), the associated-graded cofiber sequence is
\[
    F_{\mathrm{Tot}}^q\mathbf{DR}_{X/S}
    \longrightarrow
    F_{\mathrm{Tot}}^{q-1}\mathbf{DR}_{X/S}
    \longrightarrow
    \underline{\OmegaInf}^{\,q-1}(X/S)[1-q].
\]
Its connecting morphism, shifted by \(q-1\), defines
\[
\boxed{
    d_{\inf}^{\cl}:
    \underline{\OmegaInf}^{\,q-1}(X/S)
    \longrightarrow
    F_{\mathrm{Tot}}^q\mathbf{DR}_{X/S}[q]
    =
    \underline{\OmegaInf}^{\,q,\cl}(X/S;0).
}
\]
After shifting coefficients by \(n\) and taking stable sections, we obtain
\[
\boxed{
    d_{\inf}^{\cl}:
    \OmegaInf^{q-1}(X/S;n)
    \longrightarrow
    \OmegaInf^{q,\cl}(X/S;n).
}
\]
\end{construction}

\begin{proposition}[Underlying boundary of the exact-cochain classifier]
\label{prop:exact-form-underlying-boundary}
The composite
\[
    \underline{\OmegaInf}^{\,q-1}(X/S)
    \xrightarrow{d_{\inf}^{\cl}}
    \underline{\OmegaInf}^{\,q,\cl}(X/S;0)
    \xrightarrow{u_q}
    \underline{\OmegaInf}^{\,q}(X/S)
\]
is the coherent infinitesimal boundary \(d_{\inf}\).  The closedness structure carried by
\(d_{\inf}^{\cl}\) is the canonical coherent nullhomotopy of \(d_{\inf}^2\) together with all higher
compatibilities.
\end{proposition}

\begin{proof}
The coherent differential associated with a complete filtration is obtained from the connecting maps
of two consecutive associated-graded cofiber sequences.  The displayed composite is exactly this
construction.  The fact that the connecting map lands in the entire stage
\(F_{\mathrm{Tot}}^q[q]\), rather than merely its first associated-graded term, retains the complete higher
coherent closedness datum.
\end{proof}

\begin{definition}[Primitive spaces and primitive differences]
\label{def:modal-form-primitive-spectrum}
Fix \(q\geq1\) and
\[
    \omega\in
    \Omega^\infty\OmegaInf^{q,\cl}(X/S;n).
\]
Define
\[
\boxed{
    \operatorname{PrimSpace}(\omega)
    :=
    \hofib_{\omega}
    \left(
        \Omega^\infty d_{\inf}^{\cl}:
        \Omega^\infty\OmegaInf^{q-1}(X/S;n)
        \longrightarrow
        \Omega^\infty\OmegaInf^{q,\cl}(X/S;n)
    \right)
}
\]
and
\[
\boxed{
    \mathbf{PrimDiff}_{\inf}^q(X/S;n)
    :=
    \fib
    \left(
        d_{\inf}^{\cl}:
        \OmegaInf^{q-1}(X/S;n)
        \longrightarrow
        \OmegaInf^{q,\cl}(X/S;n)
    \right).
}
\]
The point \(\omega\) includes its coherent closedness structure.
\end{definition}

\begin{theorem}[Primitive torsor theorem]
\label{thm:modal-form-primitive-torsor}
The space \(\operatorname{PrimSpace}(\omega)\) is empty precisely when the homotopy fibre of
\(\Omega^\infty d_{\inf}^{\cl}\) over \(\omega\) is empty, equivalently when the component of
\(\omega\) does not lie in the image on \(\pi_0\).  When nonempty, it is canonically a torsor for
\[
    \Omega^\infty\mathbf{PrimDiff}_{\inf}^q(X/S;n).
\]
A chosen primitive \(\lambda\) gives an equivalence
\[
    \operatorname{PrimSpace}(\omega)
    \xrightarrow{\simeq}
    \Omega^\infty\mathbf{PrimDiff}_{\inf}^q(X/S;n)
\]
by subtraction of \(\lambda\).  For the zero coherently closed cochain, the zero primitive makes this
equivalence canonical.
\end{theorem}

\begin{proof}
The map \(d_{\inf}^{\cl}\) is a morphism of spectra, hence its infinite-loop map is a morphism of
grouplike \(E_\infty\)-spaces.  Translation by the fibre over zero acts on every nonempty homotopy
fibre.  Choosing one point identifies that fibre with the fibre over zero by subtraction.
\end{proof}

\begin{theorem}[Additive-delooping classifier]
\label{thm:additive-delooping-classifier}
For every \(q\geq0\) and \(n\in\mathbb Z\),
\[
\boxed{
    \OmegaInf^{q,\cl}(X/S;n)
    \simeq
    \map_{\mathcal E_S}
    \left(
        A_{\geq q}(X/S),
        \mathbb O_S^{\add}[n+q]
    \right).
}
\]
For \(n=0\),
\[
\boxed{
    \Omega^\infty\OmegaInf^{q,\cl}(X/S)
    \simeq
    \Map_{\mathcal X_{\Sp,/S,*}}
    \left(
        \mathfrak I^{\langle q\rangle}_{X/S},
        \Omega^\infty\Sigma^q\mathbb O_S^{\add}
    \right).
}
\]
For \(q=0\), this reduces canonically to
\[
\boxed{
    \Omega^\infty\OmegaInf^{0,\cl}(X/S)
    \simeq
    \Map_{\mathcal X_{\Sp,/S}}
    \left(
        Q_{X/S},
        \mathbb A^1_S
    \right).
}
\]
For \(q>0\), the collapse of the \((q-1)\)-skeleton is essential.
\end{theorem}

\begin{proof}
By Theorem~\ref{thm:skeletal-hodge-presentation},
\[
    F_{\mathrm{Tot}}^q\mathbf{DR}_{X/S}
    \simeq
    \underline{\Hom}_{\mathcal E_S}
    \left(
        A_{\geq q}(X/S),\mathbb O_S^{\add}
    \right).
\]
Shift by \(n+q\) and take stable sections to obtain the first formula.  Since
\[
    A_{\geq q}(X/S)
    \simeq
    \Sigma^\infty_S\mathfrak I^{\langle q\rangle}_{X/S},
\]
the fibrewise suspension-spectrum/infinite-loop adjunction gives the pointed mapping-space formula.
For \(q=0\), the \((-1)\)-skeleton is initial, so
\[
    \mathfrak I^{\langle0\rangle}_{X/S}\simeq Q_{X/S,+},
\]
and
\[
    \Omega^\infty\mathbb O_S^{\add}\simeq\mathbb A^1_S.
\]
Removing the disjoint basepoint gives the unpointed mapping-space formula.
\end{proof}

\begin{remark}[Why the quotient alone does not classify positive closed cochains]
For \(X=S\), the infinitesimal relation nerve is constant and every positive normalized infinitesimal cochain
spectrum vanishes.  By contrast,
\[
    \Map_S
    \left(
        S,\Omega^\infty\Sigma^q\mathbb O_S^{\add}
    \right)
\]
need not be contractible.  Thus positive closed cochains cannot be classified by unrestricted maps from
\(Q_{X/S}\).  The skeletal quotient
\(\mathfrak I^{\langle q\rangle}_{X/S}\) is exactly the normalization which removes lower-degree and
degenerate data.
\end{remark}


\section{Functoriality, base change, descent, and external products}
\label{sec:modal-de-rham-functoriality}

\begin{theorem}[Contravariant functoriality]
\label{thm:modal-de-rham-contravariant-functoriality}
A morphism \(f:X\to Y\) over \(S\) induces canonical morphisms
\[
    f_{\dR}^*:
    \underline{\mathbf C}_{\dR}^{\,\bullet}(Y/S)
    \longrightarrow
    \underline{\mathbf C}_{\dR}^{\,\bullet}(X/S),
\]
\[
    f_{\dR}^*:
    \underline{\OmegaInf}^{\,q}(Y/S)
    \longrightarrow
    \underline{\OmegaInf}^{\,q}(X/S),
\]
and
\[
    f_{\dR}^*:
    \mathbf{DR}_{Y/S}
    \longrightarrow
    \mathbf{DR}_{X/S}.
\]
They preserve the cosimplicial commutative algebra structure, the \v{C}ech-totalization tower, the coherent
\(d_{\inf}\), the natural Alexander--Whitney maps and their homotopy-category products, coherently
closed normalized cochains, exact-cochain classifiers, identities, composition, and all higher functoriality
coherences.
\end{theorem}

\begin{proof}
Proposition~\ref{prop:modal-nerve-functoriality} gives a simplicial map
\[
    R_{X/S,\bullet}\longrightarrow R_{Y/S,\bullet}.
\]
Pullback of scalar functions along its degreewise maps gives the contravariant cosimplicial map
\[
    \underline{\mathbf C}_{\dR}^{\,\bullet}(Y/S)
    \longrightarrow
    \underline{\mathbf C}_{\dR}^{\,\bullet}(X/S).
\]
Partial and complete totalization, matching fibres, associated-graded connecting maps, and the
Alexander--Whitney transformation are all functorial.  Hence the induced maps preserve every
subsequent structure.
\end{proof}

\begin{theorem}[Parameterized base change for the de Rham package]
\label{thm:base-change-modal-de-rham-package}
Let \(c:S'\to S\) and \(X'=X\times_SS'\).  There are canonical coherent equivalences
\[
    c^*Q_{X/S}\simeq Q_{X'/S'},
    \qquad
    c^*R_{X/S,\bullet}\simeq R_{X'/S',\bullet},
\]
\[
    c^*\underline{\mathbf C}_{\dR}^{\,\bullet}(X/S)
    \simeq
    \underline{\mathbf C}_{\dR}^{\,\bullet}(X'/S'),
\]
\[
    c^*\underline{\OmegaInf}^{\,q}(X/S)
    \simeq
    \underline{\OmegaInf}^{\,q}(X'/S'),
\]
\[
\boxed{
    c^*\mathbf{DR}_{X/S}
    \xrightarrow{\simeq}
    \mathbf{DR}_{X'/S'},
}
\]
and
\[
    c^*F_{\mathrm{Tot}}^q\mathbf{DR}_{X/S}
    \simeq
    F_{\mathrm{Tot}}^q\mathbf{DR}_{X'/S'}.
\]
They preserve the scalar algebra, partial totalizations, \(E_\infty\)-multiplication on total
objects, the natural Alexander--Whitney maps and their homotopy-category products on normalized
cochains, coherent boundaries, skeletal presentations, closed normalized cochains, and exact-cochain classifiers.  All equivalences satisfy identity,
composition, and higher pullback-pasting coherences.
\end{theorem}

\begin{proof}
The quotient and relation statements are
Proposition~\ref{prop:monadic-de-rham-quotient-base-change} and
Theorem~\ref{thm:modal-nerve-base-change}.  The scalar statement is
Proposition~\ref{prop:additive-scalar-spectrum-base-change}.  Right Beck--Chevalley gives the
degreewise cochain equivalence.  Since parameterized pullback preserves all limits and colimits, it
commutes with matching fibres, partial and complete totalizations, totalization fibres, skeletal cofibers,
and shifts.  Strong monoidality preserves the commutative-algebra structures on totalizations.  Naturality of
connecting morphisms and of the Alexander--Whitney maps gives the differential and normalized
cup-product compatibilities.
\end{proof}

\begin{construction}[Restriction on global normalized cochains]
\label{constr:global-form-base-change-restriction}
For \(c:S'\to S\), pullback of maps from the sphere and
Theorem~\ref{thm:base-change-modal-de-rham-package} define canonical restriction maps
\[
    \operatorname{res}_c:
    \OmegaInf^q(X/S;n)
    \longrightarrow
    \OmegaInf^q(X'/S';n),
\]
\[
    \operatorname{res}_c^{\cl}:
    \OmegaInf^{q,\cl}(X/S;n)
    \longrightarrow
    \OmegaInf^{q,\cl}(X'/S';n),
\]
and
\[
    \operatorname{res}_c^{\DR}:
    \mathbf R\Gamma_{\dR}(X/S)
    \longrightarrow
    \mathbf R\Gamma_{\dR}(X'/S').
\]
They commute with the available products, boundaries, exact-cochain classifiers, and totalization transitions.  They are not
asserted to be equivalences for arbitrary \(c\).
\end{construction}

\begin{proposition}[Finite-coproduct compatibility]
\label{prop:monadic-de-rham-finite-coproducts}
Fix \(S\in\mathcal X_{\Sp}\), let \(I\) be a finite set, and let
\(\{X_i\to S\}_{i\in I}\) be a finite family in \(\mathcal X_{\Sp,/S}\).  There are canonical
equivalences
\[
\boxed{
    \left(\coprod_{i\in I}X_i\right)_{\dR/S}
    \simeq
    \coprod_{i\in I}(X_i)_{\dR/S},
    \qquad
    Q_{(\coprod_iX_i)/S}
    \simeq
    \coprod_{i\in I}Q_{X_i/S}.
}
\]
For every \(n\geq0\), there is a canonical equivalence
\[
\boxed{
    R_{(\coprod_iX_i)/S,n}
    \simeq
    \coprod_{i\in I}R_{X_i/S,n}.
}
\]
Consequently the de Rham package carries finite coproducts in the geometric variable to finite
products of coefficient objects:
\[
\boxed{
    \underline{\mathbf C}_{\dR}^{\,\bullet}
    \left(\coprod_iX_i/S\right)
    \simeq
    \prod_i\underline{\mathbf C}_{\dR}^{\,\bullet}(X_i/S),
}
\]
\[
\boxed{
    \underline{\OmegaInf}^{\,q}
    \left(\coprod_iX_i/S\right)
    \simeq
    \prod_i\underline{\OmegaInf}^{\,q}(X_i/S),
}
\]
\[
\boxed{
    \mathbf{DR}^{\mathrm{full}}_{(\coprod_iX_i)/S}
    \simeq
    \prod_i\mathbf{DR}^{\mathrm{full}}_{X_i/S},
    \qquad
    \mathbf{DR}^{\mathrm{eff}}_{(\coprod_iX_i)/S}
    \simeq
    \prod_i\mathbf{DR}^{\mathrm{eff}}_{X_i/S},
}
\]
and, for every \(q\geq0\) and \(n\in\mathbb Z\),
\[
\boxed{
    F_{\mathrm{Tot}}^q\mathbf{DR}_{(\coprod_iX_i)/S}
    \simeq
    \prod_iF_{\mathrm{Tot}}^q\mathbf{DR}_{X_i/S},
    \qquad
    \underline{\OmegaInf}^{\,q,\cl}
    \left(\coprod_iX_i/S;n\right)
    \simeq
    \prod_i\underline{\OmegaInf}^{\,q,\cl}(X_i/S;n).
}
\]
The analogous equivalences hold for partial totalizations and, after applying stable sections, for
the corresponding global spectra.  All equivalences are natural in the finite family and compatible
with the structure maps already constructed.
\end{proposition}

\begin{proof}
We first show that finite ambient coproducts of relative de Rham-local objects are again relative de
Rham-local.  Let \(F_i\to S\) be local for every \(i\in I\).  For a finite spectral infinitesimal
thickening \(V\to T\), the pointwise lifting comparison for \(F_i\to S\) is
\[
    \Theta_{F_i/S}(V\to T):
    F_i(T)
    \longrightarrow
    S(T)\times_{S(V)}F_i(V).
\]
Its homotopy fibre over a point \(T\to S\) is precisely
\[
    \Map_S(T,F_i)
    \longrightarrow
    \Map_S(V,F_i),
\]
which is an equivalence by relative de Rham locality, composed along a finite square-zero
factorization of \(V\to T\).  Hence \(\Theta_{F_i/S}\) is an equivalence, so each
\(F_i\to S\) is formally \(\acute{e}\)tale with respect to finite spectral thickenings.

Put
\[
    F:=\coprod_{i\in I}F_i
\]
in the ambient slice.  In the infinitesimal adjunction of the preceding chapter, \(i_!\) preserves
coproducts because it is a left adjoint, while the source embedding \(i_*\) preserves all small
coproducts by the source-sheafification theorem and the fourth-adjoint construction.  Since colimits
in an \(\infty\)-topos are universal, pullback over \(i_*S\) distributes over coproducts.  Therefore
\[
\begin{aligned}
    i_!S\times_{i_*S}i_*F
    &\simeq
    i_!S\times_{i_*S}\coprod_i i_*F_i
    \\
    &\simeq
    \coprod_i\left(i_!S\times_{i_*S}i_*F_i\right),
\end{aligned}
\]
while
\[
    i_!F\simeq\coprod_i i_!F_i.
\]
Under these identifications, the comparison \(\Theta_{F/S}\) is the coproduct of the equivalences
\(\Theta_{F_i/S}\), and is therefore an equivalence.  Taking homotopy fibres of
\(\Theta_{F/S}(V\to T)\) over arbitrary points \(T\to S\) shows that
\[
    \Map_S(T,F)
    \xrightarrow{\simeq}
    \Map_S(V,F)
\]
for every affine spectral square-zero immersion \(V\to T\) over \(S\).  The slice-locality criterion
for \(L_{\dR,/S}\) shows that \(F\) is relative de Rham-local.

Apply this to the local objects \((X_i)_{\dR/S}\).  The coproduct of the localization units
\[
    \coprod_iX_i
    \longrightarrow
    \coprod_i(X_i)_{\dR/S}
\]
is a relative de Rham equivalence: for every local \(F\to S\), mapping into \(F\) gives the product
of the equivalences
\[
    \Map_S((X_i)_{\dR/S},F)
    \xrightarrow{\simeq}
    \Map_S(X_i,F).
\]
Its target is local by the preceding paragraph, so it exhibits the target as the relative de Rham
localization of \(\coprod_iX_i\).  This proves
\[
    \left(\coprod_iX_i\right)_{\dR/S}
    \simeq
    \coprod_i(X_i)_{\dR/S}.
\]

Now take the coproducts of the effective-image factorizations
\[
    X_i\twoheadrightarrow Q_{X_i/S}\hookrightarrow(X_i)_{\dR/S}.
\]
Finite coproducts of effective epimorphisms are effective epimorphisms and finite coproducts of
monomorphisms are monomorphisms in an \(\infty\)-topos.  Under the relative de Rham coproduct
equivalence just proved, their coproduct is therefore an effective-epimorphism--monomorphism
factorization of the relative de Rham unit of \(\coprod_iX_i\).  Uniqueness of image factorization
gives
\[
    Q_{(\coprod_iX_i)/S}
    \simeq
    \coprod_iQ_{X_i/S}.
\]
Using the effective \v{C}ech presentation and disjointness of coproducts,
\[
\begin{aligned}
    R_{(\coprod_iX_i)/S,n}
    &\simeq
    \underbrace{
        \left(\coprod_iX_i\right)
        \times_{\coprod_iQ_{X_i/S}}\cdots
        \times_{\coprod_iQ_{X_i/S}}
        \left(\coprod_iX_i\right)
    }_{n+1\text{ factors}}
    \\
    &\simeq
    \coprod_i
    \underbrace{
        X_i\times_{Q_{X_i/S}}\cdots\times_{Q_{X_i/S}}X_i
    }_{n+1\text{ factors}}
    \\
    &\simeq
    \coprod_iR_{X_i/S,n}.
\end{aligned}
\]

The relative suspension-spectrum functor preserves coproducts, while internal Hom converts
coproducts in its first variable to products.  The internal-function presentation therefore gives
\[
    \underline{\mathbf C}_{\dR}^{\,\bullet}
    \left(\coprod_iX_i/S\right)
    \simeq
    \prod_i\underline{\mathbf C}_{\dR}^{\,\bullet}(X_i/S).
\]
The internal-function presentations of the full and effective de Rham algebras give the corresponding
product equivalences for \(\mathbf{DR}^{\mathrm{full}}\) and \(\mathbf{DR}^{\mathrm{eff}}\).
Normalization, partial and complete totalization, the fibres defining the totalization stages, and
the coherently closed cochain construction are all formed from limits, stable fibres, and shifts.
Limits commute with finite products, so the remaining displayed equivalences follow.  Stable sections
preserve limits and therefore give the corresponding finite-product formulas for global spectra.
Naturality and compatibility with the existing structure maps follow from the functoriality of every
construction used above.
\end{proof}

\begin{lemma}[Relative de Rham equivalences are generated by square-zero immersions]
\label{lem:relative-dr-generated-by-square-zero}
Fix \(S\in\mathcal X_{\Sp}\).  Using Convention~\ref{conv:monadic-de-rham-universes}, choose a
set \(\Sigma_{\mathrm{sqz},S}\) of representatives of affine spectral square-zero immersions in the
slice \(\mathcal X_{\Sp,/S}\).  The local-equivalence
class of the relative de Rham localization
\[
    L_{\dR,/S}:\mathcal X_{\Sp,/S}\to\mathcal X_{\Sp,/S}^{\dR}
\]
is the strongly saturated class generated by \(\Sigma_{\mathrm{sqz},S}\).
\end{lemma}

\begin{proof}
Localize the presentable slice \(\mathcal X_{\Sp,/S}\) at the set
\(\Sigma_{\mathrm{sqz},S}\).  An object \(F\to S\) is local for this localization exactly when,
for every affine spectral square-zero immersion \(U\to T\) over \(S\), restriction induces an
equivalence
\[
    \Map_S(T,F)\xrightarrow{\simeq}\Map_S(U,F).
\]
The slice-locality theorem of the preceding chapter says that these are precisely the
\(L_{\dR,/S}\)-local objects.  Accessible reflective localizations with the same local objects are
canonically equivalent.  Their local-equivalence classes therefore agree, and the local-equivalence
class of a localization generated by a set is the strongly saturated class generated by that set.
\end{proof}

\begin{lemma}[Local-map criterion for left-exact localizations]
\label{lem:left-exact-local-map-criterion}
Let
\[
    L:\mathcal Y\longrightarrow\mathcal Y_L
\]
be a left-exact localization of an \(\infty\)-topos, and let \(\mathcal W\) be its class of local
equivalences.  For a morphism \(u:U\to X\) in \(\mathcal Y\), the following are equivalent:
\begin{enumerate}
\item \(u\) is right orthogonal to every morphism in \(\mathcal W\);
\item the naturality square
\[
\begin{tikzcd}[column sep=large,row sep=large]
    U \ar[r,"\eta_U"] \ar[d,"u"']
        & LU \ar[d,"Lu"]\\
    X \ar[r,"\eta_X"']
        & LX
\end{tikzcd}
\]
is Cartesian.
\end{enumerate}
\end{lemma}

\begin{proof}
If the square is Cartesian, then \(u\) is a pullback of the morphism \(Lu:LU\to LX\) between
local objects.  Every morphism between local objects is right orthogonal to local equivalences, and
right orthogonality is stable under pullback.  Thus (ii) implies (i).

Conversely, suppose (i) and put
\[
    P:=X\times_{LX}LU.
\]
Let
\[
    \alpha:U\longrightarrow P
\]
be the natural comparison and let \(p:P\to X\) be the projection.  Left exactness gives
\[
    LP
    \simeq
    LX\times_{LLX}LLU
    \simeq
    LU,
\]
so \(L\alpha\) is an equivalence and therefore \(\alpha\in\mathcal W\).  The map \(p\) is the
pullback of \(Lu\), hence is right orthogonal to \(\mathcal W\).

Apply \(\alpha\perp u\) to the square
\[
\begin{tikzcd}
    U \ar[r,"\id"] \ar[d,"\alpha"']
        & U \ar[d,"u"]\\
    P \ar[r,"p"']
        & X.
\end{tikzcd}
\]
There is a contractibly unique filler \(s:P\to U\) satisfying
\[
    s\alpha\simeq\id_U,
    \qquad
    us\simeq p.
\]
Now use \(\alpha\perp p\).  In the square
\[
\begin{tikzcd}
    U \ar[r,"\alpha"] \ar[d,"\alpha"']
        & P \ar[d,"p"]\\
    P \ar[r,"p"']
        & X
\end{tikzcd}
\]
both \(\id_P\) and \(\alpha s\) are fillers.  The filler space is contractible, so
\[
    \alpha s\simeq\id_P.
\]
Thus \(\alpha\) is an equivalence, which is exactly Cartesianity of the localization-unit square.
\end{proof}

\begin{proposition}[Formal-\(\acute{e}\)tale Cartesianity]
\label{prop:formal-etale-de-rham-cartesianity}
Let
\[
    u:U\longrightarrow X
\]
be a morphism in \(\mathcal X_{\Sp,/S}\) which is formally
\(\acute{e}\)tale relative to \(X\), in the sense that for every affine spectral
square-zero immersion
\[
    V\longrightarrow T
\]
over \(X\), restriction induces an equivalence
\[
    \Map_X(T,U)
    \xrightarrow{\simeq}
    \Map_X(V,U).
\]
Then the naturality square
\[
\begin{tikzcd}
    U \ar[r,"\eta_{U/S}"] \ar[d,"u"']
    & U_{\dR/S} \ar[d,"u_{\dR}"]
    \\
    X \ar[r,"\eta_{X/S}"']
    & X_{\dR/S}
\end{tikzcd}
\]
is Cartesian.  Consequently,
\[
\boxed{
    Q_{U/S}
    \simeq
    Q_{X/S}\times_{X_{\dR/S}}U_{\dR/S}
}
\]
and
\[
\boxed{
    U
    \simeq
    X\times_{Q_{X/S}}Q_{U/S}.
}
\]
\end{proposition}

\begin{proof}
Regard \(u:U\to X\) as an object of the slice
\(\mathcal X_{\Sp,/X}\).  By hypothesis, for every affine spectral
square-zero immersion
\[
    j:V\longrightarrow T
\]
over \(X\), restriction induces an equivalence
\[
    \Map_X(T,U)
    \xrightarrow{\simeq}
    \Map_X(V,U).
\]
Thus \(U\to X\) satisfies the slice-locality criterion for \(L_{\dR,/X}\).  Its relative de Rham
unit over \(X\) is therefore an equivalence:
\[
    U
    \xrightarrow{\simeq}
    U_{\dR/X}
    =
    X\times_{X_{\dR}}U_{\dR}.
\]

Using the definitions
\[
    X_{\dR/S}=S\times_{S_{\dR}}X_{\dR},
    \qquad
    U_{\dR/S}=S\times_{S_{\dR}}U_{\dR},
\]
finite-limit associativity gives a canonical equivalence
\[
    X\times_{X_{\dR/S}}U_{\dR/S}
    \simeq
    X\times_{X_{\dR}}U_{\dR}
    \simeq U.
\]
Hence the naturality square of relative de Rham units is Cartesian.

Pull the effective-image factorization
\[
    X\twoheadrightarrow Q_{X/S}\hookrightarrow X_{\dR/S}
\]
back along \(U_{\dR/S}\to X_{\dR/S}\).  Effective epimorphisms and monomorphisms in an
\(\infty\)-topos are stable under pullback.  By the Cartesian square just proved, the source of the
pulled-back factorization is \(U\), so it is an effective-epimorphism--monomorphism factorization of
\(U\to U_{\dR/S}\).  Uniqueness of image factorizations gives
\[
    Q_{U/S}
    \simeq
    Q_{X/S}\times_{X_{\dR/S}}U_{\dR/S}.
\]
Finally,
\[
\begin{aligned}
    X\times_{Q_{X/S}}Q_{U/S}
    &\simeq
    X\times_{Q_{X/S}}
    \left(
        Q_{X/S}\times_{X_{\dR/S}}U_{\dR/S}
    \right)
    \\
    &\simeq
    X\times_{X_{\dR/S}}U_{\dR/S}
    \\
    &\simeq U.
\end{aligned}
\]
\end{proof}

\begin{proposition}[Formal-\(\acute{e}\)tale invariance of intrinsic completion]
\label{prop:formal-etale-completion-invariance}
Let
\[
    Z\longrightarrow W\xrightarrow{u}Y\longrightarrow S
\]
be a diagram in \(\mathcal X_{\Sp}\) in which \(u\) is formally \(\acute{e}\)tale over \(S\).
Then the induced morphism
\[
\boxed{
    \widehat W^{\,\dR}_{Z/S}
    \longrightarrow
    \widehat Y^{\,\dR}_{Z/S}
}
\]
is an equivalence.  The equivalence is natural in the diagram and compatible with arbitrary base
change.
\end{proposition}

\begin{proof}
Proposition~\ref{prop:formal-etale-de-rham-cartesianity} gives a Cartesian square
\[
\begin{tikzcd}
    W \ar[r] \ar[d,"u"']
        & W_{\dR/S} \ar[d]
    \\
    Y \ar[r]
        & Y_{\dR/S}.
\end{tikzcd}
\]
Hence
\[
    W
    \simeq
    Y\times_{Y_{\dR/S}}W_{\dR/S}.
\]
Substituting this into the defining completion gives
\[
\begin{aligned}
    \widehat W^{\,\dR}_{Z/S}
    &=
    W\times_{W_{\dR/S}}Z_{\dR/S}
    \\
    &\simeq
    \left(
        Y\times_{Y_{\dR/S}}W_{\dR/S}
    \right)
    \times_{W_{\dR/S}}Z_{\dR/S}
    \\
    &\simeq
    Y\times_{Y_{\dR/S}}Z_{\dR/S}
    \\
    &=
    \widehat Y^{\,\dR}_{Z/S}.
\end{aligned}
\]
Naturality and base-change compatibility follow from those of formal-\(\acute{e}\)tale
Cartesianity and finite-limit pasting.
\end{proof}


\begin{theorem}[Formal-\(\acute{e}\)tale change of base for the intrinsic de Rham package]
\label{thm:formal-etale-change-of-base-intrinsic}
Let
\[
    g:S'\longrightarrow S
\]
be a native formally \(\acute{e}\)tale morphism, and let \(V\to S'\) be an
object of \(\mathcal X_{\Sp}\).  Regard \(V\) also as an object over \(S\) by
composition.  Then there are canonical equivalences
\[
\boxed{
    V_{\dR/S'}
    \xrightarrow{\simeq}
    V_{\dR/S},
}
\]
\[
\boxed{
    Q_{V/S'}
    \xrightarrow{\simeq}
    Q_{V/S},
    \qquad
    R_{V/S',\bullet}
    \xrightarrow{\simeq}
    R_{V/S,\bullet},
}
\]
where the right-hand objects carry their induced \(S'\)-structure.  For every
diagram
\[
    Z\longrightarrow Y\longrightarrow S'
\]
there is likewise a canonical equivalence
\[
\boxed{
    \widehat Y^{\,\dR}_{Z/S'}
    \xrightarrow{\simeq}
    \widehat Y^{\,\dR}_{Z/S}.
}
\]
All of these equivalences are compatible with identities, composition of
formally \(\acute{e}\)tale base morphisms, and arbitrary further base change.

The induced scalar-valued global de Rham invariants are canonically
equivalent.  In particular,
\[
\boxed{
    \mathbf R\Gamma_{\dR}^{\mathrm{full}}(V/S')
    \xrightarrow{\simeq}
    \mathbf R\Gamma_{\dR}^{\mathrm{full}}(V/S),
    \qquad
    \mathbf R\Gamma_{\dR}^{\mathrm{eff}}(V/S')
    \xrightarrow{\simeq}
    \mathbf R\Gamma_{\dR}^{\mathrm{eff}}(V/S),
}
\]
and for every \(q\geq0\) and \(n\in\mathbb Z\),
\[
\boxed{
    \OmegaInf^q(V/S';n)
    \xrightarrow{\simeq}
    \OmegaInf^q(V/S;n),
}
\]
\[
\boxed{
    \OmegaInf^{q,\cl}(V/S';n)
    \xrightarrow{\simeq}
    \OmegaInf^{q,\cl}(V/S;n),
}
\]
together with
\[
\boxed{
    \Gamma_{S'}^{\mathrm{st}}
    \left(
        F_{\mathrm{Tot}}^q\mathbf{DR}_{V/S'}
    \right)
    \xrightarrow{\simeq}
    \Gamma_S^{\mathrm{st}}
    \left(
        F_{\mathrm{Tot}}^q\mathbf{DR}_{V/S}
    \right).
}
\]
These equivalences preserve the full-to-effective comparison, partial
totalizations, the coherent infinitesimal differential, coherently closed
cochains, exact-cochain classifiers, and the natural Alexander--Whitney maps.
\end{theorem}

\begin{proof}
Apply Proposition~\ref{prop:formal-etale-de-rham-cartesianity} to
\(g:S'\to S\) over the terminal base.  The resulting Cartesian square gives
\[
\boxed{
    S'
    \xrightarrow{\simeq}
    S\times_{S_{\dR}}S'_{\dR}.
}
\]
Hence, by the definition of the relative de Rham object and associativity of
finite limits,
\[
\begin{aligned}
    V_{\dR/S'}
    &=
    S'\times_{S'_{\dR}}V_{\dR}
    \\
    &\simeq
    \left(
        S\times_{S_{\dR}}S'_{\dR}
    \right)
    \times_{S'_{\dR}}V_{\dR}
    \\
    &\simeq
    S\times_{S_{\dR}}V_{\dR}
    \\
    &=
    V_{\dR/S}.
\end{aligned}
\]
Under this equivalence the two relative de Rham units of \(V\) agree, because
both are induced by the same absolute unit \(V\to V_{\dR}\).  Effective
epimorphisms and monomorphisms in a slice are created in the ambient
\(\infty\)-topos.  Uniqueness of effective-image factorization therefore gives
\[
    Q_{V/S'}\simeq Q_{V/S}.
\]
Taking iterated fibre products over the identified relative de Rham target
gives
\[
    R_{V/S',n}\simeq R_{V/S,n}
\]
compatibly with every face and degeneracy map.  The completion formula follows
similarly from
\[
    Y_{\dR/S'}\simeq Y_{\dR/S},
    \qquad
    Z_{\dR/S'}\simeq Z_{\dR/S}.
\]

It remains to compare the scalar-valued invariants.  Let
\[
    R_n:=R_{V/S',n}\simeq R_{V/S,n},
\]
and write
\[
    r_n:R_n\to S',
    \qquad
    a_n:=g\circ r_n:R_n\to S.
\]
By Proposition~\ref{prop:additive-scalar-spectrum-base-change},
\[
    \mathbb O_{S'}^{\add}
    \simeq
    g^*\mathbb O_S^{\add},
\]
so
\[
    r_n^*\mathbb O_{S'}^{\add}
    \simeq
    a_n^*\mathbb O_S^{\add}.
\]
The mapping-spectrum interpretation of geometric de Rham cochains therefore
gives
\[
\begin{aligned}
    \Gamma_{S'}^{\mathrm{st}}
    \left(
        \underline{\mathbf C}_{\dR}^{\,n}(V/S')
    \right)
    &\simeq
    \map_{\mathcal E_{R_n}}
    \left(
        \mathbf1_{R_n},
        r_n^*\mathbb O_{S'}^{\add}
    \right)
    \\
    &\simeq
    \map_{\mathcal E_{R_n}}
    \left(
        \mathbf1_{R_n},
        a_n^*\mathbb O_S^{\add}
    \right)
    \\
    &\simeq
    \Gamma_S^{\mathrm{st}}
    \left(
        \underline{\mathbf C}_{\dR}^{\,n}(V/S)
    \right).
\end{aligned}
\]
These equivalences are compatible with all cosimplicial operators and with
multiplication.  Stable sections preserve limits and finite exact diagrams,
so the comparison propagates through matching fibres, partial and complete
totalizations, the fibres defining \(F_{\mathrm{Tot}}^q\), shifts, connecting
morphisms, and the exact-cochain classifier.  This proves the effective,
normalized, coherently closed, and stagewise assertions.

For the full de Rham invariant use the already identified object
\(V_{\dR/S'}\simeq V_{\dR/S}\) and the same scalar comparison on that common
object.  All identity, composition, and higher coherence statements are
inherited from the coherent formal-\(\acute{e}\)tale Cartesianity equivalence,
finite-limit associativity, scalar base change, and functoriality of the
subsequent constructions.
\end{proof}

\begin{theorem}[Formal-\(\acute{e}\)tale effective descent for effective infinitesimal quotients]
\label{thm:effective-quotient-formal-etale-descent}
Let
\[
    u:U\longrightarrow X
\]
be an effective epimorphism in \(\mathcal X_{\Sp,/S}\) which is formally \(\acute{e}\)tale over
\(S\), and let \(U_\bullet\to X\) be its \v{C}ech nerve.  Then
\[
\boxed{
    Q_{U/S}\longrightarrow Q_{X/S}
}
\]
is an effective epimorphism, and for every \(k\geq0\) there is a canonical equivalence
\[
\boxed{
    Q_{U_k/S}
    \simeq
    \underbrace{
        Q_{U/S}\times_{Q_{X/S}}\cdots\times_{Q_{X/S}}Q_{U/S}
    }_{k+1\text{ factors}}.
}
\]
Thus \(Q_{U_\bullet/S}\) is the \v{C}ech nerve of \(Q_{U/S}\to Q_{X/S}\) and
\[
    |Q_{U_\bullet/S}|\simeq Q_{X/S}.
\]
The construction is compatible with arbitrary base change.
\end{theorem}

\begin{proof}
Functoriality of effective images gives a commutative square
\[
\begin{tikzcd}
    U \ar[r,"\epsilon_{U/S}"] \ar[d,"u"']
    & Q_{U/S} \ar[d,"v"]
    \\
    X \ar[r,"\epsilon_{X/S}"']
    & Q_{X/S}.
\end{tikzcd}
\]
The composite
\[
    U\xrightarrow{u}X\xrightarrow{\epsilon_{X/S}}Q_{X/S}
\]
is an effective epimorphism, because both factors are.  It equals
\(v\epsilon_{U/S}\).  If a composite is an effective epimorphism in an \(\infty\)-topos, then its
second factor is an effective epimorphism: the image of the composite factors through the image of
the second factor.  Hence
\[
    v:Q_{U/S}\longrightarrow Q_{X/S}
\]
is an effective epimorphism.

Every \(U_k\to X\) is obtained from \(u\) by successive base changes and composition.  Formal
\(\acute{e}\)taleness is stable under composition: if \(A\to B\to C\) are formally
\(\acute{e}\)tale and \(V\to T\) is a finite spectral infinitesimal thickening, then
\[
\begin{aligned}
    A(T)
    &\simeq
    B(T)\times_{B(V)}A(V)
    \\
    &\simeq
    \left(C(T)\times_{C(V)}B(V)\right)\times_{B(V)}A(V)
    \\
    &\simeq
    C(T)\times_{C(V)}A(V),
\end{aligned}
\]
which is the formal-\(\acute{e}\)tale lifting equivalence for the composite.  Together with the
base-change stability from the preceding chapter, this shows that every \(U_k\to X\) is formally
\(\acute{e}\)tale.  Proposition~\ref{prop:formal-etale-de-rham-cartesianity} therefore gives
\[
    Q_{U_k/S}
    \simeq
    Q_{X/S}\times_{X_{\dR/S}}(U_k)_{\dR/S}.
\]
Left exactness of the relative de Rham localization gives
\[
    (U_k)_{\dR/S}
    \simeq
    \underbrace{
        U_{\dR/S}\times_{X_{\dR/S}}\cdots\times_{X_{\dR/S}}U_{\dR/S}
    }_{k+1}.
\]
Using again the formal-\(\acute{e}\)tale Cartesianity equivalence
\[
    Q_{U/S}
    \simeq
    Q_{X/S}\times_{X_{\dR/S}}U_{\dR/S},
\]
finite-limit pasting identifies the preceding expression with
\[
    Q_{U_k/S}
    \simeq
    \underbrace{
        Q_{U/S}\times_{Q_{X/S}}\cdots\times_{Q_{X/S}}Q_{U/S}
    }_{k+1}.
\]
Thus \(Q_{U_\bullet/S}\) is the \v{C}ech nerve of the effective epimorphism
\(v:Q_{U/S}\to Q_{X/S}\), so its realization is \(Q_{X/S}\).  Every construction used is compatible
with pullback, which gives the asserted arbitrary base-change compatibility.
\end{proof}

\begin{corollary}[Descent of effective infinitesimal quotients]
\label{thm:effective-quotient-nisnevich-descent}
Let \(X\to S\) be represented by a spectral scheme, let
\(\{U_\alpha\to X\}_{\alpha\in A}\) be a finite spectral Nisnevich covering family, put
\[
    U:=\coprod_{\alpha\in A}U_\alpha,
\]
and let \(U_\bullet\to X\) be the \v{C}ech nerve of the associated map \(u:U\to X\).  Then
\[
\boxed{
    Q_{U/S}\longrightarrow Q_{X/S}
}
\]
is an effective epimorphism,
\[
\boxed{
    Q_{U_k/S}
    \simeq
    \underbrace{
        Q_{U/S}\times_{Q_{X/S}}\cdots\times_{Q_{X/S}}Q_{U/S}
    }_{k+1\text{ factors}},
}
\]
and
\[
    |Q_{U_\bullet/S}|\simeq Q_{X/S}.
\]
\end{corollary}

\begin{proof}
The map \(u:U\to X\) is an effective epimorphism and a represented spectral
\(\acute{e}\)tale morphism.  Spectral \(\acute{e}\)tale morphisms are formally
\(\acute{e}\)tale by the preceding chapter, so
Theorem~\ref{thm:effective-quotient-formal-etale-descent} applies.
\end{proof}

\begin{proposition}[Descent of infinitesimal relation degrees]
\label{prop:modal-relation-degree-nisnevich-descent}
Under the hypotheses of Theorem~\ref{thm:effective-quotient-nisnevich-descent}, for every \(n\geq0\)
there are canonical equivalences
\[
\boxed{
    R_{U_k/S,n}
    \simeq
    R_{X/S,n}\times_{Q_{X/S}}Q_{U_k/S}.
}
\]
Consequently \(k\mapsto R_{U_k/S,n}\) is an effective \v{C}ech nerve with realization
\(R_{X/S,n}\).
\end{proposition}

\begin{proof}
Formal-\(\acute{e}\)tale Cartesianity gives
\[
    U_k\simeq X\times_{Q_{X/S}}Q_{U_k/S}.
\]
Using the effective \v{C}ech presentation of the relation nerve,
\[
\begin{aligned}
    R_{U_k/S,n}
    &=
    \underbrace{
        U_k\times_{Q_{U_k/S}}\cdots\times_{Q_{U_k/S}}U_k
    }_{n+1}
    \\
    &\simeq
    R_{X/S,n}\times_{Q_{X/S}}Q_{U_k/S}.
\end{aligned}
\]
The right side is the pullback of the \v{C}ech nerve of the effective epimorphism
\(Q_{U/S}\to Q_{X/S}\).
\end{proof}

\begin{theorem}[Spectral Nisnevich descent]
\label{thm:modal-de-rham-nisnevich-descent}
Under the same hypotheses, there are canonical equivalences
\[
\boxed{
    \underline{\mathbf C}_{\dR}^{\,\bullet}(X/S)
    \xrightarrow{\simeq}
    \Tot_k
    \underline{\mathbf C}_{\dR}^{\,\bullet}(U_k/S),
}
\]
\[
\boxed{
    \underline{\OmegaInf}^{\,q}(X/S)
    \xrightarrow{\simeq}
    \Tot_k
    \underline{\OmegaInf}^{\,q}(U_k/S),
}
\]
\[
\boxed{
    \mathbf{DR}_{X/S}
    \xrightarrow{\simeq}
    \Tot_k
    \mathbf{DR}_{U_k/S},
}
\]
and
\[
\boxed{
    F_{\mathrm{Tot}}^q\mathbf{DR}_{X/S}
    \xrightarrow{\simeq}
    \Tot_k
    F_{\mathrm{Tot}}^q\mathbf{DR}_{U_k/S}.
}
\]
The same holds for partial totalizations, coherently closed normalized cochains, exact-cochain classifiers, the
canonical commutative-algebra maps on totalizations, the natural normalized Alexander--Whitney maps,
and all totalization transition maps.
\end{theorem}

\begin{proof}
Fix \(n\).  By Proposition~\ref{prop:modal-relation-degree-nisnevich-descent}, the objects
\(R_{U_k/S,n}\) form an effective \v{C}ech nerve over \(R_{X/S,n}\).  Effective descent for
parameterized spectra gives
\[
    a_n^*\mathbb O_S^{\add}
    \xrightarrow{\simeq}
    \Tot_k\Pi_{v_k}v_k^*a_n^*\mathbb O_S^{\add}
\]
in \(\mathcal E_{R_{X/S,n}}\).  Applying \(\Pi_{a_n}\) and using composition of dependent products yields
\[
    \underline{\mathbf C}_{\dR}^{\,n}(X/S)
    \xrightarrow{\simeq}
    \Tot_k
    \underline{\mathbf C}_{\dR}^{\,n}(U_k/S).
\]
The construction is natural in \(n\), so it is an equivalence of cosimplicial commutative algebras.

All remaining constructions use limits, stable fibres or cofibres, and shifts.  Matching fibres use
finite limits; complete and partial totalizations use limits; totalization stages are fibres; closed normalized cochains
are shifted totalization stages.  Limits commute with limits, so the displayed degreewise descent
equivalence propagates through every construction.  The commutative-algebra structures on
totalizations, the natural Alexander--Whitney maps, and the connecting morphisms are natural and
therefore descend as well.
\end{proof}

\begin{theorem}[Nisnevich descent for full de Rham functions]
\label{thm:full-dr-nisnevich-descent}
Under the hypotheses of Theorem~\ref{thm:effective-quotient-nisnevich-descent}, there is a canonical
equivalence of commutative algebras
\[
\boxed{
    \mathbf{DR}^{\mathrm{full}}_{X/S}
    \xrightarrow{\simeq}
    \Tot_k
    \mathbf{DR}^{\mathrm{full}}_{U_k/S}.
}
\]
The comparison maps
\[
    \rho_{U_k/S}:
    \mathbf{DR}^{\mathrm{full}}_{U_k/S}
    \longrightarrow
    \mathbf{DR}^{\mathrm{eff}}_{U_k/S}
\]
assemble with descent to the comparison
\(\rho_{X/S}:\mathbf{DR}^{\mathrm{full}}_{X/S}\to\mathbf{DR}^{\mathrm{eff}}_{X/S}\).
\end{theorem}

\begin{proof}
Let \(u:U\to X\) be the effective epimorphism associated to the finite represented Nisnevich cover.
By Theorem~\ref{thm:relative-dr-preserves-effective-epis}, the induced morphism
\[
    u_{\dR/S}:
    U_{\dR/S}
    \longrightarrow
    X_{\dR/S}
\]
is an effective epimorphism in the ambient slice \(\mathcal X_{\Sp,/S}\).  Left exactness of the
relative de Rham localization identifies its \v{C}ech nerve with the relative de Rham objects of
the original \v{C}ech nerve:
\[
    (U_k)_{\dR/S}
    \simeq
    \underbrace{
        U_{\dR/S}\times_{X_{\dR/S}}\cdots\times_{X_{\dR/S}}U_{\dR/S}
    }_{k+1\text{ factors}}.
\]

Write
\[
    q_{\dR}:X_{\dR/S}\longrightarrow S,
    \qquad
    v_k:(U_k)_{\dR/S}\longrightarrow X_{\dR/S},
\]
and
\[
    q_{k,\dR}:=q_{\dR}v_k.
\]
Effective descent for parameterized spectra applied to
\[
    q_{\dR}^*\mathbb O_S^{\add}
    \in
    \mathcal E_{X_{\dR/S}}
\]
gives
\[
    q_{\dR}^*\mathbb O_S^{\add}
    \xrightarrow{\simeq}
    \Tot_k
    \Pi_{v_k}v_k^*q_{\dR}^*\mathbb O_S^{\add}.
\]
Applying \(\Pi_{q_{\dR}}\), which preserves limits, and using composition of dependent products,
we obtain
\[
\begin{aligned}
    \mathbf{DR}^{\mathrm{full}}_{X/S}
    &=
    \Pi_{q_{\dR}}q_{\dR}^*\mathbb O_S^{\add}
    \\
    &\simeq
    \Tot_k
    \Pi_{q_{\dR}}\Pi_{v_k}v_k^*q_{\dR}^*\mathbb O_S^{\add}
    \\
    &\simeq
    \Tot_k
    \Pi_{q_{k,\dR}}q_{k,\dR}^*\mathbb O_S^{\add}
    \\
    &=
    \Tot_k
    \mathbf{DR}^{\mathrm{full}}_{U_k/S}.
\end{aligned}
\]
The effective-descent equivalence is symmetric monoidal, pullback is strong symmetric monoidal,
dependent product is lax symmetric monoidal, and limits of commutative algebras are created in the
underlying category.  Hence the displayed equivalence is one of commutative algebras.
Compatibility with \(\rho\) follows from naturality of precomposition along
\(Q_{-/S}\to(-)_{\dR/S}\) and the full and effective descent equivalences.
\end{proof}

\begin{proposition}[External products]
\label{prop:modal-de-rham-external-products}
For \(X\to S\) and \(Y\to S\), there is a canonical morphism of commutative algebras
\[
\boxed{
    \mathbf{DR}_{X/S}
    \otimes
    \mathbf{DR}_{Y/S}
    \longrightarrow
    \mathbf{DR}_{X\times_SY/S}.
}
\]
For \(p,q\geq0\), the same cosimplicial pairing and ordered Alexander--Whitney normalization give
natural maps
\[
\boxed{
    \underline{\OmegaInf}^{\,p}(X/S)
    \otimes
    \underline{\OmegaInf}^{\,q}(Y/S)
    \longrightarrow
    \underline{\OmegaInf}^{\,p+q}(X\times_SY/S).
}
\]
After passage to the homotopy category these normalized external products are compatible with the de
Rham differential by the graded Leibniz identity.  No multiplicative compatibility with the \v{C}ech-totalization
filtration and no unconditional K\"unneth equivalence is asserted.
\end{proposition}

\begin{proof}
Left exactness of relative de Rham localization gives
\[
    (X\times_SY)_{\dR/S}
    \simeq
    X_{\dR/S}\times_SY_{\dR/S}.
\]
The product of the image factorizations
\[
    X\twoheadrightarrow Q_{X/S}\hookrightarrow X_{\dR/S},
    \qquad
    Y\twoheadrightarrow Q_{Y/S}\hookrightarrow Y_{\dR/S}
\]
is again an effective-epimorphism--monomorphism factorization.  Indeed, the product of the two
monomorphisms is a monomorphism, while
\[
    X\times_SY
    \longrightarrow
    Q_{X/S}\times_SY
    \longrightarrow
    Q_{X/S}\times_SQ_{Y/S}
\]
is a composite of pullbacks of the two effective epimorphisms.  Uniqueness of image factorization
therefore gives
\[
\boxed{
    Q_{X\times_SY/S}
    \simeq
    Q_{X/S}\times_SQ_{Y/S}.
}
\]
Consequently
\[
    \Sigma^\infty_{+,S}Q_{X\times_SY/S}
    \simeq
    \Sigma^\infty_{+,S}Q_{X/S}
    \otimes
    \Sigma^\infty_{+,S}Q_{Y/S}.
\]

Use the internal-function presentation of Theorem~\ref{thm:effective-quotient-de-rham}.  Closed
monoidality supplies
\[
\begin{aligned}
    &\underline{\Hom}
    \left(\Sigma^\infty_+Q_{X/S},\mathbb O_S^{\add}\right)
    \otimes
    \underline{\Hom}
    \left(\Sigma^\infty_+Q_{Y/S},\mathbb O_S^{\add}\right)
    \\
    &\qquad\longrightarrow
    \underline{\Hom}
    \left(
        \Sigma^\infty_+Q_{X/S}\otimes\Sigma^\infty_+Q_{Y/S},
        \mathbb O_S^{\add}\otimes\mathbb O_S^{\add}
    \right),
\end{aligned}
\]
followed by multiplication
\(\mathbb O_S^{\add}\otimes\mathbb O_S^{\add}\to\mathbb O_S^{\add}\).  Under the displayed
product identification of effective quotients, this is the required morphism
\[
    \mathbf{DR}_{X/S}\otimes\mathbf{DR}_{Y/S}
    \longrightarrow
    \mathbf{DR}_{X\times_SY/S}.
\]
Compatibility with convolution diagonals makes it a morphism of commutative algebras.

Degreewise the same product identification gives
\[
    R_{X\times_SY/S,n}
    \simeq
    R_{X/S,n}\times_SR_{Y/S,n}
\]
and hence a pairing of cosimplicial commutative algebras
\[
    \underline{\mathbf C}_{\dR}^{\,\bullet}(X/S)
    \otimes
    \underline{\mathbf C}_{\dR}^{\,\bullet}(Y/S)
    \longrightarrow
    \underline{\mathbf C}_{\dR}^{\,\bullet}(X\times_SY/S).
\]
Apply the paired consecutive-block Alexander--Whitney construction of
Proposition~\ref{prop:paired-ordered-monoidal-stable-dold-kan} to obtain the normalized external maps.
The homotopy-category Leibniz identity follows from the same proposition.
\end{proof}


\section{Moduli objects of normalized infinitesimal cochains}
\label{sec:monadic-de-rham-form-moduli}

The finite affine Nisnevich basis and affine density in slices were constructed in the preceding
chapter, and Theorem~\ref{thm:monadic-de-rham-affine-slice-presentation} upgraded them to the
affine-point presentation needed here.  We now apply that presentation to the normalized-cochain functors just
shown to satisfy Nisnevich descent.

\begin{definition}[Affine normalized-cochain functors]
\label{def:monadic-de-rham-form-functors}
For an affine spectral test \(T\to S\), define
\[
    \Omega_{\mathrm{inf},S}^q(T)
    :=
    \Omega^\infty\OmegaInf^q(T/S)
\]
and
\[
    \Omega_{\mathrm{inf},S}^{q,\cl}(T)
    :=
    \Omega^\infty\OmegaInf^{q,\cl}(T/S).
\]
\end{definition}

\begin{theorem}[Moduli objects of normalized and coherently closed cochains]
\label{thm:monadic-de-rham-form-moduli}
The functors \(\Omega_{\mathrm{inf},S}^q\) and \(\Omega_{\mathrm{inf},S}^{q,\cl}\) are sheaves for the affine spectral
Nisnevich topology.  Hence they determine canonical objects
\[
\boxed{
    \boldsymbol{\Omega}_{\mathrm{inf},S}^{q},
    \qquad
    \boldsymbol{\Omega}_{\mathrm{inf},S}^{q,\cl}
    \in
    \mathcal X_{\Sp,/S}.
}
\]
For every affine represented \(T\to S\),
\[
\boxed{
    \Map_S(T,\boldsymbol{\Omega}_{\mathrm{inf},S}^{q})
    \simeq
    \Omega^\infty\OmegaInf^q(T/S)
}
\]
and
\[
\boxed{
    \Map_S(T,\boldsymbol{\Omega}_{\mathrm{inf},S}^{q,\cl})
    \simeq
    \Omega^\infty\OmegaInf^{q,\cl}(T/S).
}
\]
The identity maps of the two moduli objects determine the tautological universal points of the
represented affine Nisnevich normalized-cochain sheaves.
\end{theorem}

\begin{proof}
By the finite affine basis theorem of the preceding chapter, the affine Nisnevich sheaf condition can
be checked on finite represented Nisnevich covers.  Let
\[
    \{T_\alpha\to T\}_{\alpha\in A}
\]
be such a cover, put \(U:=\coprod_{\alpha\in A}T_\alpha\), and let \(U_\bullet\to T\) be its
\v{C}ech nerve.  Theorem~\ref{thm:modal-de-rham-nisnevich-descent} gives descent for the normalized
and coherently closed cochain spectra along \(U_\bullet\to T\).  In simplicial degree \(k\),
\[
    U_k
    \simeq
    \coprod_{(\alpha_0,\ldots,\alpha_k)\in A^{k+1}}
    T_{\alpha_0}\times_T\cdots\times_TT_{\alpha_k}.
\]
Taking stable sections, which preserve limits, and applying
Proposition~\ref{prop:monadic-de-rham-finite-coproducts} identifies the resulting degree-\(k\) term
with the product of the normalized, respectively coherently closed, cochain spectra on the affine
intersections
\[
    T_{\alpha_0}\times_T\cdots\times_TT_{\alpha_k}.
\]
Thus the intrinsic descent equivalence is exactly the affine \v{C}ech descent condition.  Applying
\(\Omega^\infty\), which preserves limits, shows that
\(\Omega_{\mathrm{inf},S}^q\) and \(\Omega_{\mathrm{inf},S}^{q,\cl}\) are affine spectral
Nisnevich sheaves.  The affine-point presentation of
Theorem~\ref{thm:monadic-de-rham-affine-slice-presentation} therefore gives the two representing
objects in \(\mathcal X_{\Sp,/S}\).  Their identity maps give the tautological universal points of
the represented affine-test normalized-cochain sheaves by Yoneda.  No identification of these points
with the intrinsic normalized-cochain construction evaluated on the generally nonrepresented moduli
objects is asserted.
\end{proof}

\begin{proposition}[Moduli infinitesimal differential]
\label{prop:monadic-de-rham-moduli-differential}
For \(q\geq1\), the exact-cochain classifier and underlying-cochain map assemble to canonical morphisms
\[
\boxed{
    d_{\inf}^{\cl}:
    \boldsymbol{\Omega}_{\mathrm{inf},S}^{q-1}
    \longrightarrow
    \boldsymbol{\Omega}_{\mathrm{inf},S}^{q,\cl}
}
\]
and
\[
\boxed{
    u_q:
    \boldsymbol{\Omega}_{\mathrm{inf},S}^{q,\cl}
    \longrightarrow
    \boldsymbol{\Omega}_{\mathrm{inf},S}^{q}.
}
\]
Their composite represents the infinitesimal differential
\[
\boxed{
    d_{\inf}:
    \boldsymbol{\Omega}_{\mathrm{inf},S}^{q-1}
    \longrightarrow
    \boldsymbol{\Omega}_{\mathrm{inf},S}^{q}.
}
\]
\end{proposition}

\begin{proof}
The exact-cochain classifier and underlying-cochain map are natural under pullback and satisfy Nisnevich
descent.  Hence they define maps of the corresponding moduli sheaves.  Their composite is
\(d_{\inf}\) by Proposition~\ref{prop:exact-form-underlying-boundary}.
\end{proof}


\section{Realization on the previously constructed smooth stable site}
\label{sec:monadic-de-rham-smooth-site}

Assume now that \(S\) is represented and lies in the standing motivic scope.  Retain the notation
from the preceding chapter
\[
    \mathcal D_S:=\Sm^{\fp}_{/S},
    \qquad
    \mathcal C_S
    :=
    \operatorname{Shv}_{\Nis}(\mathcal D_S;\Sp).
\]
No new smooth-site stabilization or BNV reconstruction is introduced in this section.  We simply
restrict the intrinsic effective de Rham package to smooth represented tests and thereby obtain new coefficient
objects in the already established category \(\mathcal C_S\).
For every smooth represented test \(U\to S\), Theorem~\ref{thm:full-effective-dr-comparison}
identifies the full and effective de Rham algebras, so either presentation gives the same smooth-site
coefficient.

\begin{theorem}[Smooth-site effective de Rham and normalized-cochain coefficients]
\label{thm:monadic-de-rham-smooth-site-descent}
The assignments
\[
    U\longmapsto\mathbf R\Gamma_{\dR}(U/S),
\]
\[
    U\longmapsto
    \Gamma_S^{\mathrm{st}}
    \left(
        F_{\mathrm{Tot}}^q\mathbf{DR}_{U/S}
    \right),
\]
\[
    U\longmapsto\OmegaInf^q(U/S),
    \qquad
    U\longmapsto\OmegaInf^{q,\cl}(U/S)
\]
are spectral Nisnevich sheaves of spectra on \(\mathcal D_S\), hence objects of
\(\mathcal C_S\).  The \v{C}ech-totalization stages form a complete filtered coefficient object with associated
graded
\[
\boxed{
    U\longmapsto
    \OmegaInf^q(U/S)[-q].
}
\]
The infinitesimal boundary and exact-cochain classifier assemble to morphisms of smooth-site sheaves.
\end{theorem}

\begin{proof}
By the stable Nisnevich sheaf criterion of the preceding chapter, it is enough to verify the
empty-object condition and excision for spectral Nisnevich distinguished squares.  For a
distinguished square, the two morphisms to its terminal vertex form a finite represented spectral
Nisnevich covering family.  Form their coproduct and its \v{C}ech nerve.  Theorem~\ref{thm:modal-de-rham-nisnevich-descent}
gives intrinsic descent along that effective epimorphism.  Stable sections preserve the descent
limits, and Proposition~\ref{prop:monadic-de-rham-finite-coproducts} identifies every \v{C}ech level
with the product of the corresponding coefficient spectra on the finite intersections of the two
members of the covering family.  Thus the intrinsic descent equivalence is precisely the required
spectral Nisnevich descent, and hence the required excision square, for each assignment above.  On
the empty test the parameterized stable category is zero, so all displayed coefficient spectra
vanish.  Totalization completeness, associated graded, and the boundary maps are objectwise instances
of the intrinsic results.
\end{proof}

\begin{proposition}[Previously constructed stable Nisnevich inverse image]
\label{lem:monadic-de-rham-lan-nis-equivalences}
\label{prop:monadic-de-rham-stable-nis-base-change}
Let \(c:S'\to S\) be a morphism of represented bases in the standing motivic
scope, and let
\[
    c^*_{\Nis}:
    \mathcal C_S
    \longrightarrow
    \mathcal C_{S'}
\]
be the stable Nisnevich inverse-image functor constructed in
Chapter~\ref{chap:brave-new-infinitesimal-geometry}.  If
\[
    b_c:
    \mathcal D_S\longrightarrow\mathcal D_{S'},
    \qquad
    U\longmapsto U\times_SS',
\]
and
\[
    L_S:=L^{\mathrm{st}}_{\Nis,S},
    \qquad
    L_{S'}:=L^{\mathrm{st}}_{\Nis,S'},
\]
then the previously constructed inverse image is characterized by the canonical
localization formula
\[
\boxed{
    c^*_{\Nis}L_S
    \xrightarrow{\simeq}
    L_{S'}\operatorname{Lan}_{b_c^{\op}}.
}
\]
Equivalently, if
\[
    \iota_S:
    \mathcal C_S
    \hookrightarrow
    \Fun(\mathcal D_S^{\op},\Sp)
\]
is the inclusion, then
\[
\boxed{
    c^*_{\Nis}E
    \simeq
    L_{S'}
    \operatorname{Lan}_{b_c^{\op}}
    \iota_SE
}
\]
for \(E\in\mathcal C_S\).  The functor \(c^*_{\Nis}\) is exact and strong
symmetric monoidal, and for composable base changes
\[
    S''\xrightarrow{d}S'\xrightarrow{c}S
\]
there are canonical coherent equivalences
\[
    (c\circ d)^*_{\Nis}
    \simeq
    d^*_{\Nis}c^*_{\Nis}.
\]
No second smooth-site inverse-image functor is introduced in this chapter.
\end{proposition}

\begin{proof}
Chapter~\ref{chap:brave-new-infinitesimal-geometry} constructs
\(c^*_{\Nis}\) by left Kan extension along \(b_c^{\op}\) followed by stable
Nisnevich localization.  Base change carries stable Nisnevich excision
generators to their base changes, so the left Kan extension preserves the
stable Nisnevich local-equivalence class and factors through localization.
The universal property of \(L_S\) gives the displayed characterization.
Exactness, strong symmetric monoidality, and the coherent composition
equivalences are precisely the properties proved there for this functor.
\end{proof}


\begin{proposition}[Smooth-site base-change exchange]
\label{prop:monadic-de-rham-smooth-site-exchange}
Let \(c:S'\to S\) be a morphism of represented bases in the standing motivic scope, and let
\[
    c^*_{\Nis}:\mathcal C_S\longrightarrow\mathcal C_{S'}
\]
denote the previously constructed smooth-site inverse-image functor recalled in
Proposition~\ref{prop:monadic-de-rham-stable-nis-base-change}.  The intrinsic restriction maps
induce coherent exchange morphisms
\[
\boxed{
    c^*_{\Nis}
    \bigl(
        U\mapsto\mathbf R\Gamma_{\dR}(U/S)
    \bigr)
    \longrightarrow
    \bigl(
        U'\mapsto\mathbf R\Gamma_{\dR}(U'/S')
    \bigr)
}
\]
and analogous exchange maps for every totalization stage, normalized-cochain sheaf, and coherently closed normalized-cochain sheaf.
They preserve the available products, boundaries, totalization transitions, identities, composition, and
higher base-change pasting.  No arbitrary base-change equivalence for global sections is asserted.
\end{proposition}

\begin{proof}
Let \(\mathscr F_S\) denote any one of the smooth-site sheaves in
Theorem~\ref{thm:monadic-de-rham-smooth-site-descent}.  For every \(U\in\mathcal D_S\), intrinsic
parameterized base change followed by restriction on stable sections gives a natural morphism
\[
    \mathscr F_S(U)
    \longrightarrow
    \mathscr F_{S'}(U\times_SS').
\]
Equivalently, on stable presheaves this is a natural transformation
\[
    \iota_S\mathscr F_S
    \longrightarrow
    (-)\circ b_c^{\op}\,\iota_{S'}\mathscr F_{S'}.
\]
By the left-Kan-extension adjunction it corresponds to
\[
    \operatorname{Lan}_{b_c^{\op}}
    \iota_S\mathscr F_S
    \longrightarrow
    \iota_{S'}\mathscr F_{S'}.
\]
The target is already a stable Nisnevich sheaf, so the transformation factors canonically through
stable Nisnevich sheafification and gives
\[
    c^*_{\Nis}\mathscr F_S
    \longrightarrow
    \mathscr F_{S'}.
\]
Applying this construction to the total de Rham sheaf, every totalization stage, the normalized-cochain sheaves, and the
coherently closed-normalized-cochain sheaves gives the asserted exchange morphisms.  Naturality of the intrinsic
base-change and restriction transformations, together with the strong symmetric monoidality and
coherent functoriality of \(c^*_{\Nis}\), gives compatibility with products, boundaries, totalization transitions,
identities, composition, and higher base-change pasting.
\end{proof}

\begin{theorem}[Spectral \(\acute{e}\)tale Cartesianity of the smooth-site de Rham coefficients]
\label{thm:monadic-de-rham-smooth-site-etale-exchange}
Let
\[
    g:S'\longrightarrow S
\]
be a represented spectral \(\acute{e}\)tale morphism of finite presentation
between bases in the standing motivic scope.  For each smooth-site coefficient
system
\[
    \mathscr{DR},\qquad
    \mathscr F_{\dR}^{q},\qquad
    \mathscr\Omega_{\inf}^{q},\qquad
    \mathscr\Omega_{\inf}^{q,\cl},
\]
the exchange morphism of
Proposition~\ref{prop:monadic-de-rham-smooth-site-exchange} is an equivalence:
\[
\boxed{
    g^*_{\Nis}\mathscr F_S
    \xrightarrow{\simeq}
    \mathscr F_{S'}.
}
\]
These equivalences preserve totalization-stage transitions, the coherent
infinitesimal differential, coherently closed cochains, exact-cochain
classifiers, the available products, identities, composition, and all higher
\(\acute{e}\)tale base-change coherences.
\end{theorem}

\begin{proof}
Because \(g\) is smooth of finite presentation, the stable Nisnevich
coefficient categories carry the previously constructed adjunction
\[
    g_{\sharp,\Nis}
    \dashv
    g^*_{\Nis}.
\]
Let \(V\in\Sm^{\fp}_{/S'}\).  Since \(g\) is
\(\acute{e}\)tale of finite presentation, the composite \(V\to S'\to S\) is
again smooth of finite presentation.  On suspension representables,
\[
    g_{\sharp,\Nis}\Sigma^\infty_+h_V
    \simeq
    \Sigma^\infty_+h_V,
\]
where on the right \(V\) is regarded over \(S\) by composition with \(g\).
For any one of the four coefficient systems \(\mathscr F\),
Theorem~\ref{thm:formal-etale-change-of-base-intrinsic} gives a canonical
equivalence of its values
\[
    \mathscr F_S(V/S)
    \xrightarrow{\simeq}
    \mathscr F_{S'}(V/S').
\]
Hence
\[
\begin{aligned}
    \map_{\mathcal C_{S'}}
    \left(
        \Sigma^\infty_+h_V,
        g^*_{\Nis}\mathscr F_S
    \right)
    &\simeq
    \map_{\mathcal C_S}
    \left(
        g_{\sharp,\Nis}\Sigma^\infty_+h_V,
        \mathscr F_S
    \right)
    \\
    &\simeq
    \mathscr F_S(V/S)
    \\
    &\simeq
    \mathscr F_{S'}(V/S')
    \\
    &\simeq
    \map_{\mathcal C_{S'}}
    \left(
        \Sigma^\infty_+h_V,
        \mathscr F_{S'}
    \right).
\end{aligned}
\]
Under the left-Kan-extension adjunction this comparison is exactly the
exchange morphism already constructed in
Proposition~\ref{prop:monadic-de-rham-smooth-site-exchange}.  Integer shifts
of suspension representables generate \(\mathcal C_{S'}\), so the exchange is
an equivalence.  Compatibility with all displayed structures follows from
the corresponding compatibility in
Theorem~\ref{thm:formal-etale-change-of-base-intrinsic} and from the coherent
functoriality of the smooth-site exchange.
\end{proof}


\begin{remark}[Corrected interface with the Thom-stable differential category]
\label{rem:monadic-de-rham-corrected-interface}
The preceding chapter supplied the ambient stable Nisnevich and Thom-stable
differential machinery.  Its free level-zero functor
\[
    F_{0,S}:
    \mathcal C_S\longrightarrow\mathcal P_S^{\mathrm{pre}}
\]
lands first in the raw symmetric Thom-prespectrum presentation category.  The
actual differential realization used here is
\[
\boxed{
    \RealDiff_S
    =
    L_{\mathrm{Th},S}F_{0,S}:
    \mathcal C_S\longrightarrow\mathcal R_S^0.
}
\]
It is exact, colimit-preserving and strong symmetric monoidal.  It is not
asserted to be fully faithful.  Thus every smooth-site effective de Rham,
totalization, normalized-cochain, or coherently closed-cochain coefficient
constructed above has a canonical Thom-stable differential realization, after
which the BNV adjoint triple and the differential Umkehr formalism of
Chapter~\ref{chap:brave-new-infinitesimal-geometry} apply.
\end{remark}

\subsection{Notation for the smooth-site de Rham coefficients}
\label{subsec:monadic-de-rham-smooth-site-notation}

For later use write
\[
    \mathscr{DR}_S
    \in\mathcal C_S
\]
for the smooth-site sheaf
\[
    U\longmapsto\mathbf R\Gamma_{\dR}(U/S),
\]
and for \(q\geq0\) write
\[
    \mathscr F_{\dR,S}^q
    \in\mathcal C_S
\]
for
\[
    U\longmapsto
    \Gamma_S^{\mathrm{st}}
    \left(
        F_{\mathrm{Tot}}^q\mathbf{DR}_{U/S}
    \right).
\]
Similarly write
\[
    \mathscr\Omega_{\inf,S}^{q},
    \qquad
    \mathscr\Omega_{\inf,S}^{q,\cl}
    \in\mathcal C_S
\]
for the normalized and coherently closed smooth-site coefficient sheaves.

Their Thom-stable realizations are denoted
\[
\boxed{
    \widehat{\mathscr{DR}}_S
    :=
    \RealDiff_S\mathscr{DR}_S,
    \qquad
    \widehat{\mathscr F}_{\dR,S}^q
    :=
    \RealDiff_S\mathscr F_{\dR,S}^q,
}
\]
and
\[
\boxed{
    \widehat{\mathscr\Omega}_{\inf,S}^{q}
    :=
    \RealDiff_S\mathscr\Omega_{\inf,S}^{q},
    \qquad
    \widehat{\mathscr\Omega}_{\inf,S}^{q,\cl}
    :=
    \RealDiff_S\mathscr\Omega_{\inf,S}^{q,\cl}.
}
\]


\begin{proposition}[Exact smooth-site de Rham stage system]
\label{prop:monadic-de-rham-smooth-site-stage-exactness}
The objectwise adjacent-stage cofiber sequences assemble in
\(\mathcal C_S\) to canonical cofiber sequences
\[
\boxed{
    \mathscr F_{\dR,S}^{q+1}
    \longrightarrow
    \mathscr F_{\dR,S}^{q}
    \longrightarrow
    \mathscr\Omega_{\inf,S}^{q}[-q].
}
\]
The normalized smooth-site coefficients assemble to an actual coherent
nonnegative cochain object
\[
\boxed{
    \mathscr\Omega_{\inf,S}^{\bullet}
    \in
    \operatorname{CCh}_{\geq0}(\mathcal C_S),
}
\]
and the smooth-site exact-cochain classifiers are morphisms in
\(\mathcal C_S\) compatible with that coherent object.
\end{proposition}

\begin{proof}
For every smooth represented \(U\to S\), the intrinsic adjacent-stage
sequence
\[
    F_{\mathrm{Tot}}^{q+1}\mathbf{DR}_{U/S}
    \longrightarrow
    F_{\mathrm{Tot}}^{q}\mathbf{DR}_{U/S}
    \longrightarrow
    \underline{\OmegaInf}^{\,q}(U/S)[-q]
\]
is a cofiber sequence in \(\mathcal E_S\).  Stable sections
\[
    \Gamma_S^{\mathrm{st}}
    =
    \map_{\mathcal E_S}(\mathbf1_S,-)
\]
preserve finite limits, hence finite colimits, because the target is stable.
They therefore carry this sequence to a cofiber sequence of spectra.  The
three resulting assignments are already spectral Nisnevich sheaves by
Theorem~\ref{thm:monadic-de-rham-smooth-site-descent}.  Limits and finite
colimits in the stable sheaf category are computed objectwise, so the
objectwise cofiber sequences assemble to the displayed cofiber sequence in
\(\mathcal C_S\).

The coherent normalized-cochain object and the exact-cochain classifier are
natural in \(U\).  Applying the exact stable-section functor objectwise
preserves their specified nullhomotopies and all higher cells.  Nisnevich
descent then assembles them into the stated coherent object and morphisms in
\(\mathcal C_S\).
\end{proof}

\begin{proposition}[Exact realization of the coherent de Rham stage system]
\label{prop:monadic-de-rham-exact-thom-realization}
The functor \(\RealDiff_S\) carries every adjacent totalization-stage
fibre/cofiber sequence to the corresponding fibre/cofiber sequence in
\(\mathcal R_S^0\).  In particular,
\[
\boxed{
    \cofib
    \left(
        \widehat{\mathscr F}_{\dR,S}^{q+1}
        \longrightarrow
        \widehat{\mathscr F}_{\dR,S}^{q}
    \right)
    \simeq
    \widehat{\mathscr\Omega}_{\inf,S}^{q}[-q].
}
\]
The coherent infinitesimal differential, its specified nullhomotopy of the
square, all higher cochain coherences, and the exact-cochain classifier are
transported functorially to \(\mathcal R_S^0\).
\end{proposition}

\begin{proof}
By Proposition~\ref{prop:monadic-de-rham-smooth-site-stage-exactness}, the
adjacent-stage cofiber sequences and the coherent normalized-cochain object
already live as exact diagrams in \(\mathcal C_S\).  The functor \(F_{0,S}\)
is colimit-preserving between stable categories and Thom stabilization is an
exact reflective localization.  Hence their composite \(\RealDiff_S\) is
exact.  Applying it to those actual diagrams gives the displayed cofiber
sequences and transports the specified nullhomotopies and all higher coherent
cells.
\end{proof}

\begin{remark}[No formal preservation of the complete inverse limit]
\label{rem:monadic-de-rham-realization-no-complete-limit}
The preceding proposition concerns the already constructed stage objects and
their finite exact diagrams.  Exactness of \(\RealDiff_S\) does not by itself
identify the image of the sequential inverse limit defining complete
totalization with the inverse limit of the realized stages.  No such
identification is used in the integration theory below.
\end{remark}

\begin{proposition}[Stabilized level-zero evaluation]
\label{prop:monadic-de-rham-stabilized-level-zero-evaluation}
Let
\[
    j_{\mathrm{Th},S}:
    \mathcal R_S^0
    \hookrightarrow
    \mathcal P_S^{\mathrm{pre}}
\]
be the fully faithful right adjoint of Thom stabilization.  For every
\(A\in\mathcal C_S\) and \(E\in\mathcal R_S^0\), there is a canonical
equivalence
\[
\boxed{
    \map_{\mathcal R_S^0}
    \left(
        \RealDiff_S A,E
    \right)
    \simeq
    \map_{\mathcal C_S}
    \left(
        A,(j_{\mathrm{Th},S}E)_0
    \right).
}
\]
Consequently, for every \(U\in\Sm^{\fp}_{/S}\),
\[
\boxed{
    \map_{\mathcal R_S^0}
    \left(
        \RealDiff_S\Sigma^\infty_+h_U,E
    \right)
    \simeq
    E_0(U),
}
\]
where \(E_0\) is the zeroth level of the underlying Thom-local symmetric
prespectrum.
\end{proposition}

\begin{proof}
The two adjunctions give
\[
\begin{aligned}
    \map_{\mathcal R_S^0}(\RealDiff_SA,E)
    &=
    \map_{\mathcal R_S^0}
    (L_{\mathrm{Th},S}F_{0,S}A,E)
    \\
    &\simeq
    \map_{\mathcal P_S^{\mathrm{pre}}}
    (F_{0,S}A,j_{\mathrm{Th},S}E)
    \\
    &\simeq
    \map_{\mathcal C_S}
    \left(
        A,(j_{\mathrm{Th},S}E)_0
    \right).
\end{aligned}
\]
The represented formula is stable Yoneda.
\end{proof}

\begin{remark}[Evaluation is not full faithfulness]
\label{rem:monadic-de-rham-evaluation-not-full-faithfulness}
The preceding formula is an adjunction calculation against a Thom-local
target.  It does not imply that \(\RealDiff_S\) is fully faithful.
\end{remark}

\begin{proposition}[Smooth direct image commutes with differential realization]
\label{prop:monadic-de-rham-realization-smooth-direct-image}
Let
\[
    p:X\longrightarrow Y
\]
be smooth of finite presentation.  Then there is a canonical coherent
equivalence
\[
\boxed{
    p_\sharp^0\RealDiff_X
    \xrightarrow{\simeq}
    \RealDiff_Yp_{\sharp,\Nis}.
}
\]
There is also the already constructed pullback equivalence
\[
\boxed{
    p^{*,0}\RealDiff_Y
    \xrightarrow{\simeq}
    \RealDiff_Xp^*_{\Nis}.
}
\]
Both are compatible with identity, composition, and smooth base-change
pasting.
\end{proposition}

\begin{proof}
Before Thom stabilization, the levelwise smooth direct image on raw symmetric
Thom prespectra satisfies
\[
    p_\sharp^{\mathrm{pre}}F_{0,X}(A)
    \simeq
    F_{0,Y}(p_{\sharp,\Nis}A).
\]
The Thom-linearity theorem of
Chapter~\ref{chap:brave-new-infinitesimal-geometry} carries the stabilization
generators over \(X\) to the corresponding generators over \(Y\).  Therefore
the universal property of the Thom localization descends the preceding
equivalence to
\[
    p_\sharp^0L_{\mathrm{Th},X}F_{0,X}
    \simeq
    L_{\mathrm{Th},Y}F_{0,Y}p_{\sharp,\Nis}.
\]
This is the first formula.  The second is the pullback compatibility of the
realization functor already constructed in the preceding chapter.  All
coherences are inherited from the levelwise Kan-extension and Thom-linear
constructions.
\end{proof}

\begin{corollary}[Realized base-change exchange for the de Rham package]
\label{cor:monadic-de-rham-realized-base-change-exchange}
For every morphism \(c:S'\to S\) of represented bases in the standing
motivic scope, the smooth-site exchange morphisms of
Proposition~\ref{prop:monadic-de-rham-smooth-site-exchange} induce coherent
morphisms
\[
    c^{*,0}\widehat{\mathscr{DR}}_S
    \longrightarrow
    \widehat{\mathscr{DR}}_{S'},
\]
\[
    c^{*,0}\widehat{\mathscr F}_{\dR,S}^{q}
    \longrightarrow
    \widehat{\mathscr F}_{\dR,S'}^{q},
\]
and analogous morphisms for normalized and coherently closed cochains.
They preserve the stage transitions, exact-cochain classifiers, coherent
infinitesimal differential, identities, composition, and higher base-change
pasting.

No arbitrary base-change equivalence is asserted.  If \(c\) is represented
spectral \(\acute{e}\)tale of finite presentation, then all of the displayed
exchange morphisms, and their normalized and coherently closed analogues, are
equivalences by
Theorem~\ref{thm:monadic-de-rham-smooth-site-etale-exchange} and the pullback
compatibility of \(\RealDiff\).
\end{corollary}



\section{Basic presentations and comparison tests}
\label{sec:monadic-de-rham-basic-tests}

These tests use only the intrinsic constructions above and geometric calculations already established
in the preceding chapters.  No principal-parts algebra, represented formal-completion object, crystalline-jet model, or
HKR comparison is introduced.  The intrinsic completion used below is the pullback object of
Definition~\ref{def:intrinsic-relative-dr-completion}.

\subsection{The identity morphism}

\begin{proposition}[Terminal relative case]
\label{prop:monadic-de-rham-terminal-case}
For \(X=S\),
\[
    Q_{S/S}\simeq S,
    \qquad
    R_{S/S,\bullet}\simeq\operatorname{const}(S),
\]
and
\[
    \underline{\mathbf C}_{\dR}^{\,\bullet}(S/S)
    \simeq
    \operatorname{const}(\mathbb O_S^{\add}).
\]
Consequently
\[
\boxed{
    \underline{\OmegaInf}^{\,0}(S/S)
    \simeq
    \mathbb O_S^{\add},
    \qquad
    \underline{\OmegaInf}^{\,q}(S/S)\simeq0
    \quad(q>0),
}
\]
and
\[
\boxed{
    \mathbf{DR}_{S/S}
    \simeq
    \mathbb O_S^{\add}.
}
\]
For \(q>0\),
\[
    \mathfrak I^{\langle q\rangle}_{S/S}\simeq0
\]
and every positive coherently closed normalized-cochain spectrum vanishes.
\end{proposition}

\begin{proof}
The terminal object of the slice is fixed by relative de Rham localization, so the relative de Rham
unit is the identity of \(S\).  Its effective quotient is \(S\), and every relation degree is \(S\).
The cosimplicial function object is therefore constant.  Its normalized terms vanish in positive
degree, and the positive skeletal quotients of a constant simplicial object vanish.
\end{proof}

\subsection{Relative de Rham-local and \texorpdfstring{\(\acute{e}\)tale}{etale} morphisms}

\begin{proposition}[Vanishing of positive relative cochains]
\label{prop:monadic-de-rham-etale-case}
If \(X\to S\) is relative de Rham-local, in particular if it is represented by a spectral
\(\acute{e}\)tale morphism, then
\[
    Q_{X/S}\simeq X,
    \qquad
    R_{X/S,n}\simeq X
\]
for every \(n\).  Hence
\[
    \underline{\mathbf C}_{\dR}^{\,\bullet}(X/S)
    \simeq
    \operatorname{const}
    \left(
        \Pi_p p^*\mathbb O_S^{\add}
    \right),
\]
\[
\boxed{
    \underline{\OmegaInf}^{\,q}(X/S)\simeq0
    \qquad(q>0),
}
\]
and
\[
\boxed{
    \mathbf{DR}_{X/S}
    \simeq
    \Pi_p p^*\mathbb O_S^{\add}.
}
\]
\end{proposition}

\begin{proof}
If \(X\to X_{\dR/S}\) is an equivalence, its effective image is \(X\) and every iterated relation
term is \(X\) with identity simplicial operators.  Normalize and totalize the resulting constant
cosimplicial object.  The preceding chapter proves that represented spectral
\(\acute{e}\)tale morphisms are relative de Rham-local.
\end{proof}

\subsection{The first normalized infinitesimal cochain object}

\begin{proposition}[First normalized-cochain presentation]
\label{prop:monadic-de-rham-first-form}
There is a canonical fibre sequence
\[
\boxed{
    \underline{\OmegaInf}^{\,1}(X/S)
    \longrightarrow
    \Pi_{a_1}a_1^*\mathbb O_S^{\add}
    \xrightarrow{s^0}
    \Pi_p p^*\mathbb O_S^{\add}.
}
\]
Thus a normalized one-cochain is an additive scalar function on an infinitesimally related pair of
points together with a specified nullhomotopy of its restriction along the degenerate pair defined by
the diagonal.
\end{proposition}

\begin{proof}
This is the degree-one case of Proposition~\ref{prop:explicit-modal-normalized-terms}.
\end{proof}

\subsection{Smooth spectral schemes}

\begin{proposition}[Smooth effective presentation]
\label{prop:monadic-de-rham-smooth-test}
If \(p:X\to S\) is represented and smooth, then
\[
    Q_{X/S}\simeq X_{\dR/S}
\]
and
\[
\boxed{
    \mathbf{DR}_{X/S}
    \simeq
    \Pi_{q_{\dR}}q_{\dR}^*\mathbb O_S^{\add}.
}
\]
Equivalently,
\[
\boxed{
    \mathbf R\Gamma_{\dR}(X/S)
    \simeq
    \map_{\mathcal E_{X_{\dR/S}}}
    \left(
        \mathbf 1_{X_{\dR/S}},
        q_{\dR}^*\mathbb O_S^{\add}
    \right).
}
\]
The \v{C}ech-totalization filtration is the dual totalization/skeletal filtration of the effective \v{C}ech nerve
\[
    X^{\times_{X_{\dR/S}}(\bullet+1)}.
\]
\end{proposition}

\begin{proof}
Theorem~\ref{thm:monadic-de-rham-smooth-effectivity} identifies the effective quotient with
\(X_{\dR/S}\).  Substitute this into Theorem~\ref{thm:effective-quotient-de-rham} and
Theorem~\ref{thm:skeletal-hodge-presentation}.
\end{proof}

\subsection{Smooth local linear model of the infinitesimal simplices}

\begin{theorem}[Local linear completion model]
\label{thm:monadic-de-rham-local-linear-completion}
Let \(p:X\to S\) be a represented smooth spectral morphism and let \(x\in|X|\).  Choose an affine
open \(V\subseteq S\) containing \(p(x)\) and an affine open
\(U\subseteq X\times_SV\) containing \(x\).  After shrinking \(U\) and \(V\) if necessary, for every
\(n\geq1\) there is an open neighbourhood
\[
    W_n\subseteq U^{\times_V(n+1)}
\]
of the diagonal and an \(\acute{e}\)tale morphism
\[
    \phi_n:
    W_n
    \longrightarrow
    V_U\left(L_{U/V}^{\oplus n}\right)
\]
which carries the diagonal to the zero section and whose induced morphism on
conormal modules is canonically homotopic to the identity of
\(L_{U/V}^{\oplus n}\).  The induced maps on intrinsic completions are
equivalences, and hence
\[
\boxed{
    R_{U/V,n}
    \simeq
    \widehat{
        V_U\left(L_{U/V}^{\oplus n}\right)
    }^{\,\dR}_{U/V}.
}
\]
Thus the infinitesimal relation simplices of a smooth spectral morphism are Zariski-locally the
intrinsic de Rham completions of the zero section in the linear spaces with coordinate module
\(L_{X/S}^{\oplus n}\).
\end{theorem}

\begin{proof}
Consider first the smooth affine morphism
\[
    \operatorname{pr}_0:
    U\times_VU
    \longrightarrow
    U
\]
with diagonal section
\[
    \Delta_1:U\longrightarrow U\times_VU.
\]
Cotangent base change and transitivity for the section give the canonical conormal identification
\[
    N^*_{U/(U\times_VU)}
    \simeq
    L_{U/V}.
\]
Apply the spectral tubular neighbourhood theorem of Chapter~3, Theorem~3.3.3, to this smooth affine
morphism and section.  Since the base \(U\) is affine, the affine-base clause of that theorem allows
the base cover to consist of \(U\) itself.  Thus there is an open neighbourhood
\[
    W_1\subseteq U\times_VU
\]
of the diagonal and an \(\acute{e}\)tale morphism over \(U\)
\[
    \phi_1:
    W_1
    \longrightarrow
    V_U(L_{U/V})
\]
which carries the diagonal to the zero section and whose induced morphism on
conormal modules is canonically homotopic to
\[
    \id_{L_{U/V}}.
\]

For \(n\geq1\), there is a canonical identification over the first \(U\)-factor
\[
    U^{\times_V(n+1)}
    \simeq
    \underbrace{
        (U\times_VU)\times_U\cdots\times_U(U\times_VU)
    }_{n\text{ factors}}.
\]
Define
\[
\boxed{
    W_n
    :=
    \underbrace{
        W_1\times_U\cdots\times_UW_1
    }_{n\text{ factors}}.
}
\]
Equivalently,
\[
    W_n
    =
    \bigcap_{j=1}^n
    \operatorname{pr}_{0j}^{-1}(W_1)
    \subseteq
    U^{\times_V(n+1)}.
\]
Hence \(W_n\) is an open neighbourhood of the full diagonal.

The relative linear-space construction carries direct sums to fibre products:
\[
    \underbrace{
        V_U(L_{U/V})\times_U\cdots\times_UV_U(L_{U/V})
    }_{n\text{ factors}}
    \simeq
    V_U(L_{U/V}^{\oplus n}).
\]
Define
\[
\boxed{
    \phi_n
    :=
    \phi_1^{\times_Un}:
    W_n
    \longrightarrow
    V_U(L_{U/V}^{\oplus n}).
}
\]
Since spectral \(\acute{e}\)tale morphisms are stable under fibre products, \(\phi_n\) is
\(\acute{e}\)tale.  It carries the diagonal to the zero section.

Cotangent base change gives
\[
    L_{U^{\times_V(n+1)}/U}
    \simeq
    \bigoplus_{j=1}^n\operatorname{pr}_j^*L_{U/V},
\]
and transitivity for the diagonal therefore gives
\[
\boxed{
    N^*_{U/U^{\times_V(n+1)}}
    \simeq
    L_{U/V}^{\oplus n}.
}
\]
The conormal module of the zero section of
\(V_U(L_{U/V}^{\oplus n})\) is canonically the same module.  Because
\(\phi_n\) is the \(n\)-fold fibre product of \(\phi_1\), its induced
conormal morphism is the direct sum of the \(n\) conormal morphisms induced
by \(\phi_1\).  Each of these is canonically homotopic to
\(\id_{L_{U/V}}\), so their direct sum is canonically homotopic to
\[
    \id_{L_{U/V}^{\oplus n}}.
\]

The open immersion
\[
    W_n\hookrightarrow U^{\times_V(n+1)}
\]
and the morphism \(\phi_n\) are spectral \(\acute{e}\)tale, hence formally
\(\acute{e}\)tale.  Proposition~\ref{prop:formal-etale-completion-invariance}, applied first to the
open immersion and then to \(\phi_n\), gives
\[
\begin{aligned}
    \widehat{U^{\times_V(n+1)}}^{\,\dR}_{U/V}
    &\simeq
    \widehat{W_n}^{\,\dR}_{U/V}
    \\
    &\simeq
    \widehat{V_U(L_{U/V}^{\oplus n})}^{\,\dR}_{U/V}.
\end{aligned}
\]
Theorem~\ref{thm:relation-simplices-as-completions} identifies the left side with
\(R_{U/V,n}\), proving the displayed equivalence.  Finally,
\[
    L_{U/V}\simeq L_{X/S}|_U
\]
by open restriction and cotangent base change.  The single construction of \(W_1\) therefore
supplies the asserted models simultaneously for every \(n\geq1\), with no \(n\)-dependent
shrinking.
\end{proof}

\subsection{Brave-new affine space}

\begin{proposition}[Affine-space presentation]
\label{prop:monadic-de-rham-affine-space}
Let \(S\) be represented by a spectral scheme and let
\[
    X=\mathbb A_S^d
    =
    \Spec_S
    \Sym_{\mathcal O_S}
    \left(
        \mathcal O_S^{\oplus d}
    \right).
\]
Then \(X\to S\) is smooth and
\[
    L_{X/S}\simeq\mathcal O_X^{\oplus d}.
\]
Its intrinsic de Rham algebra is
\[
\boxed{
    \mathbf{DR}_{X/S}
    \simeq
    \Pi_{q_{\dR}}q_{\dR}^*\mathbb O_S^{\add},
}
\]
and \(\underline{\OmegaInf}^{\,q}(X/S)\) is the normalization of additive functions on
\[
    \underbrace{
        X\times_{X_{\dR/S}}\cdots\times_{X_{\dR/S}}X
    }_{q+1\text{ factors}}.
\]
The cotangent formula records the linear first-order geometry already established in the
foundational chapters.  Proposition~\ref{prop:stabilized-cotangent-is-cotangent} identifies the
first-order target of normalized one-cochains with \(\mathcal O_X^{\oplus d}\), while
Proposition~\ref{prop:normalized-one-cochains-strictly-nonlinear} shows that the full normalized
cochain generally retains higher infinitesimal terms.  No exterior-power identification is asserted.
\end{proposition}

\begin{proof}
The smoothness and cotangent statements are the brave-new affine-space calculation from the
preceding spectral geometry:
the cotangent complex of a relative free commutative algebra is the extension of scalars of its
generating module.  Apply Theorem~\ref{thm:monadic-de-rham-smooth-effectivity} and the definition of
normalized cochains.
\end{proof}

\subsection{Relative linear spaces}

\begin{proposition}[Relative linear-space presentation]
\label{prop:monadic-de-rham-linear-space}
Let \(S\) be represented by a spectral scheme, let \(P\) be a finite locally free
\(\mathcal O_S\)-module, and
\[
    X=V_S(P)=\Spec_S\Sym_{\mathcal O_S}(P).
\]
Then \(X\to S\) is smooth and
\[
    L_{X/S}\simeq p^*P.
\]
Hence
\[
\boxed{
    \mathbf{DR}_{X/S}
    \simeq
    \Pi_{q_{\dR}}q_{\dR}^*\mathbb O_S^{\add},
}
\]
while its normalized-cochain degrees and totalization stages are obtained intrinsically from the normalized infinitesimal
groupoid.
\end{proposition}

\begin{proof}
The cotangent and smoothness statements are the relative linear-space theorem already proved in the
spectral smooth geometry.  Apply the smooth effective presentation.
\end{proof}

\subsection{Discrete classical smooth input}

\begin{proposition}[Classical smooth input]
\label{prop:monadic-de-rham-classical-smooth}
Let \(X\to S\) be the discrete spectral enhancement of a smooth morphism of ordinary schemes.  The
construction of this chapter canonically associates
\[
    X\to X_{\dR/S},
    \qquad
    R_{X/S,\bullet}
    =
    X^{\times_{X_{\dR/S}}(\bullet+1)},
\]
the cosimplicial additive-function algebra
\[
    \underline{\mathbf C}_{\dR}^{\,\bullet}(X/S),
\]
the normalized infinitesimal cochain spectra
\[
    \underline{\OmegaInf}^{\,q}(X/S),
\]
the total de Rham algebra
\[
    \mathbf{DR}_{X/S},
\]
its complete \v{C}ech-totalization filtration, and the moduli objects of normalized and coherently closed cochains.  On
represented relation degrees, the de Rham cochains are spectra of functions on the intrinsic
infinitesimal simplices.  No coordinate choice and no ordinary differential-form complex enters the
definition.  The first-order linearization recovers the ordinary universal derivation on the discrete smooth
input, but this proposition does not identify the full normalized infinitesimal cochains or their
totalization filtration with the classical exterior de Rham complex or its Hodge filtration.
\end{proposition}

\begin{proof}
A discrete smooth morphism is a represented smooth spectral morphism.  Apply the smooth effective
presentation, represented scalar-function formula, totalization filtration, and moduli construction.
\end{proof}



\section{Differential exceptional integration of the realized de Rham coefficients}
\label{sec:dr-int-differential-umkehr}

Assume throughout this section that
\[
    p:X\longrightarrow Y
\]
is smooth, separated, and of finite presentation.

\subsection{Differential dualizing coefficients}

\begin{definition}[Differential dualizing coefficient]
\label{def:dr-int-dualizing-coefficient}
For \(D\in\mathcal R_Y^0\), define its differential exceptional pullback by
\[
\boxed{
    D_{X/Y}^{!,0}
    :=
    \widetilde p^{\,!}D
    =
    \Sigma^{T_p,0}p^{*,0}D.
}
\]
If \(D\) is one of
\[
    \widehat{\mathscr{DR}}_Y,\qquad
    \widehat{\mathscr F}_{\dR,Y}^{q},\qquad
    \widehat{\mathscr\Omega}_{\inf,Y}^{q},\qquad
    \widehat{\mathscr\Omega}_{\inf,Y}^{q,\cl},
\]
we call \(D_{X/Y}^{!,0}\) the corresponding
\emph{differential dualizing de Rham coefficient}.
\end{definition}

\begin{remark}[Why the dualizing coefficient is the general source]
\label{rem:dr-int-dualizing-source}
The general exceptional adjunction integrates classes with coefficients in
\(\widetilde p^{\,!}D\).  The smooth-site de Rham exchange only supplies a
morphism
\[
    p^{*,0}D
    \longrightarrow
    D_X
\]
for the intrinsic source coefficient \(D_X\), and the tangent Thom twist is
still present in \(\widetilde p^{\,!}\).  Thus no general equivalence
\[
    D_X\simeq\widetilde p^{\,!}D_Y
\]
is asserted without additional orientation or coefficient data.
\end{remark}

\begin{definition}[Differential Umkehr integration]
\label{def:dr-int-differential-integration}
For \(K\in\mathcal R_X^0\) and \(D\in\mathcal R_Y^0\), define the differential
Umkehr integration equivalence
\[
\boxed{
    \int_p^\partial:
    \map_{\mathcal R_X^0}
    \left(
        K,\widetilde p^{\,!}D
    \right)
    \xrightarrow{\simeq}
    \map_{\mathcal R_Y^0}
    \left(
        \widetilde p_!K,D
    \right)
}
\]
to be the mapping-spectrum adjunction associated with
\[
    \widetilde p_!\dashv\widetilde p^{\,!}.
\]
\end{definition}

\begin{definition}[Differential exceptional trace]
\label{def:dr-int-differential-trace}
For \(D\in\mathcal R_Y^0\), define
\[
\boxed{
    \Trdiff_{p,D}:
    \widetilde p_!\widetilde p^{\,!}D
    \longrightarrow
    D
}
\]
to be the counit of the differential exceptional adjunction.
\end{definition}

\begin{proposition}[Integration by the differential trace]
\label{prop:dr-int-integration-by-trace}
A morphism
\[
    \omega:
    K\longrightarrow\widetilde p^{\,!}D
\]
is carried by \(\int_p^\partial\) to the composite
\[
\boxed{
    \widetilde p_!K
    \xrightarrow{\widetilde p_!\omega}
    \widetilde p_!\widetilde p^{\,!}D
    \xrightarrow{\Trdiff_{p,D}}
    D.
}
\]
In particular the identity of \(\widetilde p^{\,!}D\) is carried to the
differential exceptional trace.
\end{proposition}

\begin{proof}
This is the counit description of the adjunction equivalence.
\end{proof}

\subsection{Motivic shadow}

\begin{lemma}[Homotopification preserves the smooth counit]
\label{lem:dr-int-smooth-counit-homotopification}
Let
\[
    p:X\longrightarrow Y
\]
be smooth and of finite presentation.  Write
\[
    \varepsilon_p^{\sharp,0}:
    p_\sharp^0p^{*,0}\longrightarrow\id_{\mathcal R_Y^0}
\]
for the counit of
\[
    p_\sharp^0\dashv p^{*,0},
\]
and
\[
    \varepsilon_p^\sharp:
    p_\sharp p^*\longrightarrow\id_{\SH(Y)}
\]
for the counit of the motivic smooth adjunction.  Under the canonical
comparisons
\[
    \mathbb H_Yp_\sharp^0
    \xrightarrow{\simeq}
    p_\sharp\mathbb H_X,
    \qquad
    \mathbb H_Xp^{*,0}
    \xrightarrow{\simeq}
    p^*\mathbb H_Y,
\]
there is a canonical commutative square
\[
\begin{tikzcd}[column sep=huge,row sep=large]
    \mathbb H_Yp_\sharp^0p^{*,0}
        \ar[r,"\mathbb H_Y(\varepsilon_p^{\sharp,0})"]
        \ar[d,"\simeq"']
    &
    \mathbb H_Y
        \ar[d,equal]
    \\
    p_\sharp p^*\mathbb H_Y
        \ar[r,"\varepsilon_p^\sharp\mathbb H_Y"']
    &
    \mathbb H_Y.
\end{tikzcd}
\]
Thus affine-line homotopification carries the Thom-stable smooth counit to
the motivic smooth counit.
\end{lemma}

\begin{proof}
We first prove the assertion on a motivic-local coefficient.

Let \(N\in\SH(Y)\).  For \(M\in\SH(X)\), the previously constructed
comparisons give natural equivalences
\[
\begin{aligned}
    \map_{\SH(Y)}
    \left(
        p_\sharp M,N
    \right)
    &\simeq
    \map_{\SH(Y)}
    \left(
        \mathbb H_Yp_\sharp^0\Disc_XM,N
    \right)
    \\
    &\simeq
    \map_{\mathcal R_Y^0}
    \left(
        p_\sharp^0\Disc_XM,\Disc_YN
    \right)
    \\
    &\simeq
    \map_{\mathcal R_X^0}
    \left(
        \Disc_XM,p^{*,0}\Disc_YN
    \right)
    \\
    &\simeq
    \map_{\mathcal R_X^0}
    \left(
        \Disc_XM,\Disc_Xp^*N
    \right)
    \\
    &\simeq
    \map_{\SH(X)}
    \left(
        M,p^*N
    \right).
\end{aligned}
\]
The second line is the localization adjunction
\(\mathbb H_Y\dashv\Disc_Y\), the third is
\(p_\sharp^0\dashv p^{*,0}\), and the fourth is the canonical
local-object comparison
\[
    p^{*,0}\Disc_Y
    \simeq
    \Disc_Xp^*.
\]
Hence this chain is precisely the descended adjunction
\[
    p_\sharp\dashv p^*.
\]

Take \(M=p^*N\).  Under the displayed chain, the differential smooth
counit
\[
    \varepsilon_p^{\sharp,0}(\Disc_YN):
    p_\sharp^0p^{*,0}\Disc_YN
    \longrightarrow
    \Disc_YN
\]
corresponds to the identity of
\(p^{*,0}\Disc_YN\).  After the canonical identification
\[
    p^{*,0}\Disc_YN
    \simeq
    \Disc_Xp^*N,
\]
this is the identity of \(p^*N\).  Therefore, by the defining
counit characterization of the descended adjunction,
\[
    \mathbb H_Y
    \left(
        \varepsilon_p^{\sharp,0}(\Disc_YN)
    \right)
\]
identifies with
\[
    \varepsilon_p^\sharp(N):
    p_\sharp p^*N\longrightarrow N.
\]

Now let \(D\in\mathcal R_Y^0\) be arbitrary, and let
\[
    \eta_D:
    D\longrightarrow
    \Disc_Y\mathbb H_YD
\]
be its affine-line localization unit.  Naturality of the differential
smooth counit gives
\[
\begin{tikzcd}[column sep=huge,row sep=large]
    p_\sharp^0p^{*,0}D
        \ar[r,"\varepsilon_p^{\sharp,0}(D)"]
        \ar[d,"p_\sharp^0p^{*,0}\eta_D"']
    &
    D
        \ar[d,"\eta_D"]
    \\
    p_\sharp^0p^{*,0}\Disc_Y\mathbb H_YD
        \ar[r,"\varepsilon_p^{\sharp,0}(\Disc_Y\mathbb H_YD)"']
    &
    \Disc_Y\mathbb H_YD.
\end{tikzcd}
\]
After applying \(\mathbb H_Y\), the right vertical arrow is an
equivalence.  The left vertical arrow is also an equivalence: using the
smooth-direct-image and pullback comparisons it identifies with
\[
    p_\sharp p^*\mathbb H_Y(\eta_D),
\]
and \(\mathbb H_Y(\eta_D)\) is an equivalence.  The lower horizontal
arrow has already been identified with the motivic smooth counit.
Therefore the upper horizontal arrow has the same identification.  This
proves the assertion for arbitrary \(D\).
\end{proof}

\begin{lemma}[Homotopification preserves Thom cancellation]
\label{lem:dr-int-thom-cancellation-homotopification}
Let \(X\) be in the standing motivic scope and let
\(\xi\in K(X)\).  Write
\[
    c_\xi^0:
    \Sigma^{-\xi,0}\Sigma^{\xi,0}
    \xrightarrow{\simeq}
    \id_{\mathcal R_X^0}
\]
and
\[
    c_\xi:
    \Sigma^{-\xi}\Sigma^\xi
    \xrightarrow{\simeq}
    \id_{\SH(X)}
\]
for the canonical Picard cancellation equivalences.  Under
\[
    \mathbb H_X\Sigma^{\pm\xi,0}
    \simeq
    \Sigma^{\pm\xi}\mathbb H_X,
\]
one has
\[
\boxed{
    \mathbb H_X(c_\xi^0)
    \simeq
    c_\xi\,\mathbb H_X.
}
\]
The analogous assertion holds for the opposite cancellation
\(\Sigma^{\xi,0}\Sigma^{-\xi,0}\simeq\id\).
\end{lemma}

\begin{proof}
The differential motivic \(J\)-homomorphism
\[
    J_X^0:
    K(X)\longrightarrow\Pic(\mathcal R_X^0)
\]
is a morphism of grouplike \(E_\infty\)-spaces, and its composite with
the strong symmetric monoidal homotopification functor is canonically
the motivic \(J\)-homomorphism
\[
    J_X:
    K(X)\longrightarrow\Pic(\SH(X)).
\]
The cancellation \(c_\xi^0\) is the image in
\(\Pic(\mathcal R_X^0)\) of the canonical grouplike relation
\[
    -\xi+\xi\simeq0
\]
in \(K(X)\).  Applying the induced morphism of Picard
\(E_\infty\)-groupoids carries this relation to the corresponding
relation
\[
    -\xi+\xi\simeq0
\]
for \(J_X\).  Strong monoidality identifies the resulting morphism with
the motivic cancellation
\[
    \Sigma^{-\xi}\Sigma^\xi\simeq\id.
\]
Hence the displayed comparison commutes with the canonical cancellation
equivalences.
\end{proof}

\begin{lemma}[Smooth purity identifies the exceptional counit]
\label{lem:dr-int-purity-counit}
Let
\[
    p:X\longrightarrow Y
\]
be smooth, separated, and of finite presentation.  Write
\[
    \pi_p:
    \Sigma^{T_p}p^*
    \xrightarrow{\simeq}
    p^!
\]
for smooth purity and
\[
    \chi_p:
    p_!
    \xrightarrow{\simeq}
    p_\sharp\Sigma^{-T_p}
\]
for the canonical smooth exceptional-direct-image equivalence of
Chapter~8.  Let
\[
    \varepsilon_p^!:
    p_!p^!\longrightarrow\id
\]
be the exceptional counit.

Then \(\varepsilon_p^!\) is canonically the composite
\[
\boxed{
\begin{aligned}
    p_!p^!
    &\xrightarrow{\chi_p\,p^!}
    p_\sharp\Sigma^{-T_p}p^!
    \\
    &\xrightarrow{
        p_\sharp\Sigma^{-T_p}\pi_p^{-1}
    }
    p_\sharp
    \Sigma^{-T_p}\Sigma^{T_p}p^*
    \\
    &\xrightarrow{
        p_\sharp c_{T_p}p^*
    }
    p_\sharp p^*
    \\
    &\xrightarrow{\varepsilon_p^\sharp}
    \id .
\end{aligned}
}
\]
Thus, under smooth purity, the exceptional adjunction is precisely the
smooth adjunction conjugated by the tangent Thom autoequivalence.
\end{lemma}

\begin{proof}
Put
\[
    L:=p_\sharp\Sigma^{-T_p},
    \qquad
    R:=\Sigma^{T_p}p^*.
\]
The inverse Thom twist is left adjoint to the positive Thom twist, so
composition of adjunctions gives
\[
    L\dashv R.
\]
Its counit is exactly
\[
    p_\sharp
    \Sigma^{-T_p}\Sigma^{T_p}p^*
    \xrightarrow{p_\sharp c_{T_p}p^*}
    p_\sharp p^*
    \xrightarrow{\varepsilon_p^\sharp}
    \id.
\]

Smooth purity is an equivalence of right adjoints
\[
    \pi_p:R\xrightarrow{\simeq}p^!.
\]
The canonical equivalence
\[
    \chi_p:p_!\xrightarrow{\simeq}L
\]
of Chapter~8 is the corresponding left-adjoint mate: this is precisely
the canonical equivalence obtained there from uniqueness of left
adjoints to the purity equivalence.

For adjunctions
\[
    L\dashv R,
    \qquad
    L'\dashv R'
\]
and an equivalence
\[
    \beta:R\xrightarrow{\simeq}R'
\]
with left mate
\[
    \alpha:L'\xrightarrow{\simeq}L,
\]
the defining mate identity gives
\[
    \varepsilon'
    =
    \varepsilon
    \circ
    L(\beta^{-1})
    \circ
    \alpha R'.
\]
Indeed, under the adjunction
\(L'\dashv R'\), the right-hand composite corresponds to
\[
    R'
    \xrightarrow{\beta^{-1}}
    R
    \xrightarrow{\beta}
    R',
\]
which is the identity, and hence it is the counit of
\(L'\dashv R'\).

Apply this identity with
\[
    L'=p_!,
    \qquad
    R'=p^!,
    \qquad
    \beta=\pi_p,
    \qquad
    \alpha=\chi_p.
\]
Substituting the explicit counit of \(L\dashv R\) gives the displayed
formula.
\end{proof}

\begin{theorem}[Differential exceptional integration refines motivic integration]
\label{thm:dr-int-motivic-shadow}
For every \(D\in\mathcal R_Y^0\), the homotopification of the differential
trace is canonically the motivic exceptional counit:
\[
\boxed{
    \mathbb H_Y
    \left(
        \Trdiff_{p,D}
    \right)
    \simeq
    \operatorname{tr}_p(\mathbb H_YD):
    p_!p^!\mathbb H_YD
    \longrightarrow
    \mathbb H_YD.
}
\]
More generally, if
\[
    \omega:
    K\longrightarrow\widetilde p^{\,!}D,
\]
then the homotopification of its differential adjunct is the motivic adjunct
of
\[
    \mathbb H_X\omega:
    \mathbb H_XK
    \longrightarrow
    p^!\mathbb H_YD.
\]
\end{theorem}

\begin{proof}
Write
\[
    c_{T_p}^0:
    \Sigma^{-T_p,0}\Sigma^{T_p,0}
    \xrightarrow{\simeq}
    \id
\]
for differential Thom cancellation, and
\[
    \varepsilon_p^{\sharp,0}:
    p_\sharp^0p^{*,0}\longrightarrow\id
\]
for the smooth Thom-stable counit.

The adjunction
\[
    \widetilde p_!
    =
    p_\sharp^0\Sigma^{-T_p,0}
    \dashv
    \Sigma^{T_p,0}p^{*,0}
    =
    \widetilde p^{\,!}
\]
is the composite of the adjunction
\[
    \Sigma^{-T_p,0}
    \dashv
    \Sigma^{T_p,0}
\]
with
\[
    p_\sharp^0\dashv p^{*,0}.
\]
Consequently its counit is canonically
\[
\boxed{
\begin{aligned}
    \Trdiff_{p,D}:\,
    p_\sharp^0
    \Sigma^{-T_p,0}
    \Sigma^{T_p,0}
    p^{*,0}D
    &\xrightarrow{
        p_\sharp^0c_{T_p}^0p^{*,0}
    }
    p_\sharp^0p^{*,0}D
    \\
    &\xrightarrow{
        \varepsilon_p^{\sharp,0}(D)
    }
    D.
\end{aligned}
}
\]

Apply \(\mathbb H_Y\).  The smooth-direct-image comparison,
the pullback comparison, and the differential \(J\)-homomorphism give
\[
\begin{aligned}
    \mathbb H_Y
    \widetilde p_!\widetilde p^{\,!}D
    &\simeq
    p_\sharp
    \Sigma^{-T_p}
    \Sigma^{T_p}
    p^*
    \mathbb H_YD.
\end{aligned}
\]
By
Lemma~\ref{lem:dr-int-thom-cancellation-homotopification}, the
homotopification of the first arrow in the displayed formula for
\(\Trdiff_{p,D}\) is
\[
    p_\sharp
    \Sigma^{-T_p}\Sigma^{T_p}p^*\mathbb H_YD
    \xrightarrow{
        p_\sharp c_{T_p}p^*
    }
    p_\sharp p^*\mathbb H_YD.
\]
By
Lemma~\ref{lem:dr-int-smooth-counit-homotopification}, the
homotopification of the second arrow is
\[
    \varepsilon_p^\sharp:
    p_\sharp p^*\mathbb H_YD
    \longrightarrow
    \mathbb H_YD.
\]
Therefore
\[
    \mathbb H_Y(\Trdiff_{p,D})
\]
is the Thom-conjugated smooth counit
\[
\begin{aligned}
    p_\sharp
    \Sigma^{-T_p}\Sigma^{T_p}
    p^*\mathbb H_YD
    &\longrightarrow
    p_\sharp p^*\mathbb H_YD
    \\
    &\longrightarrow
    \mathbb H_YD.
\end{aligned}
\]

Under the canonical equivalences
\[
    p_!\xrightarrow{\simeq}p_\sharp\Sigma^{-T_p},
    \qquad
    \Sigma^{T_p}p^*\xrightarrow{\simeq}p^!,
\]
Lemma~\ref{lem:dr-int-purity-counit} identifies this composite
precisely with the exceptional counit
\[
    p_!p^!\mathbb H_YD
    \longrightarrow
    \mathbb H_YD.
\]
This proves the trace statement.

Now let
\[
    \omega:
    K\longrightarrow\widetilde p^{\,!}D.
\]
Its differential adjunct is
\[
    \widetilde p_!K
    \xrightarrow{\widetilde p_!\omega}
    \widetilde p_!\widetilde p^{\,!}D
    \xrightarrow{\Trdiff_{p,D}}
    D.
\]
Homotopification preserves composition.  Naturality of
\[
    \mathbb H_Y\widetilde p_!
    \simeq
    p_!\mathbb H_X
\]
identifies the first arrow with
\[
    p_!\mathbb H_XK
    \xrightarrow{
        p_!(\mathbb H_X\omega)
    }
    p_!p^!\mathbb H_YD,
\]
where
\[
    \mathbb H_X\widetilde p^{\,!}
    \simeq
    p^!\mathbb H_Y
\]
is the direct Thom-pullback-purity comparison.  The first part of the
proof identifies the second arrow with the exceptional counit.
Their composite is therefore exactly the motivic adjunct of
\[
    \mathbb H_X\omega:
    \mathbb H_XK
    \longrightarrow
    p^!\mathbb H_YD.
\]
\end{proof}

\begin{remark}[Primary versus secondary integration]
\label{rem:dr-int-primary-secondary}
The theorem gives the primary relation required in this manuscript:
\[
    \mathbb H_Y\widetilde p_!
    \simeq
    p_!\mathbb H_X.
\]
Thus the motivic exceptional integral is the homotopy-invariant shadow of the
differential exceptional integral.  The BNV interval integral considered
below is a different, secondary operation: it is a normalized primitive of an
endpoint difference in the deformation fibre.
\end{remark}

\subsection{The de Rham stage traces}

\begin{definition}[Differential dualizing de Rham stages]
\label{def:dr-int-dualizing-dr-stages}
For \(q\geq0\), define
\[
\boxed{
    \widehat{\mathscr F}_{\dR,X/Y}^{!,q}
    :=
    \widetilde p^{\,!}
    \widehat{\mathscr F}_{\dR,Y}^{q},
}
\]
\[
\boxed{
    \widehat{\mathscr\Omega}_{\inf,X/Y}^{!,q}
    :=
    \widetilde p^{\,!}
    \widehat{\mathscr\Omega}_{\inf,Y}^{q},
    \qquad
    \widehat{\mathscr\Omega}_{\inf,X/Y}^{!,q,\cl}
    :=
    \widetilde p^{\,!}
    \widehat{\mathscr\Omega}_{\inf,Y}^{q,\cl}.
}
\]
\end{definition}

\begin{proposition}[Exactness of the dualizing de Rham stage system]
\label{prop:dr-int-dualizing-stage-exactness}
The adjacent-stage cofiber sequences pull back to canonical cofiber sequences
\[
\boxed{
    \widehat{\mathscr F}_{\dR,X/Y}^{!,q+1}
    \longrightarrow
    \widehat{\mathscr F}_{\dR,X/Y}^{!,q}
    \longrightarrow
    \widehat{\mathscr\Omega}_{\inf,X/Y}^{!,q}[-q].
}
\]
The coherent infinitesimal differential, its specified nullhomotopies and
higher coherences, and the exact-cochain classifiers are carried to the
differential dualizing coefficients.
\end{proposition}

\begin{proof}
The functor
\[
    \widetilde p^{\,!}
    =
    \Sigma^{T_p,0}p^{*,0}
\]
is exact.  Apply it to the finite exact diagrams already transported to
\(\mathcal R_Y^0\) by
Proposition~\ref{prop:monadic-de-rham-exact-thom-realization}.
\end{proof}

\begin{definition}[Differential de Rham stage trace]
\label{def:dr-int-stage-trace}
For every \(q\geq0\), define
\[
\boxed{
    \Trdiff_{p,\dR,q}:
    \widetilde p_!
    \widehat{\mathscr F}_{\dR,X/Y}^{!,q}
    \longrightarrow
    \widehat{\mathscr F}_{\dR,Y}^{q}
}
\]
to be the differential counit.  Define similarly
\[
    \Trdiff_{p,\Omega,q}:
    \widetilde p_!
    \widehat{\mathscr\Omega}_{\inf,X/Y}^{!,q}
    \longrightarrow
    \widehat{\mathscr\Omega}_{\inf,Y}^{q}
\]
and
\[
    \Trdiff_{p,\Omega^{\cl},q}:
    \widetilde p_!
    \widehat{\mathscr\Omega}_{\inf,X/Y}^{!,q,\cl}
    \longrightarrow
    \widehat{\mathscr\Omega}_{\inf,Y}^{q,\cl}.
\]
\end{definition}

\begin{theorem}[The de Rham stage trace is a coherent cochain morphism]
\label{thm:dr-int-coherent-cochain-trace}
The naturality of the differential counit gives a morphism of coherent
nonnegative cochain objects
\[
\boxed{
    \widetilde p_!\widetilde p^{\,!}
    \widehat{\mathscr\Omega}_{\inf,Y}^{\bullet}
    \longrightarrow
    \widehat{\mathscr\Omega}_{\inf,Y}^{\bullet}.
}
\]
In particular, for every \(q\geq0\) there is a canonically specified
commutative square
\[
\begin{tikzcd}[column sep=huge,row sep=large]
    \widetilde p_!\widetilde p^{\,!}
    \widehat{\mathscr\Omega}_{\inf,Y}^{q}
        \ar[r,"\Trdiff_{p,\Omega,q}"]
        \ar[d,"\widetilde p_!\widetilde p^{\,!}d_{\inf}"']
    &
    \widehat{\mathscr\Omega}_{\inf,Y}^{q}
        \ar[d,"d_{\inf}"]
    \\
    \widetilde p_!\widetilde p^{\,!}
    \widehat{\mathscr\Omega}_{\inf,Y}^{q+1}
        \ar[r,"\Trdiff_{p,\Omega,q+1}"']
    &
    \widehat{\mathscr\Omega}_{\inf,Y}^{q+1}.
\end{tikzcd}
\]
These squares preserve the specified nullhomotopies of
\(d_{\inf}^{q+1}d_{\inf}^{q}\) and all higher coherent relations.  The same
assertion holds for the exact-cochain classifier
\[
    d_{\inf}^{\cl}:
    \widehat{\mathscr\Omega}_{\inf,Y}^{q-1}
    \longrightarrow
    \widehat{\mathscr\Omega}_{\inf,Y}^{q,\cl}.
\]
\end{theorem}

\begin{proof}
Let
\[
    T_p^\partial
    :=
    \widetilde p_!\widetilde p^{\,!}:
    \mathcal R_Y^0\longrightarrow\mathcal R_Y^0.
\]
This is an exact endofunctor, and the differential exceptional counit is a
natural transformation
\[
    \varepsilon_p^\partial:
    T_p^\partial\longrightarrow\id.
\]

First apply this natural transformation to the realized de Rham stage
system.  For every \(q\) it gives a commutative square
\[
\begin{tikzcd}[column sep=huge,row sep=large]
    T_p^\partial\widehat{\mathscr F}_{\dR,Y}^{q+1}
        \ar[r]
        \ar[d,"\varepsilon_p^\partial"']
    &
    T_p^\partial\widehat{\mathscr F}_{\dR,Y}^{q}
        \ar[d,"\varepsilon_p^\partial"]
    \\
    \widehat{\mathscr F}_{\dR,Y}^{q+1}
        \ar[r]
    &
    \widehat{\mathscr F}_{\dR,Y}^{q}.
\end{tikzcd}
\]
Because \(T_p^\partial\) is exact, applying it to the adjacent-stage
cofiber sequence gives
\[
    T_p^\partial\widehat{\mathscr F}_{\dR,Y}^{q+1}
    \longrightarrow
    T_p^\partial\widehat{\mathscr F}_{\dR,Y}^{q}
    \longrightarrow
    T_p^\partial
    \widehat{\mathscr\Omega}_{\inf,Y}^{q}[-q].
\]
Naturality of the counit gives a morphism from this cofiber sequence to
\[
    \widehat{\mathscr F}_{\dR,Y}^{q+1}
    \longrightarrow
    \widehat{\mathscr F}_{\dR,Y}^{q}
    \longrightarrow
    \widehat{\mathscr\Omega}_{\inf,Y}^{q}[-q].
\]
Thus the traces on the normalized coefficients are the maps induced on
associated graded by a morphism of the entire stage system.

Equivalently, apply
\[
    \Fun_*
    \left(
        \Xi_{\geq0}^{\op},-
    \right)
\]
to the natural transformation
\(\varepsilon_p^\partial\).  Since an exact functor preserves the zero
object and all finite cofiber diagrams, this gives an actual morphism of
coherent nonnegative cochain objects
\[
\boxed{
    \widetilde p_!\widetilde p^{\,!}
    \widehat{\mathscr\Omega}_{\inf,Y}^{\bullet}
    \longrightarrow
    \widehat{\mathscr\Omega}_{\inf,Y}^{\bullet}.
}
\]
Consequently it preserves the first cochain structure maps, their
specified nullhomotopies, and every higher simplex of the coherent
cochain diagram.  This gives the stated square for \(d_{\inf}\), with
its specified higher coherence.

It remains to verify the exact-cochain classifier.  Recall that, under
the canonical identification
\[
    \widehat{\mathscr\Omega}_{\inf,Y}^{q,\cl}
    \simeq
    \widehat{\mathscr F}_{\dR,Y}^{q}[q],
\]
the classifier
\[
    d_{\inf}^{\cl}:
    \widehat{\mathscr\Omega}_{\inf,Y}^{q-1}
    \longrightarrow
    \widehat{\mathscr\Omega}_{\inf,Y}^{q,\cl}
\]
is obtained from the connecting morphism of the adjacent-stage cofiber
sequence
\[
    \widehat{\mathscr F}_{\dR,Y}^{q}
    \longrightarrow
    \widehat{\mathscr F}_{\dR,Y}^{q-1}
    \longrightarrow
    \widehat{\mathscr\Omega}_{\inf,Y}^{q-1}[1-q].
\]
The morphism of cofiber sequences constructed above induces a
commutative square on their connecting morphisms, because connecting
morphisms in a stable \(\infty\)-category are natural in morphisms of
cofiber sequences.  After shifting by \(q-1\), this is precisely the
compatibility square for \(d_{\inf}^{\cl}\).

Thus the differential exceptional trace preserves the coherent
infinitesimal differential, the exact-cochain classifier, the specified
nullhomotopies of \(d_{\inf}^2\), and all higher coherent relations.
Only the already constructed stage objects and finite exact diagrams
have been used; no preservation of the sequential inverse limit defining
complete totalization is required.
\end{proof}

\begin{corollary}[Motivic shadow of the de Rham stage trace]
\label{cor:dr-int-stage-trace-motivic-shadow}
For every \(q\geq0\),
\[
\boxed{
    \mathbb H_Y
    \left(
        \Trdiff_{p,\dR,q}
    \right)
    \simeq
    p_!p^!
    \mathbb H_Y\widehat{\mathscr F}_{\dR,Y}^{q}
    \longrightarrow
    \mathbb H_Y\widehat{\mathscr F}_{\dR,Y}^{q}.
}
\]
The analogous statement holds for the normalized and coherently closed
cochain coefficients.
\end{corollary}

\begin{proof}
Apply Theorem~\ref{thm:dr-int-motivic-shadow} coefficientwise.
\end{proof}

\begin{remark}[No unproved degree shift]
\label{rem:dr-int-no-unproved-degree-shift}
The differential exceptional source contains the Thom twist
\[
    \Sigma^{T_p,0}.
\]
No ordinary numerical degree shift and no trivialization of this twist is
inserted unless a separate orientation theorem supplies one.
\end{remark}

\section{The BNV obstruction and universal interval integration}
\label{sec:dr-int-bnv}

The affine-line cylinder and universal BNV interval integral on
\(\mathcal R_S^0\) are those of
Chapter~\ref{chap:brave-new-infinitesimal-geometry}:
\[
    \mathbb I_S=q_{S,*}^0q_S^{*,0},
\]
\[
\boxed{
    \IntBNV_{\mathbb A^1,S}:
    \mathbb I_S\mathcal Z_S
    \longrightarrow
    \mathcal A_S.
}
\]

\subsection{The evaluated obstruction}

\begin{definition}[Evaluated BNV obstruction]
\label{def:dr-int-evaluated-bnv-obstruction}
For \(E\in\mathcal R_X^0\), define
\[
\boxed{
    \ObsBNV(p;E)
    :=
    \cofib
    \left(
        \omega_p(E):
        \mathcal Z_Y\widetilde p_!E
        \longrightarrow
        \mathcal Z_Y\widetilde p_!\mathcal Z_XE
    \right).
}
\]
\end{definition}

\begin{proposition}[Evaluation of the global obstruction]
\label{prop:dr-int-evaluated-bnv-obstruction}
There is a canonical equivalence
\[
\boxed{
    \ObsBNV(p;E)
    \simeq
    \Sigma\,
    \mathcal Z_Y
    \widetilde p_!
    \mathcal S_XE.
}
\]
If the global cycle defect
\[
    \omega_p:
    \mathcal Z_Y\widetilde p_!
    \longrightarrow
    \mathcal Z_Y\widetilde p_!\mathcal Z_X
\]
is an equivalence, then
\[
    \ObsBNV(p;E)\simeq0
\]
for every coefficient \(E\).

No converse from the vanishing of a single evaluated obstruction to the global
equivalence of \(\omega_p\) is asserted.
\end{proposition}

\begin{proof}
Evaluate the functorial obstruction equivalence
\[
    \ObsBNV(p)
    \simeq
    \Sigma\mathcal Z_Y\widetilde p_!\mathcal S_X
\]
constructed in Chapter~\ref{chap:brave-new-infinitesimal-geometry}.
\end{proof}

\subsection{BNV--Umkehr compatibility}

Assume for the remainder of this subsection that the global cycle defect
\(\omega_p\) is an equivalence.  Then
Chapter~\ref{chap:brave-new-infinitesimal-geometry} constructs the cycle and
deformation exchanges and the normalized cylinder exchange.  We write
\[
    \theta_{\mathcal Z}^p:=\theta_Z,
    \qquad
    \theta_{\mathcal A}^p:=\theta_A,
    \qquad
    \theta_{\mathbb I}^p:=\theta_I^p
\]
for those previously constructed transformations.  Thus
\[
    \theta_{\mathcal Z}^p:
    \widetilde p_!\mathcal Z_X
    \longrightarrow
    \mathcal Z_Y\widetilde p_!,
\]
\[
    \theta_{\mathcal A}^p:
    \widetilde p_!\mathcal A_X
    \xrightarrow{\simeq}
    \mathcal A_Y\widetilde p_!,
\]
\[
    \theta_{\mathbb I}^p:
    \widetilde p_!\mathbb I_X
    \longrightarrow
    \mathbb I_Y\widetilde p_!.
\]
Put
\[
    \theta_{\mathbb I\mathcal Z}^p
    :=
    \mathbb I_Y(\theta_{\mathcal Z}^p)
    \circ
    \theta_{\mathbb I}^p\mathcal Z_X.
\]

\begin{theorem}[BNV interval integration is natural under differential Umkehr]
\label{thm:dr-int-bnv-umkehr}
There is a canonical commutative square
\[
\begin{tikzcd}[column sep=huge,row sep=large]
    \widetilde{p}_!\mathbb{I}_X\mathcal{Z}_X
        \ar[r,"{\widetilde{p}_!\IntBNV_{\mathbb{A}^1,X}}"]
        \ar[d,"{\theta_{\mathbb{I}\mathcal{Z}}^p}"']
    &
    \widetilde{p}_!\mathcal{A}_X
        \ar[d,"{\theta_{\mathcal{A}}^p}"']
    \\
    \mathbb{I}_Y\mathcal{Z}_Y\widetilde{p}_!
        \ar[r,"{\IntBNV_{\mathbb{A}^1,Y}}"']
    &
    \mathcal{A}_Y\widetilde{p}_!
\end{tikzcd}
\]
\end{theorem}

\begin{proof}
This is the BNV integration naturality theorem of
Chapter~\ref{chap:brave-new-infinitesimal-geometry} applied to the constructed
smooth Thom-stable lift and its constructed normalized cylinder exchange.
\end{proof}

\begin{theorem}[The differential exceptional trace preserves the BNV primitive]
\label{thm:dr-int-trace-bnv}
Assume \(\omega_p\) is an equivalence.  Let \(D\in\mathcal R_Y^0\) and put
\[
    E_D:=\widetilde p^{\,!}D.
\]
Define
\[
\begin{aligned}
    \lambda_{p,D}:\,
    \widetilde p_!\mathbb I_X\mathcal Z_X(E_D)
    &\xrightarrow{\theta_{\mathbb I\mathcal Z}^p}
    \mathbb I_Y\mathcal Z_Y(\widetilde p_!E_D)
    \\
    &\xrightarrow{
        \mathbb I_Y\mathcal Z_Y(\Trdiff_{p,D})
    }
    \mathbb I_Y\mathcal Z_Y(D),
\end{aligned}
\]
and
\[
\begin{aligned}
    \rho_{p,D}:\,
    \widetilde p_!\mathcal A_X(E_D)
    &\xrightarrow{\theta_{\mathcal A}^p}
    \mathcal A_Y(\widetilde p_!E_D)
    \\
    &\xrightarrow{
        \mathcal A_Y(\Trdiff_{p,D})
    }
    \mathcal A_Y(D).
\end{aligned}
\]
Then
\[
\begin{tikzcd}[column sep=huge,row sep=large]
    \widetilde{p}_!
    \mathbb{I}_X\mathcal{Z}_X
    \bigl(\widetilde{p}^{\,!}D\bigr)
        \ar[r,"{\widetilde{p}_!\IntBNV_{\mathbb{A}^1,X}}"]
        \ar[d,"{\lambda_{p,D}}"']
    &
    \widetilde{p}_!
    \mathcal{A}_X
    \bigl(\widetilde{p}^{\,!}D\bigr)
        \ar[d,"{\rho_{p,D}}"]
    \\
    \mathbb{I}_Y\mathcal{Z}_Y(D)
        \ar[r,"{\IntBNV_{\mathbb{A}^1,Y}}"']
    &
    \mathcal{A}_Y(D)
\end{tikzcd}
\]
commutes through a canonical specified homotopy.

For the same coefficient \(D\), the homotopification of the trace in this
square is
\[
    p_!p^!\mathbb H_YD
    \longrightarrow
    \mathbb H_YD.
\]
Thus the motivic exceptional trace is the primary shadow of the differential
trace, while universal BNV interval integration is a secondary operation
preserved by that same trace.
\end{theorem}

\begin{proof}
Theorem~\ref{thm:dr-int-bnv-umkehr} gives the square before postcomposition
with the differential counit.  Universal BNV integration is natural in
morphisms of coefficient objects.  Apply this naturality to
\[
    \Trdiff_{p,D}:
    \widetilde p_!\widetilde p^{\,!}D
    \longrightarrow
    D
\]
and paste the two squares.  The motivic statement is
Theorem~\ref{thm:dr-int-motivic-shadow}.
\end{proof}

\begin{corollary}[BNV compatibility for the realized de Rham stages]
\label{cor:dr-int-bnv-dr-stages}
Assume \(\omega_p\) is an equivalence.  Theorem~\ref{thm:dr-int-trace-bnv}
applies to
\[
    D=\widehat{\mathscr{DR}}_Y,
    \qquad
    D=\widehat{\mathscr F}_{\dR,Y}^{q},
    \qquad
    D=\widehat{\mathscr\Omega}_{\inf,Y}^{q},
    \qquad
    D=\widehat{\mathscr\Omega}_{\inf,Y}^{q,\cl}.
\]
The resulting BNV comparison squares are compatible with totalization-stage
transitions, the coherent infinitesimal differential, the exact-cochain
classifier, and all higher coherent relations.
\end{corollary}

\begin{proof}
Apply the theorem coefficientwise and use the naturality of all BNV functors
and of the differential counit in the coefficient object.
\end{proof}

\subsection{The normalized Stokes factorization}

Put
\[
    \delta_S:=1^*-0^*:
    \mathbb I_S\longrightarrow\id.
\]
Write
\[
    a_S:\mathcal A_S\longrightarrow\id,
    \qquad
    R_S:\id\longrightarrow\mathcal Z_S
\]
for the deformation inclusion and cycle quotient of
Chapter~\ref{chap:brave-new-infinitesimal-geometry}.  The universal BNV
theorem gives the normalized factorization
\[
\boxed{
    \delta_S
    \simeq
    a_S
    \circ
    \IntBNV_{\mathbb A^1,S}
    \circ
    \mathbb I_SR_S.
}
\]

\begin{theorem}[Differential Umkehr preserves the BNV Stokes factorization]
\label{thm:dr-int-stokes}
Assume \(\omega_p\) is an equivalence.  For every \(D\in\mathcal R_Y^0\),
applying \(\widetilde p_!\) to the normalized endpoint-difference
factorization of \(\widetilde p^{\,!}D\), transporting through the constructed
cycle, deformation, and cylinder exchanges, and postcomposing with
\(\Trdiff_{p,D}\) gives the normalized endpoint-difference factorization of
\(D\).

Equivalently, the differential exceptional trace preserves the universal BNV
Stokes primitive.
\end{theorem}

\begin{proof}
The normalized cylinder exchange is compatible with both endpoint maps, with
the projection transformation, and with their specified retractions.
Therefore
\[
    \delta_Y\widetilde p_!
    \circ
    \theta_{\mathbb I}^p
    \simeq
    \widetilde p_!\delta_X.
\]
The universal BNV integral is characterized by its normalized factorization of
this endpoint difference through the cycle cofiber and deformation fibre.
Theorem~\ref{thm:dr-int-trace-bnv} identifies the two factorizations after
applying the differential trace.
\end{proof}

\begin{remark}[BNV primitives and intrinsic de Rham primitives remain distinct]
\label{rem:dr-int-two-primitives}
The BNV primitive of
Theorem~\ref{thm:dr-int-stokes} lies in the deformation fibre
\(\mathcal A\).  The intrinsic de Rham primitive spaces of
Definition~\ref{def:modal-form-primitive-spectrum} are homotopy fibres of the
exact-cochain classifier
\[
    d_{\inf}^{\cl}.
\]
No canonical identification of these two kinds of primitive is asserted.
Their interaction requires a concrete coefficient map between the relevant
BNV sector and the realized intrinsic exact-cochain object.
\end{remark}

\section{\texorpdfstring{\(\acute{e}\)tale}{etale} intrinsic de Rham integration}
\label{sec:dr-int-etale}
\label{sec:dr-int-finite-etale}

For represented spectral \(\acute{e}\)tale morphisms of finite presentation,
the intrinsic de Rham coefficient exchange is an equivalence by
Theorem~\ref{thm:monadic-de-rham-smooth-site-etale-exchange}.  If the morphism
is separated, its tangent class vanishes, so the intrinsic source coefficient
is the differential exceptional pullback of the target coefficient and the
exceptional counit defines an intrinsic differential de Rham pushforward.
This applies, in particular, to every finite \(\acute{e}\)tale morphism.  The
additional Nisnevich-locally split hypothesis enters only when we invoke the
vanishing theorem for the BNV obstruction from the preceding chapter.

\subsection{The split exchange}

\begin{proposition}[Split finite \(\acute{e}\)tale exchange for the de Rham coefficient system]
\label{prop:dr-int-split-etale-exchange}
Let
\[
    p:
    \coprod_{i=1}^{d}Y
    \longrightarrow
    Y
\]
be the fold map.  For each of the smooth-site coefficient systems
\[
    \mathscr{DR},\qquad
    \mathscr F_{\dR}^{q},\qquad
    \mathscr\Omega_{\inf}^{q},\qquad
    \mathscr\Omega_{\inf}^{q,\cl},
\]
the canonical base-change exchange
\[
    p^*_{\Nis}\mathscr F_Y
    \longrightarrow
    \mathscr F_{\coprod_iY}
\]
is an equivalence.

After Thom-stable realization,
\[
\boxed{
    p^{*,0}\widehat{\mathscr F}_Y
    \xrightarrow{\simeq}
    \widehat{\mathscr F}_{\coprod_iY}.
}
\]
Under the product decomposition of the source coefficient category, the
right side is the \(d\)-tuple
\[
    (\widehat{\mathscr F}_Y,\ldots,\widehat{\mathscr F}_Y).
\]
\end{proposition}

\begin{proof}
A smooth finite-presentation spectral scheme over
\(\coprod_iY\) decomposes uniquely as a finite disjoint union of its
open-and-closed inverse images over the components.  The finite-coproduct
compatibility of
Proposition~\ref{prop:monadic-de-rham-finite-coproducts} identifies every
intrinsic de Rham construction componentwise: effective quotients, relation
degrees, cosimplicial cochains, normalized cochains, totalization stages, and
coherently closed cochains all become finite products of the corresponding
objects over the components.

The inverse-image functor along the fold map is the diagonal under the same
component decomposition of stable Nisnevich sheaves.  Hence the canonical
smooth-site exchange is componentwise the identity.  Applying
\(\RealDiff\) and the pullback compatibility of
Proposition~\ref{prop:monadic-de-rham-realization-smooth-direct-image} gives
the Thom-stable statement.
\end{proof}

\subsection{\texorpdfstring{\(\acute{e}\)tale}{Etale} exchange}

\begin{theorem}[\(\acute{e}\)tale Cartesianity of the intrinsic de Rham coefficient systems]
\label{thm:dr-int-etale-exchange}
\label{thm:dr-int-locally-split-etale-exchange}
Let
\[
    p:X\longrightarrow Y
\]
be a represented spectral \(\acute{e}\)tale morphism of finite presentation
between bases in the standing motivic scope.  Then, for each smooth-site de
Rham coefficient system
\[
    \mathscr{DR},\qquad
    \mathscr F_{\dR}^{q},\qquad
    \mathscr\Omega_{\inf}^{q},\qquad
    \mathscr\Omega_{\inf}^{q,\cl},
\]
the canonical exchange is an equivalence:
\[
\boxed{
    p^*_{\Nis}\mathscr F_Y
    \xrightarrow{\simeq}
    \mathscr F_X.
}
\]
Consequently,
\[
\boxed{
    p^{*,0}\widehat{\mathscr F}_Y
    \xrightarrow{\simeq}
    \widehat{\mathscr F}_X.
}
\]
The equivalences preserve totalization-stage transitions, coherent
infinitesimal differentials, coherently closed cochains, exact-cochain
classifiers, the available products, identities, composition, and all higher
\(\acute{e}\)tale base-change coherences.
\end{theorem}

\begin{proof}
The smooth-site equivalence is
Theorem~\ref{thm:monadic-de-rham-smooth-site-etale-exchange}.  The
pullback compatibility of the previously constructed realization functor gives
\[
\begin{aligned}
    p^{*,0}\widehat{\mathscr F}_Y
    &=
    p^{*,0}\RealDiff_Y\mathscr F_Y
    \\
    &\simeq
    \RealDiff_Xp^*_{\Nis}\mathscr F_Y
    \\
    &\simeq
    \RealDiff_X\mathscr F_X
    \\
    &=
    \widehat{\mathscr F}_X.
\end{aligned}
\]
Every structural compatibility is inherited from the smooth-site
\(\acute{e}\)tale exchange and the coherent pullback compatibility of
\(\RealDiff\).
\end{proof}

\subsection{Intrinsic differential pushforward}

Assume now that
\[
    p:X\longrightarrow Y
\]
is separated, spectral \(\acute{e}\)tale, and of finite presentation.  Since
\(p\) is \(\acute{e}\)tale,
\[
    T_p=0,
\]
hence
\[
    \widetilde p^{\,!}
    =
    p^{*,0}.
\]
For any one of the realized intrinsic de Rham coefficients
\[
    \widehat{\mathscr F}_Y
\]
Theorem~\ref{thm:dr-int-etale-exchange} therefore gives
\[
\boxed{
    \widehat{\mathscr F}_X
    \xrightarrow{\simeq}
    \widetilde p^{\,!}\widehat{\mathscr F}_Y.
}
\]

\begin{definition}[\(\acute{e}\)tale intrinsic differential de Rham pushforward]
\label{def:dr-int-etale-push}
\label{def:dr-int-finite-etale-push}
Define
\[
\boxed{
    p_{\dR,!}^{\partial}:
    \widetilde p_!
    \widehat{\mathscr F}_X
    \longrightarrow
    \widehat{\mathscr F}_Y
}
\]
as the composite
\[
    \widetilde p_!\widehat{\mathscr F}_X
    \xrightarrow{\simeq}
    \widetilde p_!\widetilde p^{\,!}\widehat{\mathscr F}_Y
    \xrightarrow{\Trdiff_{p,\widehat{\mathscr F}_Y}}
    \widehat{\mathscr F}_Y.
\]
This definition applies simultaneously to the total de Rham coefficient,
every totalization stage, every normalized cochain coefficient, and every
coherently closed cochain coefficient.
\end{definition}

\begin{theorem}[\(\acute{e}\)tale de Rham pushforward, motivic shadow, and finite-\(\acute{e}\)tale BNV integration]
\label{thm:dr-int-finite-etale-main}
Let
\[
    p:X\longrightarrow Y
\]
be separated, spectral \(\acute{e}\)tale, and of finite presentation.  Then
the intrinsic differential de Rham pushforwards of
Definition~\ref{def:dr-int-etale-push} satisfy:

\begin{enumerate}
\item they preserve the totalization-stage transitions, normalized cochains,
coherently closed cochains, exact-cochain classifiers, coherent
\(d_{\inf}\), its specified square-zero homotopies, and all higher cochain
coherences;

\item after motivic homotopification they are the corresponding exceptional
traces
\[
\boxed{
    p_!p^!
    \mathbb H_Y\widehat{\mathscr F}_Y
    \longrightarrow
    \mathbb H_Y\widehat{\mathscr F}_Y.
}
\]
\end{enumerate}

If, in addition, \(p\) is finite \(\acute{e}\)tale and
Nisnevich-locally split in the sense of
Chapter~\ref{chap:brave-new-infinitesimal-geometry}, then its BNV cycle defect
is an equivalence and the same intrinsic de Rham pushforward preserves the
universal BNV interval primitive through the canonical trace-level square of
Theorem~\ref{thm:dr-int-trace-bnv}.

Thus the primary intrinsic differential pushforward and its motivic shadow are
unconditional throughout the separated \(\acute{e}\)tale finite-presentation
class, while the secondary BNV compatibility is unconditional in the verified
Nisnevich-locally split finite \(\acute{e}\)tale class.
\end{theorem}

\begin{proof}
The coefficient identification is
Theorem~\ref{thm:dr-int-etale-exchange}.  The first assertion then follows
from Theorem~\ref{thm:dr-int-coherent-cochain-trace}, since the intrinsic
source coefficient has been identified with the differential exceptional
pullback of the target.  The motivic shadow is
Theorem~\ref{thm:dr-int-motivic-shadow}.

If \(p\) is finite \(\acute{e}\)tale and Nisnevich-locally split, Chapter
\ref{chap:brave-new-infinitesimal-geometry} proves that the global BNV cycle
defect of \(\widetilde p_!\) is an equivalence.  Applying
Theorem~\ref{thm:dr-int-trace-bnv} gives the BNV interval compatibility.
\end{proof}


\begin{theorem}[Explicit split finite \(\acute{e}\)tale formula]
\label{thm:dr-int-split-etale-codiagonal}
For the fold map
\[
    p:
    \coprod_{i=1}^{d}Y
    \longrightarrow
    Y,
\]
the source category decomposes as
\[
    \mathcal R_{\coprod_iY}^0
    \simeq
    \prod_{i=1}^{d}\mathcal R_Y^0.
\]
For every realized de Rham coefficient \(D_Y\),
\[
    \widetilde p^{\,!}D_Y
    \simeq
    (D_Y,\ldots,D_Y),
\]
\[
    \widetilde p_!(E_1,\ldots,E_d)
    \simeq
    \bigoplus_{i=1}^{d}E_i,
\]
and the intrinsic differential de Rham trace is the codiagonal
\[
\boxed{
    D_Y^{\oplus d}
    \longrightarrow
    D_Y.
}
\]
The BNV deformation and cycle sectors and the affine-line cylinder are
componentwise, so the BNV integration square is the codiagonal of the \(d\)
componentwise BNV integration squares.  Its motivic shadow is the split finite
\(\acute{e}\)tale motivic trace.
\end{theorem}

\begin{proof}
The product decomposition and direct-sum formula for \(\widetilde p_!\) are
proved in Chapter~\ref{chap:brave-new-infinitesimal-geometry}.  The intrinsic
source coefficient is the diagonal tuple by
Proposition~\ref{prop:dr-int-split-etale-exchange}.  The counit of
\[
    \bigoplus
    \dashv
    \Delta
\]
is the codiagonal.  The BNV and motivic statements follow from the
componentwise description and
Theorem~\ref{thm:dr-int-finite-etale-main}.
\end{proof}

\section{Differential transgression with explicit coefficient traces}
\label{sec:dr-int-transgression}

Consider a correspondence
\[
\begin{tikzcd}[column sep=huge]
    & Z \ar[dl,"e"'] \ar[dr,"\pi"] & \\
    X && B
\end{tikzcd}
\]
in which
\[
    \pi:Z\longrightarrow B
\]
is smooth, separated, and of finite presentation.

\begin{definition}[Differential coefficient trace datum]
\label{def:dr-int-coefficient-trace-datum}
Let
\[
    D_X\in\mathcal R_X^0,
    \qquad
    D_B\in\mathcal R_B^0.
\]
A \emph{differential coefficient trace datum} for the correspondence is a
specified morphism
\[
\boxed{
    \gamma:
    e^{*,0}D_X
    \longrightarrow
    \widetilde\pi^{\,!}D_B
}
\]
in \(\mathcal R_Z^0\).
\end{definition}

\begin{definition}[Differential transgression]
\label{def:dr-int-transgression}
Given \(\gamma\), define
\[
\boxed{
    \Transdiff_{\pi,e,\gamma}:
    \widetilde\pi_!e^{*,0}D_X
    \longrightarrow
    D_B
}
\]
by
\[
    \widetilde\pi_!e^{*,0}D_X
    \xrightarrow{\widetilde\pi_!\gamma}
    \widetilde\pi_!\widetilde\pi^{\,!}D_B
    \xrightarrow{\Trdiff_{\pi,D_B}}
    D_B.
\]
\end{definition}

\begin{proposition}[Motivic shadow of differential transgression]
\label{prop:dr-int-transgression-shadow}
The homotopification of
\(\Transdiff_{\pi,e,\gamma}\) is the motivic correspondence morphism
\[
\boxed{
    \pi_!e^*\mathbb H_XD_X
    \longrightarrow
    \mathbb H_BD_B
}
\]
obtained by adjunction from
\[
    \mathbb H_Z\gamma:
    e^*\mathbb H_XD_X
    \longrightarrow
    \pi^!\mathbb H_BD_B.
\]
\end{proposition}

\begin{proof}
Use the pullback comparison for \(e^{*,0}\), the motivic comparisons for
\(\widetilde\pi_!\) and \(\widetilde\pi^{\,!}\), and
Theorem~\ref{thm:dr-int-motivic-shadow}.
\end{proof}

\begin{proposition}[BNV naturality of differential transgression]
\label{prop:dr-int-transgression-bnv}
Assume the global BNV cycle defect of \(\widetilde\pi_!\) is an equivalence.
Then differential transgression preserves the universal BNV interval primitive
through the canonical pasting of:

\begin{enumerate}
\item the BNV--Umkehr square for \(\widetilde\pi_!\);
\item naturality of BNV integration under
\(\widetilde\pi_!\gamma\);
\item naturality of BNV integration under the differential trace
\(\Trdiff_{\pi,D_B}\).
\end{enumerate}

No additional BNV compatibility datum for \(\gamma\) is required: \(\gamma\)
is an actual morphism in the Thom-stable differential category.
\end{proposition}

\begin{proof}
Apply Theorem~\ref{thm:dr-int-bnv-umkehr} to
\(\widetilde\pi_!\), and paste with the ordinary naturality squares of the
universal BNV integration transformation for the two coefficient morphisms.
\end{proof}

\begin{remark}[No transgression from an unspecified orientation]
\label{rem:dr-int-transgression-no-formal-orientation}
The exceptional adjunction does not by itself construct
\[
    \gamma:
    e^{*,0}D_X
    \longrightarrow
    \widetilde\pi^{\,!}D_B.
\]
For intrinsic de Rham coefficients such a map must come from actual geometric
or coefficient data.  In particular, no blanket orientation, regulator,
differential six-functor axiom, or analytic realization is inserted to force
its existence.
\end{remark}

\begin{proposition}[Transgression of a coherent de Rham stage system]
\label{prop:dr-int-coherent-transgression}
Suppose a coherent family of coefficient trace data
\[
    \gamma_q:
    e^{*,0}\widehat{\mathscr F}_{\dR,X}^{q}
    \longrightarrow
    \widetilde\pi^{\,!}
    \widehat{\mathscr F}_{\dR,B}^{q}
\]
has been constructed and is compatible with the stage transitions.  Then
there are stagewise differential transgressions
\[
    \widetilde\pi_!
    e^{*,0}\widehat{\mathscr F}_{\dR,X}^{q}
    \longrightarrow
    \widehat{\mathscr F}_{\dR,B}^{q}.
\]
If the family is compatible with the adjacent-stage exact diagrams, it induces
transgression maps on normalized and coherently closed cochains which commute
with \(d_{\inf}\), \(d_{\inf}^{\cl}\), their specified nullhomotopies, and all
higher coherent relations.

If the BNV obstruction of \(\widetilde\pi_!\) vanishes, these transgressions
preserve the universal BNV interval primitive.
\end{proposition}

\begin{proof}
Apply the exact functors \(e^{*,0}\) and \(\widetilde\pi_!\) to the coherent
stage diagrams, use the family \(\gamma_q\), and postcompose with the
differential exceptional trace.  Exactness preserves the adjacent finite
cofiber diagrams.  Naturality of the trace preserves the coherent
differential.  The final assertion is
Proposition~\ref{prop:dr-int-transgression-bnv}.
\end{proof}

\begin{remark}[Scope of the integration theorem]
\label{rem:dr-int-integration-scope}
The chapter has therefore proved four statements of increasing specificity.

For every smooth separated finite-presentation \(p\), the differential
exceptional trace exists on differential dualizing coefficients and refines
the motivic exceptional trace.

If the constructed BNV cycle defect of such a \(p\) is an equivalence, that
differential trace preserves universal BNV interval integration.

For every separated spectral \(\acute{e}\)tale morphism of finite
presentation, the intrinsic source de Rham coefficients themselves identify
with the differential exceptional pullbacks.  Hence the intrinsic differential
de Rham pushforward and its motivic exceptional shadow are unconditional in
that class.

For the verified Nisnevich-locally split finite \(\acute{e}\)tale class, the
BNV obstruction also vanishes, so the intrinsic differential de Rham
pushforward preserves universal BNV interval integration unconditionally.
\end{remark}



\section{Structural synthesis}
\label{sec:monadic-de-rham-synthesis}

\begin{theorem}[Brave-new de Rham, differential cochain, and Umkehr package]
\label{thm:monadic-de-rham-package}
Retain the relative de Rham localization, relative de Rham objects, formal infinitesimal geometry,
and differential-cohomological ambient formalism of the preceding chapter.  For every
\(X\to S\) in \(\mathcal X_{\Sp}\), the present chapter then canonically constructs:
\begin{enumerate}
\item the effective image and effective infinitesimal groupoid
\[
    X\twoheadrightarrow Q_{X/S}\hookrightarrow X_{\dR/S},
    \qquad
    R_{X/S,\bullet}
    \simeq
    \check C_\bullet(X\to Q_{X/S}),
    \qquad
    |R_{X/S,\bullet}|\simeq Q_{X/S};
\]
\item for every \(Z\to Y\to S\), the intrinsic relative de Rham completion
\[
\boxed{
    \widehat Y^{\,\dR}_{Z/S}
    =
    Y\times_{Y_{\dR/S}}Z_{\dR/S},
}
\]
with arbitrary base change and canonical factorization of affine finite spectral infinitesimal
thickenings and finite composites of global structured square-zero thickenings, together with the
diagonal-completion identification
\[
\boxed{
    R_{X/S,n}
    \simeq
    \widehat{X^{\times_S(n+1)}}^{\,\dR}_{X/S};
}
\]
\item the parameterized stable categories
\[
    \mathcal E_T=\Sp(\mathcal X_{\Sp,/T})
\]
with \(\Sigma_f\dashv f^*\dashv\Pi_f\), Beck--Chevalley, projection, closed base change, and effective
descent;
\item the additive scalar algebra obtained by canonically equipping the existing brave-new affine line
with its additive group law and stabilizing it,
\[
    \mathbb O_S^{\add}\in\CAlg(\mathcal E_S),
    \qquad
    \Omega^\infty\mathbb O_S^{\add}\simeq\mathbb A^1_S;
\]
\item the geometric cosimplicial commutative algebra
\[
\boxed{
    \underline{\mathbf C}_{\dR}^{\,n}(X/S)
    =
    \Pi_{a_n}a_n^*\mathbb O_S^{\add}
    \simeq
    \underline{\Hom}_{\mathcal E_S}
    \left(
        \Sigma^\infty_{+,S}R_{X/S,n},
        \mathbb O_S^{\add}
    \right);
}
\]
\item the infinitesimal descent monad
\[
    \mathbb M_{X/S}=\Pi_{\epsilon_{X/S}}\epsilon_{X/S}^*
\]
and the monadic presentation
\[
\boxed{
    \underline{\mathbf C}_{\dR}^{\,\bullet}(X/S)
    \simeq
    \Pi_{b_{X/S}}
    \Am_{\mathbb M_{X/S}}^\bullet
    \left(
        b_{X/S}^*\mathbb O_S^{\add}
    \right);
}
\]
\item the full or modal de Rham algebra
\[
\boxed{
    \mathbf{DR}^{\mathrm{full}}_{X/S}
    \simeq
    \Pi_{q_{\dR}}q_{\dR}^*\mathbb O_S^{\add}
    \simeq
    \underline{\Hom}_{\mathcal E_S}
    \left(
        \Sigma^\infty_{+,S}X_{\dR/S},
        \mathbb O_S^{\add}
    \right),
}
\]
and the effective or \v{C}ech de Rham algebra
\[
\boxed{
    \mathbf{DR}^{\mathrm{eff}}_{X/S}
    \simeq
    \Pi_{b_{X/S}}b_{X/S}^*\mathbb O_S^{\add}
    \simeq
    \underline{\Hom}_{\mathcal E_S}
    \left(
        \Sigma^\infty_{+,S}Q_{X/S},
        \mathbb O_S^{\add}
    \right),
}
\]
together with the canonical comparison
\[
\boxed{
    \rho_{X/S}:
    \mathbf{DR}^{\mathrm{full}}_{X/S}
    \longrightarrow
    \mathbf{DR}^{\mathrm{eff}}_{X/S};
}
\]
for represented smooth spectral morphisms, smooth effectivity identifies \(Q_{X/S}\) with
\(X_{\dR/S}\) and makes \(\rho_{X/S}\) an equivalence;
\item the complete \v{C}ech-totalization filtration
\[
    F_{\mathrm{Tot}}^q\mathbf{DR}_{X/S}
    =
    \fib
    \left(
        \mathbf{DR}^{\mathrm{eff}}_{X/S}
        \to
        \mathbf{DR}_{X/S}^{<q}
    \right),
\]
with normalized infinitesimal de Rham cochains
\[
\boxed{
    \gr_{\mathrm{Tot}}^q\mathbf{DR}_{X/S}
    \simeq
    \underline{\OmegaInf}^{\,q}(X/S)[-q]
}
\]
and skeletal presentation
\[
\boxed{
    F_{\mathrm{Tot}}^q\mathbf{DR}_{X/S}
    \simeq
    \underline{\Hom}_{\mathcal E_S}
    \left(
        A_{\geq q}(X/S),
        \mathbb O_S^{\add}
    \right);
}
\]
\item the coherent infinitesimal differential
\[
    d_{\inf}:
    \underline{\OmegaInf}^{\,q}(X/S)
    \longrightarrow
    \underline{\OmegaInf}^{\,q+1}(X/S)
\]
with its specified coherent nullhomotopy of \(d_{\inf}^2\) and all higher relations, together with
the natural ordered Alexander--Whitney maps whose images in the homotopy category are associative and
unital and satisfy the graded Leibniz identity;
\item for represented \(X\to S\), the first structured infinitesimal neighbourhood of the diagonal
\[
    X\hookrightarrow D^{(1)}_{X/S}\longrightarrow X\times_SX,
\]
the stabilized cotangent coefficient
\[
    \mathbb L^{\mathrm{st}}_{X/S}
\]
whose affine values are those of \(L_{X/S}\), and the first-order linearization
\[
\boxed{
    \lambda^1_{X/S}:
    \underline{\OmegaInf}^{\,1}(X/S)
    \longrightarrow
    \Pi_p\mathbb L^{\mathrm{st}}_{X/S}
}
\]
with
\[
\boxed{
    \lambda^1_{X/S}d_{\inf}^0
    \simeq
    d^{\mathrm{st}}_{X/S};
}
\]
for the discrete affine line over \(\mathbb Q\), the class \((t_1-t_0)^2\) gives a nonzero element in
the kernel of \(\pi_0\Gamma_S^{\mathrm{st}}(\lambda^1_{X/S})\), so the canonical first-order
linearization is not an equivalence;
\item the coherently closed normalized-cochain spectra
\[
    \underline{\OmegaInf}^{\,q,\cl}(X/S;n)
    =
    F_{\mathrm{Tot}}^q\mathbf{DR}_{X/S}[n+q],
\]
the exact-cochain classifier, primitive spaces, and the additive-delooping presentation
\[
\boxed{
    \Omega^\infty\OmegaInf^{q,\cl}(X/S)
    \simeq
    \Map_{\mathcal X_{\Sp,/S,*}}
    \left(
        \mathfrak I^{\langle q\rangle}_{X/S},
        \Omega^\infty\Sigma^q\mathbb O_S^{\add}
    \right);
}
\]
\item the moduli objects
\[
    \boldsymbol{\Omega}_{\mathrm{inf},S}^{q},
    \qquad
    \boldsymbol{\Omega}_{\mathrm{inf},S}^{q,\cl}
    \in
    \mathcal X_{\Sp,/S},
\]
their tautological universal affine-test points, and the moduli infinitesimal differential;
\item contravariant functoriality, arbitrary parameterized base change, spectral Nisnevich descent for
the effective cochains and their totalization, spectral Nisnevich descent for full de Rham functions,
external products, formal-\(\acute{e}\)tale Cartesianity, formal-\(\acute{e}\)tale invariance of
intrinsic completions, and invariance of the scalar-valued intrinsic de Rham package under insertion
of a formally \(\acute{e}\)tale intermediate base;
\item for represented smooth \(X\to S\), Zariski-local linear models
\[
\boxed{
    R_{U/V,n}
    \simeq
    \widehat{V_U(L_{U/V}^{\oplus n})}^{\,\dR}_{U/V},
}
\]
obtained from the smooth-affine tubular linearization of the preceding chapter;
\item the realization of the effective de Rham, totalization, normalized-cochain, and coherently
closed-cochain assignments as objects of the stable smooth-site Nisnevich
coefficient category \(\mathcal C_S\), followed by the Thom-stable differential
realization
\[
\boxed{
    \RealDiff_S
    =
    L_{\mathrm{Th},S}F_{0,S}:
    \mathcal C_S\longrightarrow\mathcal R_S^0;
}
\]
the functor \(\RealDiff_S\) is exact, colimit-preserving and strong symmetric
monoidal, but no full-faithfulness statement is made; on smooth represented
tests the full and effective de Rham coefficient objects coincide;

\item the stabilized evaluation formula
\[
\boxed{
    \map_{\mathcal R_S^0}
    \left(
        \RealDiff_SA,E
    \right)
    \simeq
    \map_{\mathcal C_S}
    \left(
        A,(j_{\mathrm{Th},S}E)_0
    \right),
}
\]
and for smooth finite-presentation \(p:X\to Y\), the coherent smooth exchange
\[
\boxed{
    p_\sharp^0\RealDiff_X
    \simeq
    \RealDiff_Yp_{\sharp,\Nis},
    \qquad
    p^{*,0}\RealDiff_Y
    \simeq
    \RealDiff_Xp^*_{\Nis};
}
\]

\item for every smooth separated finite-presentation morphism
\(p:X\to Y\), the differential exceptional adjunction of the preceding
chapter acts on the realized de Rham coefficients:
\[
\boxed{
    \widetilde p_!
    =
    p_\sharp^0\Sigma^{-T_p,0}
    \dashv
    \Sigma^{T_p,0}p^{*,0}
    =
    \widetilde p^{\,!};
}
\]
its counit
\[
    \Trdiff_{p,D}:
    \widetilde p_!\widetilde p^{\,!}D
    \longrightarrow
    D
\]
is the differential exceptional trace;

\item the motivic shadow of the differential exceptional trace is exactly
\[
\boxed{
    p_!p^!\mathbb H_YD
    \longrightarrow
    \mathbb H_YD,
}
\]
so brave-new motivic exceptional integration is the primary
homotopy-invariant shadow of the differential integration;

\item the differential trace on the realized de Rham stage system is a
morphism of coherent cochain objects and therefore preserves the coherent
\(d_{\inf}\), the exact-cochain classifier, the specified nullhomotopies of
\(d_{\inf}^2\), and all higher cochain coherences; no preservation of the
complete totalization inverse limit is inferred merely from exactness;

\item the remaining BNV compatibility is controlled by the already
constructed obstruction
\[
\boxed{
    \ObsBNV(p)
    \simeq
    \Sigma\mathcal Z_Y\widetilde p_!\mathcal S_X.
}
\]
No general smooth vanishing theorem is asserted.  Whenever this obstruction
vanishes, the same differential exceptional trace preserves the universal BNV
interval primitive;

\item for every represented spectral \(\acute{e}\)tale morphism of finite
presentation, the smooth-site de Rham exchange is an equivalence
\[
\boxed{
    p^*_{\Nis}\mathscr F_Y
    \xrightarrow{\simeq}
    \mathscr F_X
}
\]
for the total de Rham coefficient, every totalization stage, normalized
cochains, and coherently closed cochains; after realization
\[
    p^{*,0}\widehat{\mathscr F}_Y
    \xrightarrow{\simeq}
    \widehat{\mathscr F}_X;
\]

\item for every separated spectral \(\acute{e}\)tale morphism of finite
presentation, \(T_p=0\), the intrinsic source coefficient is therefore the
differential exceptional pullback of the target, and the counit defines an
intrinsic differential de Rham pushforward
\[
\boxed{
    p_{\dR,!}^{\partial}:
    \widetilde p_!\widehat{\mathscr F}_X
    \longrightarrow
    \widehat{\mathscr F}_Y.
}
\]
Its motivic shadow is the corresponding motivic exceptional trace;

\item for the verified Nisnevich-locally split finite
\(\acute{e}\)tale class, the BNV obstruction vanishes, so the same intrinsic
differential de Rham pushforward also preserves universal BNV interval
integration;

\item for the split fold
\[
    p:\coprod_{i=1}^dY\to Y,
\]
the intrinsic differential de Rham trace is explicitly the codiagonal
\[
\boxed{
    D^{\oplus d}\longrightarrow D,
}
\]
with componentwise BNV integration and motivic finite
\(\acute{e}\)tale shadow;

\item a smooth correspondence
\[
    X\longleftarrow Z\xrightarrow{\pi}B
\]
admits differential transgression once a concrete coefficient trace
\[
    e^{*,0}D_X
    \longrightarrow
    \widetilde\pi^{\,!}D_B
\]
has been constructed.  No blanket orientation, regulator package, analytic
realization, or differential six-functor axiom is inserted to manufacture
such a map.
\end{enumerate}

The chapter introduces no second de Rham localization, no second brave-new affine line, no second
stable motivic category, and no independent differential-cohomological formalism.  The effective
infinitesimal cochains and their coherent differential are extracted from the inherited infinitesimal
geometry by effective \v{C}ech descent, additive stabilization, totalization, stable Dold--Kan, and
the complete-filtration/coherent-cochain correspondence.  Their first-order linearization is the
cotangent differential, but the normalized cochains themselves retain higher infinitesimal orders.
Accordingly, the \v{C}ech-totalization filtration is not identified with the Hodge filtration, and no
exterior-cotangent presentation or spectral HKR theorem is assumed.  The ordered Alexander--Whitney
maps are used with their associative, unital, and Leibniz identities in the homotopy category; no
separate homotopy-coherent multiplicative totalization refinement is asserted.  No principal-parts
algebra, crystalline-jet construction, represented formal-completion model, regulator package, or
analytic realization is used in the definition; the only completion object introduced is the
intrinsic pullback
\[
    Y\times_{Y_{\dR/S}}Z_{\dR/S}
\]
formed from the already established relative de Rham reflector.

The integration statement is likewise separated by logical strength.  For every smooth separated
finite-presentation morphism, differential exceptional integration exists on the differential
dualizing coefficient and has motivic \(p_!\) as its shadow.  BNV naturality is asserted when the
constructed BNV obstruction vanishes.  For every separated spectral \(\acute{e}\)tale morphism of
finite presentation, the intrinsic source de Rham coefficient itself is identified with the
exceptional pullback, so the intrinsic differential de Rham pushforward and its motivic shadow are
unconditional.  For the verified Nisnevich-locally split finite \(\acute{e}\)tale class, the BNV
obstruction also vanishes, and the same intrinsic differential de Rham pushforward preserves the BNV
interval integral unconditionally.
\end{theorem}

\begin{proof}
Combine the constructions and theorems of the preceding sections.
\end{proof}
